\documentclass[11pt,a4paper]{article}

\usepackage[T1]{fontenc}
\usepackage[utf8]{inputenc}
\usepackage[british]{babel}
\usepackage{lmodern}
\usepackage{microtype}
\usepackage[a4paper,margin=1in]{geometry}
\usepackage{amsmath,amssymb,amsthm,mathtools}
\usepackage{enumitem}
\usepackage{booktabs,tabularx,array}
\usepackage{xcolor}
\usepackage{hyperref}
\hypersetup{
  colorlinks=true,
  linkcolor=blue!60!black,
  citecolor=blue!60!black,
  urlcolor=blue!60!black,
  hypertexnames=false,
  pdftitle={Trace-Tempered Choice: A Self-Contained Axiomatisation with Explicit Payoffs and Records},
  pdfauthor={Daniele Condorelli},
  pdfsubject={A behavioural characterization of expected utility plus mutual information}
}
\setlength{\emergencystretch}{2em}
\interfootnotelinepenalty=10000

\newtheorem{theorem}{Theorem}
\newtheorem{proposition}[theorem]{Proposition}
\newtheorem{lemma}[theorem]{Lemma}
\newtheorem{corollary}[theorem]{Corollary}
\theoremstyle{definition}
\newtheorem{definition}[theorem]{Definition}
\newtheorem{remark}[theorem]{Remark}

\newcommand{\D}{\Delta}
\newcommand{\E}{\mathbb E}
\newcommand{\I}{I}
\newcommand{\Sh}{H}
\newcommand{\supp}{\operatorname{supp}}
\newcommand{\Id}{\operatorname{Id}}
\newcommand{\Rnum}{\mathbb R}
\newcommand{\R}{\mathbb R}
\newcommand{\KL}{D_{\mathrm{KL}}}
\newcommand{\1}{\mathbf 1}

\title{Trace-Tempered Choice\\[0.35em]
\large A Self-Contained Axiomatisation with Explicit Payoffs and Records}
\author{Daniele Condorelli\\University of Warwick}
\date{Draft of \today}

\begin{document}
\maketitle

\begin{abstract}
\noindent
An action can matter twice: through the payoff it produces and through the visible trace it leaves.  We separate these
roles in the primitive domain.  A finite channel maps actions into pairs consisting of a payoff-bearing outcome and an
explicit record.  Continuous weak order, coherent comparison environments, data processing on both sides of the
channel, and a branchwise sure-thing principle characterize
\[
  \E[u(O)]+\lambda \I(A;O,R),\qquad \lambda>0.
\]
No primitive payoff--trace separability axiom is required.  It is derived from the temporal asymmetry of the domain:
records may arrive first, while payoff is terminal.  The proof isolates that derivation in a calibrated branch-recursion
lemma.  A complete proof of the mutual-information characterization for the pure-trace restriction is included as an
appendix, so the paper is logically independent of earlier versions of the traceability paper.  Every primitive axiom is
shown to be independent of the others.
\end{abstract}

% ===============================================================
\section{Introduction}
% ===============================================================

\begin{quote}
\emph{Shower upon him every earthly blessing \ldots\ He would even risk his cakes and would deliberately desire the
most fatal rubbish, the most uneconomical absurdity \ldots\ simply in order to prove to himself---as though that were
so necessary---that men still are men and not the keys of a piano.\ \ldots\ If you say that all this, too, can be
calculated and tabulated \ldots\ then man would purposely go mad in order to be rid of reason and gain his point!}

\hfill --- Fyodor Dostoevsky, \emph{Notes from the Underground}, Part I, Chapter VIII \cite{Dostoevsky1864}
\end{quote}

Dostoevsky's provocation has the structure of a choice problem.  The cakes represent material payoff, and the
Underground Man is willing to risk them so that his conduct can attest to something: that it is his own rather than
the mechanical implication of a table.  Yet the proof he seeks is self-defeating.  An act fixed in advance cannot
provide it, and deliberate perversity offers no escape: rebellion, too, can be “calculated and tabulated.”  His
final resort is therefore paradoxical: he would “purposely go mad in order to be rid of reason”---not because
madness is what he values, but because it is his last imagined way of preventing the table from closing before he
acts.  This paper takes that evidentiary idea literally, in a deliberately narrow statistical sense.  In Dostoevsky's
example, conduct is observed directly: the act supplies its own evidence.  We consider the more general case in which
an act is observed only through a possibly noisy payoff-bearing outcome and an explicit visible record.

For visible consequences to carry information about an act, two conditions are needed.  There must be uncertainty
beforehand about which act will occur, and the eventual payoff--record pair must discriminate statistically among the
possible acts.  Either condition can fail independently.  A lottery over acts leaves no trace when every act generates
the same distribution of visible consequences; conversely, even a consequence that perfectly reveals the act conveys
no new information when the act was already known.  We call the non-instrumental preference for this conjunction a
preference for \emph{traceable agency}.  The paper axiomatizes it jointly with ordinary preferences over payoff risk.

The idea has a simple formal expression.  Let $A$ be a finite set of actions, $O$ a finite set of payoff-bearing
outcomes, $R$ a finite set of explicit visible records, and $K:A\to\Delta(O\times R)$ the channel relating actions to
visible consequences.  An action lottery $q\in\Delta(A)$ gives the distribution of the realised action; conditional
on that action, $K$ gives the joint distribution of payoff and record.  For each fixed channel $K$, the primitive is a
weak order $\succeq_K$ over action lotteries.  Thus $q\succeq_Kp$ means that $q$ is evaluated at least as highly as
$p$ in environment $K$.  Record processing may rewrite what is visible while leaving payoff fixed, and action
processing may coarsen the description of the realised action; neither operation can increase traceability.  The
remaining axioms impose order, continuity, and comparison coherence, together with a dominance condition for records
that arrive in stages.

The main result is an if-and-only-if characterization.  A family $\{\succeq_K\}_K$ satisfies the eight axioms if and
only if there are a nonconstant payoff utility $u:O\to\mathbb R$ and $\lambda>0$ such that, for every finite joint
channel $K$ and every $q,p\in\Delta(A)$,
\[
 q\succeq_Kp
 \quad\Longleftrightarrow\quad
 \E_{q,K}[u(O)]+\lambda I_{q,K}(A;O,R)
 \ge
 \E_{p,K}[u(O)]+\lambda I_{p,K}(A;O,R).
\]
Reversing the three trace-facing clauses yields the trace-averse dual with $\lambda<0$; replacing them by
indifference yields expected utility, or $\lambda=0$.

The information term has a simple interpretation.  Observing $(O,R)$ induces a posterior belief about the realised
action.  Its traceability is the expected reduction in uncertainty about $A$:
\[
 I_{q,K}(A;O,R)
 =
 \Sh(q)-\E_{q,K}\!\left[\Sh\bigl(q(\cdot\mid O,R)\bigr)\right].
\]
If visible consequences are independent of actions, the value is zero.  If the action is known in advance because
$q$ is degenerate, the value is also zero.  If each action in the support of $q$ produces a distinct visible
consequence, the value is $\Sh(q)$---Dostoevsky's case, in which conduct is observed directly.  Thus the criterion is
not a taste for randomisation as such.  Randomisation matters only because it creates uncertainty about agency for the
visible consequence to resolve.

A more mundane example fixes ideas.  A donor can route a gift to a named project or to a pooled fund; the material
effect is payoff-bearing, while the public attribution is an explicit record.  Holding the distribution of material
effects fixed, if the two routes generate the same distribution over records, the record is independent of the act
and its additional traceability is zero, however uncertain the giving.  If the routes report differently,
traceability is positive exactly when the realised route is genuinely uncertain beforehand---the donor wavers, or
delegates the decision---and it reaches $\Sh(q)$ when the visible consequence identifies the route.  A donor certain
to give by name leaves no trace in this sense: a perfectly attributed gift that was never in doubt yields a record
that only confirms what was already known.  If material consequences themselves differ across routes, they too form
part of the trace.  The criterion prices the news a visible consequence carries about conduct, not knowledge of it.

The example ranks lotteries within a single fixed environment; the representation also speaks across environments.
The environment itself is never an object of choice: the theory begins neither with a preference over pairs $(q,K)$
nor with a numerical index on channels.  When lotteries carried by different channels must be compared, the channels
are embedded as labelled blocks of a single larger environment, and the comparison is an ordinary one between two
block-supported lotteries in a fixed problem.  The theorem represents all such comparisons by the same $u$ and
$\lambda$ on one scale: it does not merely say that each environment carries a ranking ordinally equivalent to the
displayed criterion; it identifies common payoff and trace units across environments.

Traceable agency cannot be reduced to any preference over consequences alone, because here the action is part of
the consequence: what is valued is the joint distribution of the action with the visible pair, not the distribution
of the visible pair by itself.  Consider a channel with three actions, a constant payoff, and two records.  The
first two actions produce different records with certainty, while the third produces either record with equal
probability.  A fair lottery over the first two actions and a certain choice of the third induce the same joint
distribution of payoff and record.  Yet in the first case the record identifies the realised action, whereas in the
second it conveys no information about it.  Every index of the induced consequence distribution---expected utility
over payoffs, or over payoff--record pairs---is indifferent, while the traceability component assigns one bit to
the first lottery and zero to the second.

This difference does not make traceable agency a substitute for expected utility.  The point of the present theorem
is to axiomatize their joint decomposition rather than assume it:
\[
 V(q,K)=\E_{q,K}[u(O)]+\lambda I_{q,K}(A;O,R).
\]
The payoff--record distinction is essential to that identification.  It permits records to be processed without
permitting payoff to be garbled, and it gives branchwise continuation a temporal asymmetry: records may resolve first,
whereas payoff is terminal.  Those restrictions derive common payoff units, the additive decomposition, and matching
neutrality.  No primitive payoff--trace separability axiom is imposed.

A willingness to pay for decision rights and payoff autonomy is well documented.  Bartling, Fehr, and Herz infer a
positive intrinsic value of decision rights for 83 percent of participants assigned to the principal role; the average
control premium is 16.7 percent and significantly positive in all ten parameterisations, and an independent
replication closely reproduces their individual-level results
\cite{BartlingFehrHerz2014,BuffatPraxmarerSutter2023}.  Owens, Grossman, and Fackler elicit beliefs separately:
expected-money maximisation predicts betting on oneself in 56.4 percent of decisions, yet subjects do so in
64.9 percent, sacrificing 8--15 percent of expected earnings \cite{OwensGrossmanFackler2014}.  Neither design varies
what consequences reveal about the act, so neither identifies the traceability channel.

Evidence closer to the action--consequence channel comes from two studies, the second incentivised, in which subjects
prefer action menus with greater Jensen--Shannon divergence among consequence distributions even when current expected
payoff and outcome entropy are fixed \cite{MistryLiljeholm2016,NortonLiljeholm2020}.  Under their uniform prior over
two actions, this index is exactly mutual information.  Both designs, however, let token values change after the menu
is chosen, so divergence also provides instrumental flexibility; a later model comparison finds that dynamic expected
utility best explains the choices \cite{Liljeholm2022}.  These studies therefore do not isolate a non-instrumental
value of mutual information and, to our knowledge, no existing experiment does.

The model suggests the following direct test.  Fix a channel with four actions, $a_1,\dots,a_4$.  Every action yields
the same monetary payoff together with a binary non-payoff record.  Actions $a_1$ and $a_2$ each produce a different
symbol; actions $a_3$ and $a_4$ each produce either symbol with equal probability.  Participants choose between a fair
coin flip over $\{a_1,a_2\}$ and one over $\{a_3,a_4\}$.  The two flips induce the same joint distribution of
payoff and record and involve the same degree of deliberate randomisation; only the first makes the record
informative about the realised action.  Traceable agency therefore predicts a strict preference for it.  Because the
record affects neither material payoffs nor any later decision, this comparison already excludes instrumental
flexibility.\footnote{A garbling control completes the design.  Repeat the choice in a second treatment with a channel
that writes a fair record independently of the realised action: both flips then have zero traceability, and the
preference should disappear.  A preference surviving the garbling would indicate salience or framing effects rather
than a value of trace.}

Let us now discuss the eight axioms.  The structural axioms \textup{(A1)--(A3)} provide the comparison framework.  Each
channel carries a continuous weak order over action lotteries.  When lotteries associated with different channels
must be compared, the channels are embedded as disjoint components of one larger environment; duplicating a component
or adding components not involved in the comparison must not affect the ranking.

Axioms~\textup{(A4)} and \textup{(A5)} state the two sides of data processing.  The agent never gains when her
record is rewritten after the fact, holding payoffs fixed, nor when her action is reported noisily.  Because action processing is stated on the joint law, reported actions reached with zero probability impose
no constraint: their irrelevance is a consequence of the axiom rather than a convention appended to it.

Axiom~\textup{(A6)} is a sure-thing principle for records that arrive in stages.  If one continuation is weakly
preferred after every first-stage record that can occur, then the branch-by-branch replacement cannot make the
compound environment worse; a strict improvement on a positive-probability branch makes it strictly better.  This
condition supplies the system's cardinal discipline.  It implies invariance to common public randomisation and
calibrates continuation values across posterior action lotteries.

The two orientation axioms separate payoff from trace.  Axiom~\textup{(A7)} requires some pair of payoff outcomes to
be strictly ranked.  Axiom~\textup{(A8)} requires full revelation of a genuinely uncertain action to be strictly
preferred to an uninformative record when payoff is fixed.  The first makes $u$ nonconstant; the second selects
$\lambda>0$.  Neither imposes a numerical scale.

A sketch of the proof follows the theorem in Section~\ref{sec:theorem}; complete arguments, including a self-contained proof of the
pure-record characterisation on which the trace component rests, are in the Appendix.

% ===============================================================
\section{Related Literature}
% ===============================================================

The paper belongs to the decision-theoretic tradition that enriches the object of preference to capture motives that
expected utility over final payoffs alone leaves out.  Kreps \cite{Kreps1979} and Dekel, Lipman, and Rustichini
\cite{DekelLipmanRustichini2001} represent flexibility and subjective uncertainty about future tastes through
preferences over menus, and the freedom-of-choice literature initiated by Pattanaik and Xu \cite{PattanaikXu1990}
ranks opportunity sets by the freedom they embody; here decision rights and feasible actions are held fixed, and what
varies is ex-post authorship rather than ex-ante opportunity: whether the realised visible consequence statistically
records the action.  Gul and Pesendorfer \cite{GulPesendorfer2001} represent temptation and commitment; Kreps and
Porteus \cite{KrepsPorteus1978} enrich the domain to temporal lotteries so that the timing of the resolution of
uncertainty can itself carry utility; and Grant, Kajii, and Polak \cite{GrantKajiiPolak1998} characterise within that
domain an intrinsic preference for information, valued independently of any planning use.  The information valued
here flows in the opposite direction---not what signals teach the agent about the world, but what visible consequences
teach the world about the agent.  Deliberate randomisation has also been axiomatised through a hedging motive
\cite{CerreiaVioglioDillenbergerOrtolevaRiella2019}; unlike that motive, traceability depends on the observability of
the action--consequence record.  The nearest neighbours lie in the signaling tradition descending from Spence
\cite{Spence1973}, in which actions are informative about hidden characteristics: social-image models explain
attribution-seeking behaviour, such as named rather than anonymous giving, through observers' beliefs entering utility
\cite{GlazerKonrad1996,BenabouTirole2006,AndreoniBernheim2009}; self-image models direct the same inference inward,
towards beliefs about one's own latent traits \cite{BenabouTirole2002}; and in Bernheim's conformity model
\cite{Bernheim1994} agents \emph{suppress} the legibility of their predispositions, an instrumental counterpart of
the deniability orientation developed below.  In all of these, what is inferred matters through its
consequences---esteem, sanction, or self-regard; here no observer's response enters utility, beliefs figure only as
the statistical yardstick against which the trace is read, and the preference survives anonymity.  Attribution,
expressive value, and reputation are possible microfoundations for the behaviour the axioms describe, but they are
not primitives of the theorem; what none of these literatures has made the object of preference is the trace itself.

The relation to axiomatic information theory is different in kind.  Shannon \cite{Shannon1948} derived entropy from
continuity, monotonicity on uniform distributions, and a grouping recursion; Khinchin \cite{Khinchin1957} and Faddeev
\cite{Faddeev1956} progressively weakened these postulates, Faddeev reducing them to symmetry, continuity of the
binary entropy, and the recursion that appears here as a derived object.  R\'enyi \cite{Renyi1961} showed that
relaxing the grouping postulate to general quasi-arithmetic means, while retaining additivity over independent
experiments, yields the one-parameter family bearing his name, in which Shannon entropy is the unique member still
satisfying the recursion; Acz\'el and Dar\'oczy \cite{AczelDaroczy1975} give a systematic treatment by functional
equations, and Csisz\'ar \cite{Csiszar2008} surveys the field.  Mutual information itself is typically characterised
indirectly, as entropy reduction or as the relative entropy between the joint law and the product of its marginals:
Hobson \cite{Hobson1969} characterises relative entropy by grouping and additivity conditions, and Baez and Fritz
\cite{BaezFritz2014}, extending the treatment of entropy as a measure of information loss in Baez, Fritz, and Leinster
\cite{BaezFritzLeinster2011}, derive it as the essentially unique functorial, convex-linear, lower semicontinuous
invariant of Bayesian inference.

All of these characterisations take a numerical functional as primitive and impose structural equations on
distributions, random variables, or information-loss maps.  The present theorem starts instead from ordinal preference
relations over lotteries in fixed joint environments and derives the functional: common cardinal structure, the
expected-utility component, the information chain rule, and finally Faddeev's recursion are consequences of
behavioural axioms rather than postulates on a given function.  The Shannon part of the exercise is thus the
information-theoretic analogue of expected-utility theory, in which the numerical form of the criterion is the output
of a representation theorem rather than its input; what is characterised is not an abstract measure of information
but the preference of an agent for whom that measure is the value of traceable agency.

A smaller literature takes information structures themselves as the ranked domain.  Gilboa and Lehrer
\cite{GilboaLehrer1991} characterise which functions on partitions are values of information for some Bayesian
decision problem, and Azrieli and Lehrer \cite{AzrieliLehrer2008} extend the characterisation to stochastic
structures; in both, value is instrumental by construction.  Cabrales, Gossner, and Serrano
\cite{CabralesGossnerSerrano2013} come closest to the present representation from the instrumental side: for
ruin-averse investors, the ranking of information structures is represented by entropy reduction, so that mutual
information emerges from a specific class of decision problems rather than from ordinal axioms.  On the
non-instrumental side, Caplin and Leahy \cite{CaplinLeahy2001} put beliefs into the utility function, and Ely, Frankel,
and Kamenica \cite{ElyFrankelKamenica2015} derive the entertainment value of suspense and surprise from preferences
over the path of beliefs.  The present paper combines the two strands: its domain consists of lotteries in fixed
payoff--record channels, its trace valuation is non-instrumental, and, unlike either strand, the functional is not
posited or induced by a decision problem but derived ordinally and pinned down jointly with material utility.

The closest formal literature axiomatises the cost of information.  Pomatto, Strack, and Tamuz
\cite{PomattoStrackTamuz2023} take as primitive a cardinal cost functional on Blackwell experiments
\cite{Blackwell1953}, defined without reference to a prior, and show that Blackwell monotonicity, continuity,
additivity over independent experiments, and constant marginal cost characterise the log-likelihood ratio costs:
weighted sums of Kullback--Leibler divergences between the state-conditional signal distributions.  Caplin, Dean, and
Leahy \cite{CaplinDeanLeahy2022} and Denti \cite{Denti2022} work from state-dependent stochastic choice data generated
by rationally inattentive agents and characterise posterior-separable cost functions, with the Shannon cost singled
out by invariance conditions; Frankel and Kamenica \cite{FrankelKamenica2019} characterise valid measures of
information as expected reductions of measures of uncertainty, with entropy reduction selected by additivity
requirements.  Throughout this literature, information is instrumental: it concerns an exogenous state and is costly,
or valuable, as an input into a downstream decision.  This has two structural consequences.  Blackwell monotonicity is
a theorem or an unquestioned desideratum, since a decision-maker who can condition on a signal is weakly better off
with a finer one; and the primitive object is cardinal, or is disciplined by an optimisation model that supplies
cardinal structure.

The present paper shares the destination of that literature---coherence over mixing and sequential resolution forces
the Shannon scale---but differs in primitive, in direction, and in what is derived rather than assumed.  The primitive
here is a family of ordinal weak orders over lotteries on the agent's own actions; there is no downstream decision,
and prior dependence is essential rather than excluded: the trace component is the whole map
$(q,K)\mapsto I_{q,K}(A;O,R)$.  Indeed, the two families are disjoint rather than nested: mutual information is
prior-dependent and subadditive over independent experiments, so it violates the additivity that defines the
log-likelihood ratio class, while fully revealing experiments, which have infinite log-likelihood ratio cost, receive
the finite value $\Sh(q)$ here.  Because the primitive is ordinal and fibred by the prior, the cardinal structure that
the cost literature postulates must here be constructed; branchwise continuation consistency performs that
calibration jointly with the payoff scale.  Finally, record and action data processing,
Axioms~\textup{(A4)} and \textup{(A5)}, are primitive postulates of taste rather than theorems of instrumental
rationality; as the sign trichotomy (Corollary~\ref{cor:signduality}) makes precise, their joint reversal is equally
coherent and characterises a preference for deniability.  Where the two traditions do meet---at the Shannon
functional, reached from opposite ends of the channel---they axiomatise different objects: there, the price of learning
about the world; here, a taste for being learned about, tempered by ordinary payoff utility.

% ===============================================================
\section{The Model}
% ===============================================================

Fix a finite set $O$ of \emph{payoff-bearing outcomes}, with $|O|\ge2$.  Let $A$ be a nonempty finite set of
\emph{actions} and $R$ a nonempty finite set of \emph{explicit visible records}.  A joint channel is a stochastic
kernel
\[
  K:A\longrightarrow\D(O\times R).
\]
For each finite joint channel $K$, the primitive behavioural datum is a binary relation
\[
 \succeq_K\quad\text{on}\quad\D(A).
\]
The interpretation is that $q$ is a lottery over actions used by the agent and $q\succeq_Kp$ means that, in the fixed
environment $K$, $q$ is evaluated at least as highly as $p$.  The lottery is the procedure being evaluated; the
channel mapping actions into payoff--record pairs is not an object of choice.

\bigskip
Axioms are imposed on the family $\{\succeq_K\}_K$.

\begin{description}[leftmargin=*,style=sameline]

\item[(A1) Weak order.]
For every finite joint channel $K$, $\succeq_K$ is complete and transitive on $\D(A)$.

\item[(A2) Continuity.]
For each fixed triple of alphabets $(A,O,R)$, the set
\[
 \{(K,q,p):q\succeq_Kp\}
 \subseteq \D(O\times R)^A\times\D(A)^2
\]
is closed.

\end{description}

\paragraph{Comparison environments.}
The remaining axioms compare a channel with a garbling of it and a compound experiment with its branches.  Each
relation $\succeq_K$ ranks only lotteries, so these comparisons need a domain in which two channels can appear in
one preference statement.  Labelled blocks provide it, much as objective lotteries enrich Savage's domain in
Anscombe and Aumann \cite{AnscombeAumann1963}, and without making the channel an object of choice.

For
$K:A\to\D(O\times R)$ and $L:B\to\D(O\times S)$, let $K\sqcup L$ be the block-diagonal joint channel with action set
$(A\times\{0\})\cup(B\times\{1\})$, payoff set $O$, and record set
$(R\times\{0\})\cup(S\times\{1\})$, defined by
\[
 (K\sqcup L)(o,(r,0)\mid(a,0))=K(o,r\mid a),
 \qquad
 (K\sqcup L)(o,(s,1)\mid(b,1))=L(o,s\mid b),
\]
with all cross-block probabilities zero.  For $q\in\D(A)$ and $p\in\D(B)$, let $q^0$ and $p^1$ denote their
block-supported copies, defined on the tagged action set by
\[
\begin{aligned}
 q^0(a,0)&:=q(a), & q^0(b,1)&:=0,\\
 p^1(a,0)&:=0,    & p^1(b,1)&:=p(b),
\end{aligned}
\qquad(a\in A,\ b\in B).
\]
When the blocks have the same action alphabet, $q^0$ and $q^1$ are
the two tagged copies of the same lottery.  For a finite family
$\{K_i:A_i\to\D(O\times R_i)\}_{i\in I}$, write $\bigsqcup_{i\in I}K_i$ for the analogous block environment and
$q_i^i$ for the lottery that assigns $q_i(a)$ to $(a,i)$ and zero to every other block.  Write
\[
 (q,K)\succeq(p,L)
 \quad\Longleftrightarrow\quad
 q^0\succeq_{K\sqcup L}p^1.
\]
Each instance of this shorthand is a single comparison inside the two-block environment $K\sqcup L$, which is
itself built from the two pairs; $\succ$ and $\sim$ denote its asymmetric and symmetric parts.  The pair $(q,K)$
has no meaning in isolation.\footnote{A complete numerical example.  Let $A=\{a_1,a_2\}$, $B=\{b_1,b_2\}$,
$O=\{o,o'\}$, and let the record alphabets be $R=\{r,r'\}$ and $S=\{s\}$.  With rows indexed by actions and
columns by payoff--record pairs,
\[
 K=\begin{array}{c|cccc}
    & (o,r) & (o',r) & (o,r') & (o',r')\\\hline
 a_1& 0.5 & 0 & 0 & 0.5\\
 a_2& 0 & 0.7 & 0.3 & 0
 \end{array}
 \qquad\quad
 L=\begin{array}{c|cc}
    & (o,s) & (o',s)\\\hline
 b_1& 0.3 & 0.7\\
 b_2& 0.6 & 0.4
 \end{array}
\]
The block environment $K\sqcup L$ has tagged action set $\{(a_1,0),(a_2,0),(b_1,1),(b_2,1)\}$, payoff set $O$,
and tagged record set $\{(r,0),(r',0),(s,1)\}$; its kernel is block-diagonal, $K$ on the first two rows, $L$ on
the last two, and every cross-block entry zero.  For $q=(0.3,0.7)$ and $p=(0.8,0.2)$, the block-supported copies
are $q^0=(0.3,\,0.7\,;\,0,\,0)$ and $p^1=(0,\,0\,;\,0.8,\,0.2)$, the semicolon separating the blocks.  Then
$(q,K)\succeq(p,L)$ abbreviates $q^0\succeq_{K\sqcup L}p^1$: one comparison of two ordinary lotteries in this one
environment.  Both block-supported lotteries deliver $o$ with probability $0.36$, so under the representation of
Theorem~\ref{thm:main} the comparison turns on the trace terms alone---approximately $0.88$ against $0.04$
bits---and is strict.}

\begin{description}[leftmargin=*,style=sameline]

\item[(A3) Block-comparison coherence.]
\emph{Duplication:} for every $K:A\to\D(O\times R)$ and all $q,p\in\D(A)$,
\[
 q\succeq_Kp
 \quad\Longleftrightarrow\quad
 (q,K)\succeq(p,K).
\]
\emph{Irrelevant blocks:} for every finite block environment $\bigsqcup_{k\in I}K_k$, every ordered pair of
distinct blocks $i,j\in I$, and all $q_i\in\D(A_i)$ and $q_j\in\D(A_j)$,
\[
 q_i^i\succeq_{\bigsqcup_{k\in I}K_k}q_j^j
 \quad\Longleftrightarrow\quad
 (q_i,K_i)\succeq(q_j,K_j).
\]

Block labels should not alter within-environment rankings, nor should unused blocks affect the comparison between
active blocks.

\end{description}

\paragraph{Processing and compounding.}
The next three axioms use two ways of transforming an environment and one way of composing environments.  A
\emph{record processor} is a stochastic kernel $T(r'\mid o,r)$ that may depend on the payoff but leaves it
unchanged; it transforms $K:A\to\D(O\times R)$ into
\begin{equation}
 (KT)(o,r'\mid a):=\sum_r K(o,r\mid a)T(r'\mid o,r).
 \label{eq:record-processing}
\end{equation}
An \emph{action report} is a stochastic kernel $S:A\to\D(B)$; it transforms a lottery $q\in\D(A)$ into
$(qS)(b):=\sum_aq(a)S(b\mid a)$ and $K$ into any channel $\widehat K:B\to\D(O\times R)$ satisfying
\begin{equation}
 (qS)(b)\widehat K(o,r\mid b)
 =\sum_a q(a)S(b\mid a)K(o,r\mid a)
 \quad\text{for every }(b,o,r).
 \label{eq:action-processing}
\end{equation}
If $(qS)(b)=0$, equation~\eqref{eq:action-processing} leaves the corresponding row of $\widehat K$ unrestricted,
and every statement below applies to every consistent completion.  Both transformations preserve the $O$-marginal.
Finally, a first-stage record channel $P:A\to\D(Y)$ and a profile of continuation channels
$K^y:A\to\D(O\times R_y)$ compose into the \emph{compound environment}
\begin{equation}
 (P\triangleright\{K^y\})(o,(y,r)\mid a)
 :=P(y\mid a)K^y(o,r\mid a).
 \label{eq:compound}
\end{equation}
Given $q\in\D(A)$, a branch $y$ is \emph{reached} if $m(y):=\sum_aq(a)P(y\mid a)>0$, and its branch posterior is
$q_y(a):=q(a)P(y\mid a)/m(y)$.

\begin{description}[leftmargin=*,style=sameline]

\item[(A4) Record data processing.]
For every $K:A\to\D(O\times R)$, every record processor $T$, and every $q\in\D(A)$,
\[
 (q,K)\succeq(q,KT).
\]
The agent never gains when her record is garbled.

\item[(A5) Action data processing.]
For every $K:A\to\D(O\times R)$, every action report $S:A\to\D(B)$, every $q\in\D(A)$, and every channel
$\widehat K$ satisfying \eqref{eq:action-processing},
\[
 (q,K)\succeq(qS,\widehat K).
\]
Nor does she gain when her action is reported noisily, holding the payoff lottery fixed.

\item[(A6) Branchwise continuation consistency.]
For every first-stage record channel $P:A\to\D(Y)$, all continuation profiles $\{K^y\}$ and $\{L^y\}$, and every
$q\in\D(A)$: if
\[
 (q_y,K^y)\succeq(q_y,L^y)
 \quad\text{for every reached }y,
\]
then
\[
 \bigl(q,\,P\triangleright\{K^y\}\bigr)\succeq\bigl(q,\,P\triangleright\{L^y\}\bigr),
\]
strictly if one reached branch comparison is strict.

In a compound environment the branching variable is a record, and payoffs arrive only in the continuations.  The
axiom is a sure-thing principle across the branches of an interim record: a profile weakly better on every reached
branch is weakly better overall, with unreached branches unrestricted exactly as Savage's null events are.

\item[(A7) Material relevance.]
For $o\in O$, let $K_o:\{\ast\}\to\D(O\times\{\ast\})$ deliver $o$ with an uninformative record,
$K_o(o,\ast\mid\ast)=1$.  For some $o^+,o^-\in O$,
\[
 (\delta_\ast,K_{o^+})\succ(\delta_\ast,K_{o^-}).
\]
The agent cares about material outcomes, and not only about their trace.

\item[(A8) Positive trace orientation.]
For every finite $A$ with $|A|\ge2$, every full-support $q\in\D(A)$, and every $o\in O$,
\[
 (q,K^{\mathrm{id}}_o)\succ(q,K^{0}_o),
\]
where $K^{\mathrm{id}}_o(o,a\mid a)=1$, with record alphabet $A$, and $K^0_o(o,\ast\mid a)=1$.
At a constant payoff, a record that reveals a genuinely uncertain action is strictly better than one that reveals
nothing: trace is sought, not shunned.

\end{description}

% ===============================================================
\section{The Representation Theorem}\label{sec:theorem}
% ===============================================================

For $q\in\D(A)$ and $K:A\to\D(O\times R)$, let
\[
 p_{qK}(o,r):=\sum_a q(a)K(o,r\mid a)
\]
be the induced distribution of visible consequences, and define
\begin{equation}
 \I_{qK}(A;O,R)
 :=\sum_{a,o,r}q(a)K(o,r\mid a)
 \log\frac{K(o,r\mid a)}{p_{qK}(o,r)}.
 \label{eq:mi}
\end{equation}
Summands with $q(a)K(o,r\mid a)=0$ are set to zero, the logarithm has any fixed base, and expectations below are
taken under the joint distribution $q(a)K(o,r\mid a)$.

\begin{theorem}[Trace-tempered choice]\label{thm:main}
For a family $\{\succeq_K\}_K$ as described above, the following are equivalent.
\begin{enumerate}[label=\textup{(\roman*)},leftmargin=*]
\item The family satisfies Axioms \textup{(A1)--(A8)}.
\item There are a nonconstant function $u:O\to\mathbb R$ and a number $\lambda>0$ such that, for every nonempty
finite $A,R$, every $K:A\to\D(O\times R)$, and every $q,p\in\D(A)$,
\begin{equation}
 q\succeq_Kp
 \quad\Longleftrightarrow\quad
 V(q,K)\ge V(p,K),
 \qquad
 V(q,K):=\E_{q,K}[u(O)]+\lambda\I_{qK}(A;O,R).
 \label{eq:representation}
\end{equation}
\end{enumerate}
Moreover, under these equivalent conditions, block-supported cross-channel comparisons are represented on the same
scale: for every finite block environment $\bigsqcup_{k\in I}K_k$, every two distinct blocks $i,j\in I$, and every
$q_i\in\D(A_i)$ and $q_j\in\D(A_j)$,
\[
 q_i^i\succeq_{\bigsqcup_{k\in I}K_k}q_j^j
 \quad\Longleftrightarrow\quad
 V(q_i,K_i)\ge V(q_j,K_j).
\]
\end{theorem}

The proof of the equivalence of \textup{(i)} and \textup{(ii)} is in Appendix~\ref{app:main-proof}; the moreover
clause follows because a block-supported lottery has a constant block label.

\paragraph{Proof sketch.}
That the representation satisfies the axioms is a direct verification.  The converse has four steps.

\emph{Step 1: reduction to terminal laws.}  Fix a full-support $q$.  Two channels that induce the same joint law
of realised payoff and posterior about the action are record processings of one another, so Axioms~\textup{(A1)},
\textup{(A3)}, and \textup{(A4)} rank them together (Lemma~\ref{lem:terminal-suff}); Axiom~\textup{(A5)} deletes
null actions, and applied to bijective reports in both directions it makes relabellings neutral.  Each pair $(q,K)$ is
thereby identified with its \emph{terminal law}, the joint distribution of realised payoff and posterior.

\emph{Step 2: cardinalisation.}  A public coin is a first stage that reveals nothing, so every branch posterior
equals $q$.  Axiom~\textup{(A6)} applied to such stages yields mixture independence over terminal laws
(Lemma~\ref{lem:public-coin}), and with continuity, Axiom~\textup{(A2)}, the mixture-space theorem gives a
continuous affine representation on each fixed-prior fibre.  All cardinal content enters here: every other axiom
is an ordinal, environment-by-environment condition.

\emph{Step 3: two scales.}  On singleton action sets the affine representation is expected utility, nonconstant by
Axiom~\textup{(A7)}; action reports transport the payoff scale across priors and alphabets, and block coherence
splices the fibre representations into one common scale (Lemmas~\ref{lem:payoff-scale}
and~\ref{lem:affine-bridge}).  On the constant-payoff restriction the axioms collapse to the pure-record
conditions of Appendix~\ref{pt:sec:environment}, whose characterisation identifies the residual order with mutual
information; Axiom~\textup{(A8)} signs its weight (Lemma~\ref{lem:trace-engine}).

\emph{Step 4: assembly.}  Axiom~\textup{(A6)} enters once more: inserting continuations into a reached branch is
an order isomorphism, and calibrating branches with terminal payoffs forces values to aggregate across branches at
their reached probabilities (Lemma~\ref{lem:branch-recursion}).  Sequentialising an arbitrary channel---announce
$(o,r)$ as a first-stage record, deliver $o$ terminally---is neutral by Axiom~\textup{(A4)}, and the recursion
returns $\E_{q,K}[u(O)]+\lambda\I_{qK}(A;O,R)$.

\paragraph{Uniqueness} To state uniqueness, fix one representation $V$ supplied by Theorem~\ref{thm:main}.  A real-valued index $W$ on all
finite lottery--channel pairs is a \emph{common-scale representation} if, for every channel
$K:A\to\D(O\times R)$ and every $q,p\in\D(A)$,
\[
 q\succeq_Kp
 \quad\Longleftrightarrow\quad
 W(q,K)\ge W(p,K),
\]
and, for every finite block environment $\bigsqcup_{k\in I}K_k$, every two distinct blocks $i,j\in I$, and all
$q_i\in\D(A_i)$ and $q_j\in\D(A_j)$,
\[
 q_i^i\succeq_{\bigsqcup_{k\in I}K_k}q_j^j
 \quad\Longleftrightarrow\quad
 W(q_i,K_i)\ge W(q_j,K_j).
\]

Call a common-scale representation $W$ \emph{branch-additive} if, for every first-stage record channel $P$, every
$q\in\D(A)$, and every two continuation profiles $\{K^y\}$ and $\{L^y\}$ admitted in Axiom~\textup{(A6)},
\begin{equation}
 W\bigl(q,P\triangleright\{K^y\}\bigr)
 -W\bigl(q,P\triangleright\{L^y\}\bigr)
 =\sum_{y:m(y)>0}m(y)
 \bigl[W(q_y,K^y)-W(q_y,L^y)\bigr].
 \label{eq:branch-additive-index}
\end{equation}
Thus continuation-value differences aggregate by their reached-branch probabilities; no arbitrary zero for material
utility is required.

\begin{proposition}[Uniqueness]\label{prop:uniqueness}
Suppose the equivalent conditions of Theorem~\ref{thm:main} hold.  For any real-valued index $W$ on all finite
lottery--channel pairs:
\begin{enumerate}[label=\textup{(\roman*)},leftmargin=*]
\item $W$ is a common-scale representation if and only if $W=\varphi\circ V$ for a single strictly increasing
$\varphi:[\underline u,\infty)\to\mathbb R$, where $\underline u:=\min_{o\in O}u(o)$.
\item A common-scale representation $W$ is branch-additive if and only if
$W=\alpha V+\beta$ for some $\alpha>0$ and $\beta\in\mathbb R$.
\end{enumerate}
\end{proposition}

The primitive remains ordinal: part \textup{(i)} permits any common increasing transformation.  Branch additivity is
an accounting requirement on an index, not another behavioural axiom; part \textup{(ii)} selects the affine gauge.
Consequently, if $(u,\lambda)$ and $(\widetilde u,\widetilde\lambda)$ are two trace-tempered representations, then
\[
 \widetilde u=\alpha u+\beta,
 \qquad
 \widetilde\lambda=\alpha\lambda
\]
for some $\alpha>0$ and $\beta\in\mathbb R$.  Once two payoff outcomes are assigned numerical utilities, $\lambda$ is
therefore uniquely measured in those payoff units per unit of information.

\paragraph{Sign} Uniqueness concerns scale; a complementary observation concerns sign.  Axioms~\textup{(A4)}, \textup{(A5)}, and
\textup{(A8)} orient the trace component.  Write $(\mathrm{A}4^-)$, $(\mathrm{A}5^-)$, and $(\mathrm{A}8^-)$ for
the clauses obtained by reversing their displayed comparisons, and write $(\mathrm{A}4^0)$,
$(\mathrm{A}5^0)$, and $(\mathrm{A}8^0)$ for the corresponding indifference clauses.

\begin{corollary}[Sign trichotomy]\label{cor:signduality}
Suppose Axioms~\textup{(A1)--(A3)}, \textup{(A6)}, and \textup{(A7)} hold.  Then:
\begin{enumerate}[label=\textup{(\roman*)},leftmargin=*]
\item under Axioms~\textup{(A4)}, \textup{(A5)}, and \textup{(A8)}, the family has the representation in
Theorem~\ref{thm:main}, with $\lambda>0$;
\item under $(\mathrm{A}4^-)$, $(\mathrm{A}5^-)$, and $(\mathrm{A}8^-)$, it has the same representation with
$\lambda<0$;
\item under $(\mathrm{A}4^0)$, $(\mathrm{A}5^0)$, and $(\mathrm{A}8^0)$, it is represented by expected utility,
equivalently by the same formula with $\lambda=0$.
\end{enumerate}
Thus the trace component has exactly the positive, negative, and neutral orientations.
\end{corollary}

\begin{proof}
Part \textup{(i)} is Theorem~\ref{thm:main}.  For parts \textup{(ii)} and \textup{(iii)}, the material calibration and
branch-recursion arguments are unchanged.  On the constant-payoff restriction, reversing both processing directions
and full-revelation orientation gives $-I$, while replacing them by indifference gives zero, by
Corollary~\ref{pt:cor:signduality}.  The converse follows directly from data processing and the two chain rules.
\end{proof}

For $q\in\D(A)$ and $K:A\to\D(O\times R)$, write
\begin{align*}
 \ell_{qK}(o)&:=\sum_r p_{qK}(o,r),\\
 \pi^{qK}_{or}(a)&:=\frac{q(a)K(o,r\mid a)}{p_{qK}(o,r)}
 \quad\text{when }p_{qK}(o,r)>0,\\
 \mu_{qK}&:=\sum_{(o,r):p_{qK}(o,r)>0}
 p_{qK}(o,r)\,\delta_{\pi^{qK}_{or}}.
\end{align*}

\begin{corollary}[Derived matching neutrality]\label{cor:matching}
Suppose Axioms \textup{(A1)--(A8)} hold.  Fix $A$ and $q\in\D(A)$.  If channels $K$ and $L$ satisfy
\[
 \sum_o\ell_{qK}(o)u(o)=\sum_o\ell_{qL}(o)u(o),
 \qquad
 \mu_{qK}=\mu_{qL},
\]
then $(q,K)\sim(q,L)$.  Thus the matching of payoff labels to posterior atoms has no additional value.
\end{corollary}

Indeed, $\ell_{qK}$ is the payoff marginal, $\mu_{qK}$ is the law of reached posteriors, and
\[
 \I_{qK}(A;O,R)=\Sh(q)-\int\Sh(\pi)\,d\mu_{qK}(\pi).
\]
In words: the agent values what she obtains and what the visible consequence reveals about what she did, but not
the coincidence of the two.  Winning and being seen to act are priced separately; there is no premium for the good
outcome arriving through the revealing record rather than alongside it.  Any such premium would be a
payoff--trace complementarity, and the corollary leaves no room for one.  This is the behavioural face of the
additive representation---the empirical content of the separability that the theorem derives rather than assumes.
The separation is achieved in the domain and in Axiom~\textup{(A6)}: compound environments may branch on interim
records, never on payoffs.

\begin{proposition}[Independence of the axioms]\label{prop:independence}
For each $i\in\{1,\ldots,8\}$ there is a family satisfying the other seven axioms and violating
Axiom~\textup{(A$i$)}.
\end{proposition}

The eight families are constructed in Appendix~\ref{app:main-proof}.  Three mark the substantive frontier of the
system.  A conditional-entropy penalty satisfies every axiom except record data processing and values the removal of
extraneous record noise.  A Gini posterior-reduction criterion satisfies every axiom except action data processing.
The satiation criterion $\E[u(O)]+\min\{I(A;O,R),c\}$, for $c>0$, satisfies every axiom except branchwise
continuation consistency and ceases to value trace once $c$ units of information are secured.

Taken together, the results separate orientation, scale, and behavioural content.  The sign trichotomy determines
whether trace is sought, avoided, or ignored; branch additivity selects the affine gauge; and no displayed axiom is
redundant.  Unlike the pure-record model, the magnitude of $\lambda$ is not merely a choice of information unit once
payoff utility has been calibrated: it is the identified payoff--trace trade-off.

% ===============================================================
\section{Choice implications}\label{sec:implications}
% ===============================================================

This section works out what an agent with the representation of Theorem~\ref{thm:main} does: which finite designs
falsify the model directly (\S\ref{sec:tracetest}), when she deliberately randomises (\S\ref{sec:randomisation}),
what she will pay for changes to the record technology (\S\ref{sec:record-price}), and what her behaviour can
never disclose about the environment (\S\ref{sec:cannot-reveal}).
Throughout, $\lambda$ is measured in payoff units per unit of information in the fixed base.  Proofs for this
section, together with full statements of the supporting results that the text presents in summary form, are
collected in Appendix~\ref{app:choice-proofs}; only one-line verifications are kept here.

\subsection{The trace test}\label{sec:tracetest}

The functional $V$ reads the joint distribution of the action and its visible consequence.  Existing theories of
choice read less: the distribution of consequences, the lottery, or both.  The following experiment separates
them.

In the first treatment the agent chooses between two fair coins over actions.  One flips between two actions
that leave distinct marks; the other flips between two actions that produce noise.  Both coins yield the same
payoff for certain and the same fifty--fifty distribution over records; only the first coin's record identifies
the realised action.  In the second treatment the record technology is disabled: the same choice is offered in a
channel where every action produces noise.

\begin{proposition}[The trace test]\label{prop:signature}
Four actions pay the same outcome $o_0$.  Two of them leave distinct marks: $a_1$ always produces record $r_1$,
and $a_2$ always produces the distinct record $r_2$.  The other two are indistinguishable: $a_3$ and $a_4$ each
produce $r_1$ or $r_2$ with equal probability.  Call this channel $K$, and let $q'$ be the fair coin over
$\{a_1,a_2\}$ and $q''$ the fair coin over $\{a_3,a_4\}$.
\begin{enumerate}[label=\textup{(\roman*)},leftmargin=*]
\item The two coins induce the same joint distribution of visible consequences, $p_{q'K}=p_{q''K}$; yet
\[
 V(q',K)-V(q'',K)=\lambda\log2>0,
 \qquad\text{so}\qquad
 q'\succ_Kq''.
\]
\item Let $K^\circ$ be the channel in which every action pays $o_0$ and produces $r_1$ or $r_2$ with equal
probability.  Then $q'\sim_{K^\circ}q''$.
\end{enumerate}
\end{proposition}

The prediction is strict preference for the revealing coin in the first treatment and exact indifference in the
second.  Part \textup{(i)} alone rules out every consequentialist theory.  The two coins induce the same
distribution over payoff--record pairs, so any criterion that evaluates only that distribution is indifferent
between them: expected utility over payoffs, expected utility that values records directly as consequences
through some $v(o,r)$, and every non-linear functional of $p_{q'K}$.  What separates the coins is the joint
distribution of the action with its visible consequence, and that is precisely what consequentialist theories do
not read.

Part \textup{(i)} restricts only how a theory reads consequences.  A theory that also reads the lottery can
match part \textup{(i)}.  Existing theories of deliberate randomisation are of this kind: in the hedging model
of Cerreia-Vioglio, Dillenberger, Ortoleva, and Riella \cite{CerreiaVioglioDillenbergerOrtolevaRiella2019}, and
in the preference for randomisation documented by Agranov and Ortoleva \cite{AgranovOrtoleva2017}, the value of
mixing depends only on the induced distribution and the lottery.  Call an index $W$ \emph{channel-blind} if it
reads only these two objects: there are functions $F$ and $h:\D(A)\to\Rnum$, fixed across channels, such that
\[
 W(q,K)=F\bigl(p_{qK},h(q)\bigr).
\]

The second treatment rules these out.  Nothing a channel-blind index reads changes between the treatments: in
all four lottery--channel combinations the distribution of visible consequences is uniform on
$\{(o_0,r_1),(o_0,r_2)\}$, and the lotteries are the same two lotteries.  A channel-blind index therefore
separates the coins in both treatments or in neither, so none reproduces the pattern of
Proposition~\ref{prop:signature}: strict preference under $K$ and indifference under
$K^\circ$.\footnote{If $h$ does not depend on how the actions are labelled---an entropy bonus, say---the first
treatment is enough: the permutation that swaps $a_1$ with $a_3$ and $a_2$ with $a_4$ turns $q'$ into $q''$ and
leaves $p_{qK}$ unchanged, so $W(q',K)=W(q'',K)$.  The second treatment is needed for arbitrary $h$.}

The test also rejects.  The first treatment separates the three cases of Corollary~\ref{cor:signduality}: the
revealing coin if $\lambda>0$, the other coin if $\lambda<0$, indifference if $\lambda=0$.  In the second
treatment every representation assigns both coins the same value, whatever $u$ and $\lambda$, so a strict
preference there rejects the model.

\subsection{Deliberate randomisation}\label{sec:randomisation}

A degenerate lottery leaves no trace, so an agent with $\lambda>0$ may deliberately randomise: she gives up
expected payoff to make her act genuinely uncertain, because only then can its consequence reveal it.  In a
payoff-based theory a sufficiently large payoff advantage makes a single action optimal.  Here it need not: an
action that leaves a mark no other action can leave is never wholly abandoned, however poorly it pays, because
the first grain of mass on it is infinitely informative per unit, and no finite payoff difference outweighs
that.  The next proposition states this, and gives the exact terms on which a pure action survives: its payoff
advantage must cover, for every rival, the price of the rival's distinctiveness from it.

Some notation.  Let $Z:=O\times R$ and write $K_Z(z\mid a):=K(o,r\mid a)$ for $z=(o,r)$; let
\[
 \bar u(a):=\sum_{o,r}K(o,r\mid a)u(o),
 \qquad
 m_{q,K}(z):=\sum_aq(a)K_Z(z\mid a)
\]
be the expected payoff of an action and the distribution of visible consequences induced by a lottery.  The
Kullback--Leibler divergence $\KL(P\,\|\,Q):=\sum_zP(z)\log\bigl(P(z)/Q(z)\bigr)$ between distributions on $Z$
measures how sharply consequences drawn from $P$ stand out against a background of $Q$; it equals $+\infty$ when
$P$ puts mass where $Q$ does not.

\begin{proposition}[Exercised distinctiveness]\label{cor:support-pure}
Let $Z_K:=\{z\in Z:K_Z(z\mid a)>0\text{ for some }a\in A\}$.
\begin{enumerate}[label=\textup{(\roman*)},leftmargin=*]
\item Every optimal lottery $q^*$ satisfies $m_{q^*,K}(z)>0$ for every $z\in Z_K$.  Consequently, if every action $a$ has
a private consequence---some $z_a$ with $K_Z(z_a\mid a)>0$ and $K_Z(z_a\mid b)=0$ for all $b\ne a$---then every
optimal lottery has full support.
\item The pure lottery $\delta_a$ is optimal if and only if
\[
 \bar u(b)+\lambda
 \KL\!\left(K_Z(\cdot\mid b)\,\middle\|\,K_Z(\cdot\mid a)\right)
 \le \bar u(a)
 \qquad\text{for every }b\in A,
\]
and uniquely optimal if these inequalities are strict for $b\ne a$.  In particular, if some rival reaches a
consequence outside the support of $K_Z(\cdot\mid a)$, then $\delta_a$ is not optimal, however large its payoff
advantage.
\end{enumerate}
\end{proposition}

With two actions, the threshold below which the agent mixes is a case of
Proposition~\ref{cor:support-pure}\textup{(ii)}.  Let $A=\{a,b\}$, write $K_a:=K_Z(\cdot\mid a)$ and
$K_b:=K_Z(\cdot\mid b)$, and let $\Delta:=\bar u(a)-\bar u(b)\ge0$ be the payoff advantage of $a$.  Then
$\delta_a$ is optimal if and only if $\Delta\ge\lambda\,\KL(K_b\,\|\,K_a)$.  The threshold is directional: what
protects the pure action is the detectability of the rival against the background of its own signature.  It is
infinite when $K_b$ puts mass where $K_a$ does not---the two-action face of
Proposition~\ref{cor:support-pure}\textup{(i)}.

Below the threshold the agent mixes (Proposition~\ref{prop:binary}, Appendix~\ref{app:choice-proofs}).  The
optimal mix is unique, and it puts more mass on the inferior action the larger is $\lambda$ and less the larger
is $\Delta$.  Mixing therefore stops exactly at $\Delta=\lambda\,\KL(K_b\,\|\,K_a)$, and since the divergence is
known once the channel is designed, a single switching experiment identifies $\lambda$.  At equal payoffs the
mix is the capacity-achieving input of the channel, whatever $\lambda$: an agent who values trace at all,
however little, randomises at zero stakes so as to maximise the information the visible consequence carries
about her act.  The comparative statics restate the wedge of the trace test (\S\ref{sec:tracetest}):
randomisation here is instrumental to trace, so it strengthens with the informativeness of the channel and
vanishes on uninformative ones, whereas a taste for mixing as such is channel-blind.

With more than two actions, the same logic is captured by a first-order condition.  The objective
$q\mapsto V(q,K)$ is concave in $q$, so an optimum exists and Kuhn--Tucker conditions characterise it: at an
optimum, every action in use offers the same total of expected payoff and priced distinctiveness.

\begin{proposition}[Distinctiveness condition]\label{prop:optimal-choice}
Let $\lambda>0$.  A lottery $q^*$ is optimal if and only if there is a constant $c$ such that
\[
 \bar u(a)+\lambda\KL\!\left(K_Z(\cdot\mid a)\,\middle\|\,m_{q^*,K}\right)=c
 \quad\text{for }a\in\supp(q^*),
 \qquad
 \le c\quad\text{for }a\notin\supp(q^*).
\]
\end{proposition}

The proposition gives a \emph{distinctiveness premium}.  Each action is credited, on top of its expected payoff,
with $\lambda$ times how far its visible consequences depart from the average visible behaviour that the lottery
itself generates; at an optimum this total is equalised across the actions in use.  An action can therefore receive
more probability because it pays more or because its consequences are more distinctive.  The
condition mirrors rational-inattention first-order conditions with the channel direction and sign reversed
\cite{Sims2003,MatejkaMcKay2015}.

Which lotteries are optimal depends on the coupling of actions to consequences, not on consequences alone; and
optimal lotteries need not be unique.  Yet every optimal lottery induces the same distribution of visible
consequences.  The model therefore predicts the record an observer sees, not the mix that generated it.

\begin{proposition}[Optimal marginal and multiplicity]\label{prop:optimal-marginal}
Write $P_a:=K_Z(\cdot\mid a)$ and $C^*(K):=\max_{q\in\D(A)}V(q,K)$.  There is a unique $m^*\in\D(Z)$ such that
$m_{q^*,K}=m^*$ for every optimal lottery $q^*$.  With
\[
 \mathcal A^*
 :=\bigl\{a\in A:
 \bar u(a)+\lambda
 \KL\!\left(P_a\,\middle\|\,m^*\right)=C^*(K)
 \bigr\},
\]
the set of optimal lotteries is exactly
\[
 \bigl\{q\in\D(A):
 m_{q,K}=m^*,\ \supp(q)\subseteq\mathcal A^*
 \bigr\},
\]
and the optimum is unique whenever the rows $\{P_a:a\in\mathcal A^*\}$ are affinely independent.  If actions are
partitioned into classes with identical rows, $V(q,K)$ depends on $q$ within each class only through the class's
total mass: every redistribution of an optimal lottery within such a class is optimal, and the optimal lotteries are
exactly the lifts of the optimal lotteries of the channel that collapses the classes.
\end{proposition}

The proposition's final claim---that moving mass between actions with identical rows changes nothing---is a
test.  Split an action into two copies with the same row.  The record cannot tell the copies apart, so the
value, the induced marginal $m^*$, and the pair's total probability do not depend on how the mass is divided
between them.  Any model that gives each action its own weight, as multinomial logit does, predicts instead that
duplication raises the pair's combined share.  Trace-tempered choice predicts exact invariance.

\begin{remark}[Logit boundary]\label{rem:logit}
Suppose every action leaves its own mark: the supports of the rows are pairwise disjoint, so the visible
consequence identifies the action and $\I_{qK}(A;O,R)=\Sh(q)$.  The unique optimum is then multinomial logit,
\[
 q^*(a)
 =\frac{\exp\bigl(\bar u(a)/\lambda\bigr)}
 {\sum_{a'}\exp\bigl(\bar u(a')/\lambda\bigr)},
\]
with $1/\lambda$ as the payoff sensitivity (information in natural logarithms; a change of base rescales
$\lambda$).  If instead no action leaves any mark---identical rows---every lottery is optimal.  Between the two
extremes shares solve the equalisation condition of Proposition~\ref{prop:optimal-choice}, in which
distinctiveness is measured against the average behaviour the optimum itself generates; there is no closed form.
Multinomial logit is thus the fully revealing special case of trace-tempered choice, not its general prediction.
In rational inattention the same formula arises for the opposite reason, as the optimal tilting of a chosen
channel at a fixed prior over states \cite{MatejkaMcKay2015}; here the channel is fixed and the source is chosen.
\end{remark}

\subsection{The price of being on the record}\label{sec:record-price}

Because Theorem~\ref{thm:main} represents comparisons across environments by one $V$, the model prices not only
lotteries but changes to the record technology itself: what is it worth to the agent that her coin is observed,
that her record is coarsened, that a further signal about her is disclosed?  Each answer below is an exact amount
of information multiplied by $\lambda$, and each corresponds to a block comparison in the primitive domain.  The
first result concerns the mixing device.

\begin{corollary}[The recorded coin]\label{cor:recorded-coin}
Let $K,L:A\to\D(O\times R)$ and $t\in(0,1)$.  Let $tK\oplus(1-t)L$ be the public mixture that draws an
action-independent coin $Y$ with $\Pr(Y=0)=t$, applies $K$ if $Y=0$ and $L$ otherwise, and includes the coin in
the record; let $M_t:=tK+(1-t)L$ be the rowwise mixture, which hides it.  Then, for every $q\in\D(A)$,
\begin{enumerate}[label=\textup{(\roman*)},leftmargin=*]
\item $V\bigl(q,\,tK\oplus(1-t)L\bigr)=t\,V(q,K)+(1-t)\,V(q,L)$;
\item $V\bigl(q,\,tK\oplus(1-t)L\bigr)-V(q,M_t)=\lambda\,\I(A;Y\mid O,R)\ge0$, with equality if and only if $A$
and $Y$ are conditionally independent given $(O,R)$.
\end{enumerate}
\end{corollary}

\begin{proof}
Part \textup{(i)} is the chain rule with an action-independent first stage: $\I(A;Y)=0$ and each reached branch
reproduces the joint law of the corresponding constituent.  Part \textup{(ii)} follows from Proposition~\ref{prop:garbling-price} below, since deleting the coin coordinate is a deterministic payoff-preserving
record processor carrying the public mixture to $M_t$.
\end{proof}

Part \textup{(i)} is the value-level face of derived public-coin independence (Lemma~\ref{lem:public-coin}): the
agent has no intrinsic taste for or against randomising over environments, so a publicly randomised environment is
never ranked strictly above or below both of its constituents.  Part \textup{(ii)} says what she does care about:
whether the coin is on the record.  Hiding the randomisation device costs $\lambda\,\I(A;Y\mid O,R)$, which is
strictly positive whenever knowing the operating regime would sharpen inference about the action given the visible
consequences---even though the coin is independent of the action ex ante.  A trace-tempered agent weakly prefers
her mixing device publicised; theories that attach value to randomisation as such are silent on this visibility
margin.

Behind part \textup{(ii)} is the exact price of any record garbling.  Axiom~\textup{(A4)} asserts only a weak
preference against garbling; the representation upgrades the axiom to an exact price and characterises when
the price is zero.

\begin{proposition}[Price of a record garbling]\label{prop:garbling-price}
Let $K:A\to\D(O\times R)$, let $T(r'\mid o,r)$ be a record processor as in \eqref{eq:record-processing}, and write
$L:=KT$.  Fix $q\in\D(A)$ and let $(A,O,R,R')$ carry the joint law $q(a)K(o,r\mid a)T(r'\mid o,r)$.  Then
\[
 V(q,K)-V(q,L)=\lambda\,\I\!\left(A;R\mid O,R'\right)\ge0,
\]
and the following are equivalent:
\textup{(a)} $V(q,L)=V(q,K)$;
\textup{(b)} $A$ and $(O,R)$ are conditionally independent given $(O,R')$;
\textup{(c)} $\mu_{qL}=\mu_{qK}$.
\end{proposition}

Rewriting the record costs $\lambda$ times the information about the action that the original record retains beyond
the rewritten one, and a garbling is exactly free when it pools only visible consequences that already induced the
same posterior about the action: an agent with a trace motive does not pay for gratuitous detail, only for detail
that discriminates among her own actions.  Symmetrically, appending to the record any further signal
$S$ drawn given $(a,o,r)$ raises $V(q,\cdot)$ by exactly $\lambda\,\I(A;S\mid O,R)\ge0$: more evidence about
oneself never hurts, a conclusion the deniability dual of Corollary~\ref{cor:signduality} reverses sign for sign.
Since the visible pair always includes the payoff coordinate, the total value of trace in any environment is
bounded by $\lambda\min\{\Sh(q),\log(|O|\,|R|)\}$.

These prices invite measurement: how much payoff would the agent give up for a more revealing record?  The
difficulty is that the domain has no money---the finite outcome set contains no outcome worth $u(o)+\tau$---so a
compensating offset must be manufactured from the two outcomes that Axiom~\textup{(A7)} strictly ranks, using
comparisons that stay inside the primitive domain.  Let $o^+$ and $o^-$ attain the
maximum and minimum of $u$, write $\Delta u:=u(o^+)-u(o^-)>0$, and for $\beta\in[0,1]$ let
$w(\beta):=\beta u(o^+)+(1-\beta)u(o^-)$.

\begin{definition}[Sweetened dilution]\label{def:dilution}
For $K:A\to\D(O\times R)$, $\varepsilon\in[0,1)$, and $\beta\in[0,1]$, let $P_\varepsilon:A\to\D(\{0,1\})$ be the
public coin $P_\varepsilon(1\mid a)=\varepsilon$, and let $B_\beta:A\to\D(O\times\{\ast\})$ have the constant rows
$B_\beta(o^+,\ast\mid a)=\beta$ and $B_\beta(o^-,\ast\mid a)=1-\beta$.  The
\emph{$(\varepsilon,\beta)$-sweetened dilution} of $K$ is the compound channel
\[
 D_{\varepsilon,\beta}K:=P_\varepsilon\triangleright\{K,\,B_\beta\},
\]
as in \eqref{eq:compound}, listing the branch-$0$ and branch-$1$ continuations in order.
\end{definition}

\begin{proposition}[Willingness to pay for trace]\label{prop:wtp}
Fix $q\in\D(A)$, $p\in\D(B)$, and channels $K:A\to\D(O\times R)$ and $L:B\to\D(O\times S)$.
\begin{enumerate}[label=\textup{(\roman*)},leftmargin=*]
\item For every $\varepsilon\in[0,1)$ and $\beta\in[0,1]$,
\[
 V\bigl(q,\,D_{\varepsilon,\beta}K\bigr)=(1-\varepsilon)\,V(q,K)+\varepsilon\,w(\beta).
\]
\item If $\varepsilon\in(0,1)$ and $\beta,\beta'\in[0,1]$ satisfy
$\bigl(q,D_{\varepsilon,\beta}K\bigr)\sim\bigl(p,D_{\varepsilon,\beta'}L\bigr)$, then
\[
 \frac{\varepsilon}{1-\varepsilon}\bigl[w(\beta)-w(\beta')\bigr]
 =V(p,L)-V(q,K).
\]
When the material parts agree, $\E_{q,K}[u(O)]=\E_{p,L}[u(O)]$, the elicited offset equals
$\lambda\bigl[\I_{pL}(B;O,S)-\I_{qK}(A;O,R)\bigr]$: information is priced linearly at rate $\lambda$.
\item Such a pair $(\beta,\beta')$ exists if and only if
$\lvert V(p,L)-V(q,K)\rvert\le\frac{\varepsilon}{1-\varepsilon}\Delta u$; a boundary response therefore brackets
the value difference instead of measuring it.
\end{enumerate}
\end{proposition}

The public coin is incentive-neutral: it rescales both blocks' values identically, so the elicited odds gap
measures only the value difference, whatever the dilution.  The price it recovers is, moreover, unique and
portable.

\begin{proposition}[Single information price]\label{prop:single-price}
Write $E_{qK}:=\E_{q,K}[u(O)]$ and $\phi(q,K):=\bigl(\I_{qK}(A;O,R),\,E_{qK}\bigr)$, and let
$\underline u:=\min_ou(o)$ and $\overline u:=\max_ou(o)$.
\begin{enumerate}[label=\textup{(\roman*)},leftmargin=*]
\item The image of $\phi$ over all finite pairs $(q,K)$ is exactly $[0,\infty)\times[\underline u,\overline u]$, and
the indifference classes of the block-comparison relation are, in these coordinates, the nonempty intersections of
the lines $E=c-\lambda\I$ with that strip.
\item If $(q,K)\sim(p,L)$ and $\I_{qK}\ne\I_{pL}$, then
\[
 \frac{E_{pL}-E_{qK}}{\I_{qK}-\I_{pL}}=\lambda,
\]
and $\lambda$ is the unique number with this property: for every $\lambda'\ne\lambda$ there are finite pairs
$(q,K)\sim(p,L)$ with $E_{qK}+\lambda'\I_{qK}\ne E_{pL}+\lambda'\I_{pL}$.
\end{enumerate}
\end{proposition}

At any fixed calibration of payoff units, the cross-environment indifference map in the information--payoff plane
consists of parallel straight lines of slope $-\lambda$: bits carry one price, the same at every prior, every
information level, and in every finite environment, as Proposition~\ref{prop:uniqueness} transports it.  This is
the sharpest testable fingerprint of branchwise continuation consistency.  The satiation criterion of
Proposition~\ref{prop:independence}, which fails exactly Axiom~\textup{(A6)}, produces kinked indifference curves
whose measured slope collapses to zero beyond the cap; constancy of the elicited slope across designs with
different priors or alphabets is therefore a direct over-identification test.
Appendix~\ref{app:choice-proofs} records a complete two-treatment protocol implementing this elicitation
(Remark~\ref{rem:elicitation}), with the information gap fixed by design and the garbling control built in.

Choices, and not only indifferences, bound the same price.  Two behaviours ordered by $\lambda$ are ordered in
what they choose: a higher weight buys more information and less payoff, and the rate of exchange between the two
brackets the weight itself.

\begin{remark}[Trace demand]\label{rem:trace-demand}
Fix $K$, write $e(q):=\E_{q,K}[u(O)]$ and $i(q):=\I_{qK}(A;O,R)$, and let $q_j$ maximise $e+\lambda_ji$ for
$0<\lambda_1<\lambda_2$.  Comparing each optimum at the other weight gives $i_1\le i_2$, $e_1\ge e_2$, and,
whenever behaviour responds,
\[
 \lambda_1\le\frac{e_1-e_2}{i_2-i_1}\le\lambda_2:
\]
the payoff surrendered per additional unit of information brackets the trace weight, whatever selection among
optima the data reflect.  The indirect value is convex in $\lambda$, with the chosen information as its derivative
wherever it is differentiable, and its extreme regimes are expected utility ($\lambda\downarrow0$, with information
breaking material ties) and maximal distinctiveness ($\lambda\uparrow\infty$, with payoff breaking information
ties).  Exact statements are Proposition~\ref{prop:trace-demand} and Corollary~\ref{cor:lambda-limits} in
Appendix~\ref{app:choice-proofs}.
\end{remark}

\subsection{What preferences cannot reveal}\label{sec:cannot-reveal}

Suppose an agent ranks environment $K$ above environment $L$ at \emph{every} lottery.  Has she thereby revealed
that $L$ is a coarsening of $K$---that some garbling maps the one record technology into the other?  For a
Blackwell decision-maker the answer would be yes: dominance in every decision problem is equivalent to garbling
\cite{Blackwell1953}.  Here the answer is no, and the gap can be stated exactly.

\begin{definition}\label{def:channel-orders}
Let $K:A\to\D(O\times R)$ and $L:A\to\D(O\times R')$ have identical payoff kernels,
$\sum_rK(o,r\mid a)=\sum_{r'}L(o,r'\mid a)$ for every $(a,o)$.  Write $K\succeq_gL$ if $L=KT$ for some record
processor $T$ as in \eqref{eq:record-processing}, and $K\succeq_cL$ if
$\I_{qK}(A;O,R)\ge\I_{qL}(A;O,R')$ for every $q\in\D(A)$; the latter is the \emph{more capable} order on the
visible-consequence channels \cite{KornerMarton1977}.
\end{definition}

\begin{proposition}[Uniform dominance is more-capable, not garbling]\label{prop:blackwell}
Let $K$ and $L$ have identical payoff kernels.
\begin{enumerate}[label=\textup{(\roman*)},leftmargin=*]
\item $K\succeq_gL\Rightarrow K\succeq_cL$, and $K\succeq_cL$ holds if and only if $V(q,K)\ge V(q,L)$ for every
$q\in\D(A)$.
\item The converse of the first implication fails: there are channels with identical payoff kernels such that
$V(q,K)\ge V(q,L)$ for every $q$ and every representation $(u,\lambda)$, yet $L\ne KT$ for every record processor
$T$.
\end{enumerate}
\end{proposition}

The reason Blackwell's equivalence fails here is that his test family is rich---all state-dependent utilities---
while the axioms pin the attitude to records down to the single functional $\E[u]+\lambda\I$, so environments are
assessed only through mutual information at priors.  The revealed informativeness order is then exactly the
information theorists' more-capable order, strictly coarser than garbling.  And the counterexample is
utility-robust---the payoff terms cancel identically---so observing that an agent always ranks one environment
above another licenses no inference that the second is a censored version of the first.


% ===============================================================
\section{Interpretation and scope}
% ===============================================================

\paragraph{Evidence, not authorship.}
The name promises agency; the functional rewards uncertainty.  An agent certain of her act earns zero trace value,
however perfectly the record attributes it, while one who hands the decision to a fair coin earns the most the
channel allows.  We concede what this invites: the theory prices evidence---the news a visible consequence carries
about conduct---and news requires antecedent doubt, which is Dostoevsky's own logic.  The zero at certainty is a
discriminating prediction rather than an embarrassment.  Image models value the level of an observer's regard, so
anticipated attribution keeps its worth for an agent certain of her act; here trace dies with the doubt.  Two
readings of $q$ are consistent with this: a mixing device the agent operates, which makes the randomisation results
of Section~\ref{sec:implications} literal, or an observer's prior or population frequency against which a realised act is evaluated.
Under the second reading $q$ is exogenous, and the theorem is silent on how the observer forms it.

\paragraph{Trace or control.}
Mutual information is symmetric, and the domain inherits the symmetry.  A taste for visible consequences that reveal
the act is behaviourally indistinguishable from a taste for acts that statistically shape visible consequences: the
theorem pins down the value of the coupling between action and consequence, and trace and control are two readings
of that one object.  No data in a channel-indexed domain can separate them, much as Savage's data cannot separate
belief from state-dependent taste.  Separation would require interventional primitives---an audience that can be
removed while the channel is left intact---and these lie outside the theory.  Every choice implication in
Section~\ref{sec:implications} therefore tests both readings at once.

\paragraph{Empowerment.}
At constant material utility, maximising $V$ yields the Shannon capacity of the visible-consequence channel: the
empowerment index of artificial-intelligence and complex-systems research
\cite{KlyubinPolaniNehaniv2005,SalgeGlackinPolani2014}.  The connection is stronger than a shared optimum.  The
theorem gives a preference foundation for the entire map $(q,K)\mapsto I_{q,K}$, so it covers constrained behaviour
and choices tempered by material payoffs, and it identifies $\lambda$, the payoff price of a unit of empowerment.
The connection is also narrow in an instructive way: the criterion depends only on actions assigned positive
probability, so it captures control that is exercised, not options held in reserve---a genuine option value would
require preferences over menus, outside this domain.  Each side gains something: agency research acquires a
behavioural test for its objective, and decision theory acquires an objective already in use in another field.

\paragraph{A belief-based utility, located.}
At a fixed prior the trace term is a belief-based utility: the expected reduction in the entropy of the posterior
about the act, a posterior-separable felicity in the manner of \cite{CaplinLeahy2001}---though the belief concerns
the agent's own conduct, not an external state.  What the axioms add is the location of the functional form in
behaviour.  Entropy is nowhere assumed; branchwise continuation consistency carries the whole cardinal content, and
it is a substantive claim, not a convenience.  Its fingerprint is that bits carry one price: cross-environment
indifference curves in the information--payoff plane are parallel straight lines
(Proposition~\ref{prop:single-price}).  An agent whose taste for trace satiates violates exactly this axiom,
produces kinked curves, and satisfies every other.  The theorem's scope is agents for whom that constancy of price
holds.

\paragraph{One-shot, with no audience in the model.}
Beliefs enter only as the yardstick for how much the realised consequence narrows uncertainty about the action.  No
observer responds: the representation predicts neither praise nor sanction nor any return to reputation, and the
staged records of Axiom~\textup{(A6)} introduce sequencing of evidence, not repeated interaction.  This is the
right first step, not an oversight: it isolates the intrinsic value of trace, so that a strategic or dynamic
extension can add an audience's response without confusing it with the motive identified here.

\paragraph{The split is scaffolding.}
The payoff--record split is scaffolding, not substance.  Identification needs two independently variable directions
and a temporal order: records may be garbled while payoff stays intact, and terminal payoff anchors a common scale.
But the object valued is the undivided pair; the theorem separates $O$ from $R$ only to recombine them in
$I(A;O,R)$.  The Underground Man, on this reading, was not asking to be rid of reason: he was asking for a lottery
genuinely open beforehand and a consequence able to tell his acts apart, and the cakes he stood ready to risk are
what $\lambda$ measures.

\begin{thebibliography}{99}
\raggedright

\bibitem{AczelDaroczy1975}
J\'anos Acz\'el and Zolt\'an Dar\'oczy.
\newblock {\em On Measures of Information and Their Characterizations}.
\newblock Academic Press, New York, 1975.

\bibitem{AgranovOrtoleva2017}
Marina Agranov and Pietro Ortoleva.
\newblock Stochastic choice and preferences for randomization.
\newblock {\em Journal of Political Economy}, 125(1):40--68, 2017.

\bibitem{AndreoniBernheim2009}
James Andreoni and B. Douglas Bernheim.
\newblock Social image and the 50--50 norm: A theoretical and experimental analysis of audience effects.
\newblock {\em Econometrica}, 77(5):1607--1636, 2009.

\bibitem{AnscombeAumann1963}
Francis J. Anscombe and Robert J. Aumann.
\newblock A definition of subjective probability.
\newblock {\em Annals of Mathematical Statistics}, 34(1):199--205, 1963.

\bibitem{AzrieliLehrer2008}
Yaron Azrieli and Ehud Lehrer.
\newblock The value of a stochastic information structure.
\newblock {\em Games and Economic Behavior}, 63(2):679--693, 2008.

\bibitem{BaezFritz2014}
John C. Baez and Tobias Fritz.
\newblock A {B}ayesian characterization of relative entropy.
\newblock {\em Theory and Applications of Categories}, 29(16):421--456, 2014.

\bibitem{BartlingFehrHerz2014}
Bj\"orn Bartling, Ernst Fehr, and Holger Herz.
\newblock The intrinsic value of decision rights.
\newblock {\em Econometrica}, 82(6):2005--2039, 2014.

\bibitem{BaezFritzLeinster2011}
John C. Baez, Tobias Fritz, and Tom Leinster.
\newblock A characterization of entropy in terms of information loss.
\newblock {\em Entropy}, 13(11):1945--1957, 2011.

\bibitem{BenabouTirole2002}
Roland B\'enabou and Jean Tirole.
\newblock Self-confidence and personal motivation.
\newblock {\em Quarterly Journal of Economics}, 117(3):871--915, 2002.

\bibitem{BenabouTirole2006}
Roland B\'enabou and Jean Tirole.
\newblock Incentives and prosocial behavior.
\newblock {\em American Economic Review}, 96(5):1652--1678, 2006.

\bibitem{Bernheim1994}
B. Douglas Bernheim.
\newblock A theory of conformity.
\newblock {\em Journal of Political Economy}, 102(5):841--877, 1994.

\bibitem{Blackwell1953}
David Blackwell.
\newblock Equivalent comparisons of experiments.
\newblock {\em Annals of Mathematical Statistics}, 24(2):265--272, 1953.

\bibitem{BuffatPraxmarerSutter2023}
Justin Buffat, Matthias Praxmarer, and Matthias Sutter.
\newblock The intrinsic value of decision rights: A replication and an extension to team decision making.
\newblock {\em Journal of Economic Behavior \& Organization}, 209:560--571, 2023.

\bibitem{CabralesGossnerSerrano2013}
Antonio Cabrales, Olivier Gossner, and Roberto Serrano.
\newblock Entropy and the value of information for investors.
\newblock {\em American Economic Review}, 103(1):360--377, 2013.

\bibitem{CaplinDeanLeahy2022}
Andrew Caplin, Mark Dean, and John Leahy.
\newblock Rationally inattentive behavior: Characterizing and generalizing {S}hannon entropy.
\newblock {\em Journal of Political Economy}, 130(6):1676--1715, 2022.

\bibitem{CaplinLeahy2001}
Andrew Caplin and John Leahy.
\newblock Psychological expected utility theory and anticipatory feelings.
\newblock {\em Quarterly Journal of Economics}, 116(1):55--79, 2001.

\bibitem{CerreiaVioglioDillenbergerOrtolevaRiella2019}
Simone Cerreia-Vioglio, David Dillenberger, Pietro Ortoleva, and Gil Riella.
\newblock Deliberately stochastic.
\newblock {\em American Economic Review}, 109(7):2425--2445, 2019.

\bibitem{CoverThomas2006}
Thomas M. Cover and Joy A. Thomas.
\newblock {\em Elements of Information Theory}.
\newblock Wiley-Interscience, Hoboken, second edition, 2006.

\bibitem{Csiszar2008}
Imre Csisz\'ar.
\newblock Axiomatic characterizations of information measures.
\newblock {\em Entropy}, 10(3):261--273, 2008.

\bibitem{DekelLipmanRustichini2001}
Eddie Dekel, Barton L. Lipman, and Aldo Rustichini.
\newblock Representing preferences with a unique subjective state space.
\newblock {\em Econometrica}, 69(4):891--934, 2001.

\bibitem{Denti2022}
Tommaso Denti.
\newblock Posterior separable cost of information.
\newblock {\em American Economic Review}, 112(10):3215--3259, 2022.

\bibitem{Dostoevsky1864}
Fyodor Dostoevsky.
\newblock {\em Notes from the Underground}, Part I.
\newblock 1864.  Translated by Constance Garnett, 1918.

\bibitem{ElyFrankelKamenica2015}
Jeffrey Ely, Alexander Frankel, and Emir Kamenica.
\newblock Suspense and surprise.
\newblock {\em Journal of Political Economy}, 123(1):215--260, 2015.

\bibitem{Faddeev1956}
D. K. Faddeev.
\newblock On the concept of entropy of a finite probabilistic scheme.
\newblock {\em Uspekhi Matematicheskikh Nauk}, 11(1):227--231, 1956.
\newblock In Russian.

\bibitem{FrankelKamenica2019}
Alexander Frankel and Emir Kamenica.
\newblock Quantifying information and uncertainty.
\newblock {\em American Economic Review}, 109(10):3650--3680, 2019.

\bibitem{GilboaLehrer1991}
Itzhak Gilboa and Ehud Lehrer.
\newblock The value of information---an axiomatic approach.
\newblock {\em Journal of Mathematical Economics}, 20(5):443--459, 1991.

\bibitem{GlazerKonrad1996}
Amihai Glazer and Kai A. Konrad.
\newblock A signaling explanation for charity.
\newblock {\em American Economic Review}, 86(4):1019--1028, 1996.

\bibitem{GrantKajiiPolak1998}
Simon Grant, Atsushi Kajii, and Ben Polak.
\newblock Intrinsic preference for information.
\newblock {\em Journal of Economic Theory}, 83(2):233--259, 1998.

\bibitem{GulPesendorfer2001}
Faruk Gul and Wolfgang Pesendorfer.
\newblock Temptation and self-control.
\newblock {\em Econometrica}, 69(6):1403--1435, 2001.

\bibitem{HersteinMilnor1953}
I. N. Herstein and John Milnor.
\newblock An axiomatic approach to measurable utility.
\newblock {\em Econometrica}, 21(2):291--297, 1953.

\bibitem{Hobson1969}
Arthur Hobson.
\newblock A new theorem of information theory.
\newblock {\em Journal of Statistical Physics}, 1(3):383--391, 1969.

\bibitem{Khinchin1957}
Aleksandr I. Khinchin.
\newblock {\em Mathematical Foundations of Information Theory}.
\newblock Dover, New York, 1957.

\bibitem{KlyubinPolaniNehaniv2005}
Alexander S. Klyubin, Daniel Polani, and Chrystopher L. Nehaniv.
\newblock Empowerment: A universal agent-centric measure of control.
\newblock In {\em Proceedings of the 2005 IEEE Congress on Evolutionary Computation}, volume~1, pages 128--135. IEEE,
  2005.

\bibitem{KornerMarton1977}
J\'anos K\"orner and Katalin Marton.
\newblock Comparison of two noisy channels.
\newblock In Imre Csisz\'ar and Peter Elias, editors, {\em Topics in Information Theory}, Colloquia Mathematica
  Societatis J\'anos Bolyai, volume~16, pages 411--423. North-Holland, Amsterdam, 1977.

\bibitem{Kreps1979}
David M. Kreps.
\newblock A representation theorem for preference for flexibility.
\newblock {\em Econometrica}, 47(3):565--577, 1979.

\bibitem{KrepsPorteus1978}
David M. Kreps and Evan L. Porteus.
\newblock Temporal resolution of uncertainty and dynamic choice theory.
\newblock {\em Econometrica}, 46(1):185--200, 1978.

\bibitem{Liljeholm2022}
Mimi Liljeholm.
\newblock Flexible control as surrogate reward or dynamic reward maximization.
\newblock {\em Cognition}, 229:105262, 2022.

\bibitem{MatejkaMcKay2015}
Filip Mat\v{e}jka and Alisdair McKay.
\newblock Rational inattention to discrete choices: A new foundation for the multinomial logit model.
\newblock {\em American Economic Review}, 105(1):272--298, 2015.

\bibitem{MistryLiljeholm2016}
Prachi Mistry and Mimi Liljeholm.
\newblock Instrumental divergence and the value of control.
\newblock {\em Scientific Reports}, 6:36295, 2016.

\bibitem{NortonLiljeholm2020}
Kaitlyn G. Norton and Mimi Liljeholm.
\newblock The rostrolateral prefrontal cortex mediates a preference for high-agency environments.
\newblock {\em Journal of Neuroscience}, 40(22):4401--4409, 2020.

\bibitem{OwensGrossmanFackler2014}
David Owens, Zachary Grossman, and Ryan Fackler.
\newblock The control premium: A preference for payoff autonomy.
\newblock {\em American Economic Journal: Microeconomics}, 6(4):138--161, 2014.

\bibitem{PattanaikXu1990}
Prasanta K. Pattanaik and Yongsheng Xu.
\newblock On ranking opportunity sets in terms of freedom of choice.
\newblock {\em Recherches \'Economiques de Louvain}, 56(3--4):383--390, 1990.

\bibitem{PomattoStrackTamuz2023}
Luciano Pomatto, Philipp Strack, and Omer Tamuz.
\newblock The cost of information: The case of constant marginal costs.
\newblock {\em American Economic Review}, 113(5):1360--1393, 2023.

\bibitem{Renyi1961}
Alfr\'ed R\'enyi.
\newblock On measures of entropy and information.
\newblock In {\em Proceedings of the Fourth Berkeley Symposium on Mathematical Statistics and Probability}, volume~1,
  pages 547--561. University of California Press, 1961.

\bibitem{SalgeGlackinPolani2014}
Christoph Salge, Cornelius Glackin, and Daniel Polani.
\newblock Empowerment---an introduction.
\newblock In Mikhail Prokopenko, editor, {\em Guided Self-Organization: Inception}, pages 67--114. Springer, Berlin,
  2014.

\bibitem{Shannon1948}
Claude E. Shannon.
\newblock A mathematical theory of communication.
\newblock {\em Bell System Technical Journal}, 27:379--423, 623--656, 1948.

\bibitem{Sims2003}
Christopher A. Sims.
\newblock Implications of rational inattention.
\newblock {\em Journal of Monetary Economics}, 50(3):665--690, 2003.

\bibitem{Spence1973}
Michael Spence.
\newblock Job market signaling.
\newblock {\em Quarterly Journal of Economics}, 87(3):355--374, 1973.

\bibitem{WynerZiv1973}
Aaron D. Wyner and Jacob Ziv.
\newblock A theorem on the entropy of certain binary sequences and applications: Part {I}.
\newblock {\em IEEE Transactions on Information Theory}, 19(6):769--772, 1973.

\end{thebibliography}



\clearpage
\appendix

\part*{Proofs of the main results}
\addcontentsline{toc}{part}{Proofs of the main results}

\section{Proof of the representation theorem}\label{app:main-proof}
The sufficiency proof is modular.  The first lemma removes inessential record labels.  The second records the precise
consequence of Axiom~\textup{(A6)} that was an axiom in many earlier formulations.  The next lemmas calibrate material
payoffs, import the pure-trace engine, and align continuation scales.  Sequentialising an arbitrary joint channel then
yields the representation in one line.

\paragraph{Proof notation.}
For a prior $q$ and joint channel $K$, define the \emph{terminal law}
\begin{equation}
 \eta_{qK}:=\sum_{(o,r):p_{qK}(o,r)>0}
 p_{qK}(o,r)\,\delta_{(o,\pi^{qK}_{or})}.
 \label{eq:terminal-law}
\end{equation}
It records the payoff and the posterior about the realised action after everything visible has been observed; its
posterior marginal is $\mu_{qK}$, as defined immediately after the main theorem.  Bayes plausibility gives
\[
 \int_{O\times\D(A)}\pi\,d\eta_{qK}(o,\pi)=q.
\]
Let $\mathcal E_q$ be the set of finite-support laws $\eta$ on $O\times\D(A)$ satisfying this barycentre condition.
Every $\eta=\sum_i m_i\delta_{(o_i,\pi_i)}\in\mathcal E_q$ is implementable: take records $i$ and, for $q(a)>0$,
set
\begin{equation}
 K_\eta(o,i\mid a)
 :=\mathbf 1_{\{o=o_i\}}\frac{m_i\pi_i(a)}{q(a)}.
 \label{eq:implementation}
\end{equation}
Bayes plausibility makes the rows sum to one; null rows may be completed arbitrarily.  Thus terminal laws range over
the entire finite Bayes-plausible domain.

For channels $K,L$ on the same action set and $t\in[0,1]$, let $tK\oplus(1-t)L$ denote the channel that draws a
public coin independent of the action, applies the selected channel, and includes the coin in the record.  Then
\begin{equation}
 \eta_{q,tK\oplus(1-t)L}=t\eta_{qK}+(1-t)\eta_{qL}.
 \label{eq:public-mixture}
\end{equation}

\subsection{Terminal-law sufficiency and a derived public coin}

\begin{lemma}[Terminal-law sufficiency]\label{lem:terminal-suff}
Assume Axioms \textup{(A1), (A3), (A4)}.  Fix a full-support prior $q$ on $A$.  If
$\eta_{qK}=\eta_{qL}$, then $(q,K)\sim(q,L)$.  Boundary priors satisfy the same conclusion after deleting null
actions under Axiom~\textup{(A5)}.
\end{lemma}

\begin{proof}
For every reached $(o,r)$, replace $r$ by the posterior $\pi^{qK}_{or}$.  This is a deterministic record processor
allowed to depend on $o$.  The resulting canonical channel is
\[
 K^{\eta}(o,\pi\mid a)
 =\frac{\eta_{qK}(o,\pi)\pi(a)}{q(a)}.
\]
Conversely, conditional on $(o,\pi)$, the canonical record can draw an original record $r$ with probability
$p_{qK}(o,r)/\eta_{qK}(o,\pi)$ among those records inducing $\pi$.  This draw is independent of $a$: Bayes' formula
gives $q(a)K(o,r\mid a)=p_{qK}(o,r)\pi(a)$.  Hence $K$ and $K^\eta$ are record processings of one another.  Axiom
\textup{(A4)} ranks each weakly above the other; weak order and block coherence give indifference.  The same canonical
channel is obtained from $L$ when the terminal laws agree.  For a boundary $q$, Axiom~\textup{(A5)} projects to
$\supp q$ and embeds back, making null rows irrelevant.
\end{proof}

\begin{corollary}[Derived pure-action consequentialism]\label{cor:pure-action}
Let $q=\delta_a$ and let $\ell_a(o):=\sum_rK(o,r\mid a)$.  Then $(\delta_a,K)$ is indifferent to the one-action,
uninformative-record channel that pays according to $\ell_a$.  Thus no separate pure-action axiom is needed.
\end{corollary}

\begin{proof}
Axiom~\textup{(A5)} first restricts the problem to the singleton support $\{a\}$ and embeds it back.  On a singleton,
Axiom~\textup{(A4)} may delete the record conditional on $o$; in the reverse direction it redraws $r$ from the
conditional distribution induced by $K(\cdot\mid a)$.  The two record processings and the two action maps give
indifference by weak order and block coherence.
\end{proof}

\begin{lemma}[Derived public-coin independence]\label{lem:public-coin}
Assume Axioms \textup{(A1), (A3), (A4), (A6)}.  For a full-support $q$, channels $K,L,M$ on $A$, and
$t\in(0,1)$,
\[
 (q,K)\succeq(q,L)
 \quad\Longleftrightarrow\quad
 (q,tK\oplus(1-t)M)\succeq(q,tL\oplus(1-t)M).
\]
\end{lemma}

\begin{proof}
Use as first stage an action-independent binary record with probabilities $(t,1-t)$.  Both branch posteriors equal
$q$.  Put $K$ versus $L$ on the first branch and the common continuation $M$ on the second.  Axiom~\textup{(A6)}
gives the forward implication.  The compounds have the terminal laws in \eqref{eq:public-mixture}, so
Lemma~\ref{lem:terminal-suff} identifies them with the stated public mixtures.

For the converse, if $K\not\succeq_qL$, completeness gives $L\succ_qK$.  Apply the strict clause of
Axiom~\textup{(A6)} with the continuations exchanged.  It makes the reverse public mixture strictly preferred, which
contradicts the displayed weak comparison.
\end{proof}

This is the old public-coin axiom derived from the continuation axiom.  It is not a ninth primitive.

\subsection{Material and product calibration}

\begin{lemma}[Common material scale and dummy neutrality]\label{lem:payoff-scale}
Assume Axioms \textup{(A1)--(A7)}.  There is a nonconstant vNM index $u:O\to\mathbb R$ such that every channel whose
payoff lottery $\ell$ is independent of the action and whose record is uninformative is ranked by
$\sum_o\ell(o)u(o)$, independently of its prior and action alphabet.

Moreover, independent dummy actions are neutral.  If $p\in\D(B)$ and
$K^\uparrow(o,r\mid a,b):=K(o,r\mid a)$, then
\begin{equation}
 (q\otimes p,K^\uparrow)\sim(q,K).
 \label{eq:dummy-neutrality}
\end{equation}
\end{lemma}

\begin{proof}
At a singleton action, Lemma~\ref{lem:public-coin} gives mixture independence over payoff lotteries.  The public coin
is visible, but Lemma~\ref{lem:terminal-suff} flattens it to the ordinary mixed payoff lottery because every terminal
posterior is the singleton prior.  Axiom~\textup{(A2)} supplies the Archimedean condition, so the finite mixture theorem gives expected utility.  Axiom
\textup{(A7)} makes the index nonconstant.

For arbitrary priors $q$ and $p$, let an action report ignore its input and draw the reported action from $p$.  Applied
to an action-independent payoff channel, \eqref{eq:action-processing} produces the same payoff lottery at prior $p$.
Axiom~\textup{(A5)} gives one weak comparison; repeating the construction from $p$ to $q$ gives the reverse.  Thus the
payoff-lottery order is independent of the prior.  Block coherence and two-way relabelling transport it across finite
action alphabets.  Inserting two payoff lotteries into the same reached branch and using the weak and strict parts of
Axiom~\textup{(A6)} shows that the insertion is an order isomorphism.  The two distinct payoff anchors supplied by
nonconstant $u$ will later fix its cardinal multiplier.

For dummy neutrality, project $(a,b)$ to $a$.  Equation \eqref{eq:action-processing} and Axiom~\textup{(A5)} give
$(q\otimes p,K^\uparrow)\succeq(q,K)$.  In the other direction, report $(a,b)$ from $a$ by independently drawing
$b\sim p$; the Bayesian pushforward is $K^\uparrow$, so Axiom~\textup{(A5)} gives the reverse comparison.  Weak order
and block coherence yield \eqref{eq:dummy-neutrality}.
\end{proof}

\begin{lemma}[Affine terminal representation and block bridge]\label{lem:affine-bridge}
For every full-support prior $q$, preferences on $\mathcal E_q$ have a continuous affine representative $W_q$.
These representatives may be normalized so that
\begin{equation}
 W_q\!\left(\sum_o\ell(o)\delta_{(o,q)}\right)=\sum_o\ell(o)u(o).
 \label{eq:payoff-normalization}
\end{equation}
With this normalization, dummy lifting preserves numerical value and the family $\{W_q\}$ represents every
block-supported comparison, including comparisons whose priors or action alphabets differ.
\end{lemma}

\begin{proof}
Terminal-law sufficiency turns the order at $q$ into an order on the mixture set $\mathcal E_q$.  Public-coin
independence is mixture independence there.  Along any mixture segment the union of the two finite supports is fixed,
so Axiom~\textup{(A2)} gives mixture continuity.  The mixture-space theorem therefore supplies a continuous affine
representative directly on all of $\mathcal E_q$.  Normalize it by the two distinct payoff values in
Lemma~\ref{lem:payoff-scale}, which gives
\eqref{eq:payoff-normalization}.

Dummy lifting is an affine ordinal isomorphism between a $q$-face and its copy in the $q\otimes p$-fibre.  Affine
uniqueness initially permits a transformation $x\mapsto cx+d$.  The two payoff anchors have the same numerical values
on both faces by \eqref{eq:payoff-normalization}; hence $c=1$ and $d=0$.

Finally, lift an alternative at $q\in\D(A)$ by an independent $p\in\D(B)$ and lift an alternative at $p$ by an
independent $q$.  Both lifted alternatives now have prior $q\otimes p$ on $A\times B$.  Dummy neutrality, block
coherence, and transitivity in a common multiple-block environment allow each original alternative to be replaced by
its lift.  Their block comparison is therefore represented by the single affine index $W_{q\otimes p}$, whose two
restrictions equal $W_q$ and $W_p$.  Boundary priors follow by support restriction and a nondegenerate dummy lift.
\end{proof}

\subsection{The pure-trace engine}

\begin{lemma}[Pure-record reduction and affine scale]\label{lem:trace-engine}
Fix $o_0\in O$ and restrict attention to channels
\[
 K_P(o,r\mid a):=\mathbf 1_{\{o=o_0\}}P(r\mid a).
\]
Under Axioms \textup{(A1)--(A6), (A8)}, every within-channel and block-supported comparison on this restriction is
ordered by $\I_{qP}(A;R)$.  After the common material scale and dummy normalisation of
Lemmas~\ref{lem:payoff-scale}--\ref{lem:affine-bridge}, the affine value of this restriction is
\begin{equation}
 W(q,K_P)=u(o_0)+\lambda\I_{qP}(A;R)
 \label{eq:pure-trace-scale}
\end{equation}
for one $\lambda>0$, common to all priors and finite action sets.
\end{lemma}

\begin{proof}
On a constant-payoff restriction, Axioms \textup{(A1)--(A6), (A8)} induce Conditions
\textup{(P1)--(P5)} of the pure-record model developed in Appendix~\ref{pt:sec:environment}: weak order and positive
orientation, continuity, block coherence, record and action data processing, and branchwise continuation.  The
self-contained pure-record characterization, Proposition~\ref{pt:prop:main}, therefore gives the ordinal conclusion
$\I(A;R)$.

By Lemma~\ref{lem:affine-bridge}, $W-u(o_0)$ is a common-scale affine representation of the ordinal $I$ order on the
constant-payoff domain, so it is
$\varphi(I)$ for a strictly increasing $\varphi$ with $\varphi(0)=0$.  Publicly mixing two experiments at the same
prior makes both their posterior laws and their $I$ values mix linearly; affinity gives
\[
 \varphi(t x+(1-t)y)=t\varphi(x)+(1-t)\varphi(y).
\]
Erasure channels at the uniform prior on $n$ actions attain the whole interval $[0,\log n]$.  These intervals are
nested and unbounded as $n$ grows, so continuity gives $\varphi(x)=\lambda x$ on $\mathbb R_+$.  Axiom
\textup{(A8)} gives $\lambda>0$, proving \eqref{eq:pure-trace-scale}.
\end{proof}

\subsection{Calibrated branch recursion}

\begin{lemma}[Calibrated branch recursion]\label{lem:branch-recursion}
The common-scale representative $W$ of Lemma~\ref{lem:affine-bridge} satisfies
\begin{equation}
 W(q,P\triangleright\{K^y\})
 =\lambda\I_q(A;Y)+\sum_{y:m(y)>0}m(y)W(q_y,K^y).
 \label{eq:branch-recursion}
\end{equation}
For a continuation that pays $o$ and reveals nothing further, $W=u(o)$.
\end{lemma}

\begin{proof}
Hold all continuations except one reached branch $y$ fixed.  Axiom~\textup{(A6)} and its strict converse say that
branch insertion is an order isomorphism.  Between affine faces, the change in compound value must therefore equal a
positive multiplier $c_y$ times the change in continuation value.

This multiplier is independent of the fixed continuations in the other branches.  Indeed, if their terminal sublaw is
$\eta_0$ and the reached branch has mass $m=m(y)$, affinity gives
\[
 W_q(\eta_0+m\eta_K)-W_q(\eta_0+m\eta_L)
 =m\,L_q(\eta_K-\eta_L),
\]
where $L_q$ is the linear part of $W_q$; the background $\eta_0$ cancels.

By the background-independence just shown, reset every other continuation to the deterministic payoff $o_0$ fixed in
Lemma~\ref{lem:trace-engine} and an uninformative record.  Now choose in branch $y$ two continuations with that same
payoff and different, nontrivial
record experiments.  Both the local and aggregate comparisons lie in the pure-record restriction.  Equation
\eqref{eq:pure-trace-scale} and the
mutual-information chain rule give
\[
 \Delta W_{\mathrm{compound}}
 =\lambda m(y)\Delta I_{\mathrm{branch}}
 =m(y)\Delta W_{\mathrm{branch}}.
\]
Thus $c_y=m(y)$.  The multiplier belongs to the insertion map, so it applies to every pair of continuations, including
payoff changes and mixed payoff--record changes.  If $q_y$ is degenerate and has no nontrivial trace direction, append
an independent nondegenerate dummy action, use the preceding calibration there, and project it away using
\eqref{eq:dummy-neutrality}; the conclusion is unchanged.

Changing branches one at a time now shows that the two sides of \eqref{eq:branch-recursion} differ by a term depending
only on the fixed first stage.  Evaluate that term with every continuation paying $o_0$ and revealing nothing.
The compound then lies in the pure-record restriction and has value $u(o_0)+\lambda I_q(A;Y)$, while the weighted
continuation values sum to $u(o_0)$.  The missing term is therefore exactly $\lambda I_q(A;Y)$.
\end{proof}

\subsection{Completion of the representation}

\begin{proof}[Proof of Theorem~\ref{thm:main}]
Assume Axioms \textup{(A1)--(A8)} and fix an arbitrary joint channel $K$.  Write $Z:=O\times R$ and sequentialise
$K$ as follows.  The first-stage record channel is
\[
 P_K(z\mid a):=K(o,r\mid a),\qquad z=(o,r).
\]
After the record $z=(o,r)$, the continuation pays $o$ for sure and reveals nothing further.  This does not mean that
payoff is paid early: the first record announces the payoff that the terminal continuation subsequently delivers.

The compound channel is indifferent to $K$.  Starting from $K$, a record processor can copy the already observed
pair $(o,r)$ into the compound record; starting from the compound, another processor deletes the copy.  Both
processors leave the payoff coordinate fixed, so Axiom~\textup{(A4)} applies in both directions.  The calibrated
branch recursion now gives
\[
 W(q,K)
 =\lambda\I_q(A;Z)+\sum_zp_{qK}(z)u(o_z)
 =\lambda\I_{qK}(A;O,R)+\E_{q,K}[u(O)].
\]
This is \eqref{eq:representation}, including block comparisons by Lemma~\ref{lem:affine-bridge}.  Axiom
\textup{(A7)} makes $u$ nonconstant and Axiom~\textup{(A8)} makes $\lambda>0$.

Conversely, let preferences be represented by \eqref{eq:representation}.  Weak order and continuity are immediate.
Block tags are constant under a block-supported lottery, so they change neither expected utility nor mutual
information, proving block coherence.  Record processing preserves the payoff marginal and weakly reduces mutual
information by data processing.  Equation \eqref{eq:action-processing} also preserves the payoff marginal, and the
Markov chain $(O,R)-A-B$ gives $\I(A;O,R)\ge\I(B;O,R)$.  Thus Axioms \textup{(A4)} and \textup{(A5)} hold.

For a compound channel, the two chain rules are
\begin{align*}
 \E[u(O)]&=\sum_y m(y)\E_{q_y,K^y}[u(O)],\\
 \I(A;Y,O,R)&=\I(A;Y)+\sum_y m(y)\I_{q_y,K^y}(A;O,R).
\end{align*}
The first-stage term $\lambda\I(A;Y)$ is common to the two continuation profiles.  Their aggregate value difference is
therefore the positive-probability weighted sum of their branch value differences.  This proves both clauses of
Axiom~\textup{(A6)}.  Nonconstant $u$ gives Axiom~\textup{(A7)}, and $\lambda>0$ gives Axiom~\textup{(A8)}.
\end{proof}

\begin{proof}[Proof of Proposition~\ref{prop:uniqueness}]
Choose $o_*\in O$ with $u(o_*)=\underline u$.  First note that the range of $V$ is
$[\underline u,\infty)$.  The lower bound is immediate.  For $q\in\D(A)$ and $t\in[0,1]$, let a constant-$o_*$
record channel reveal the realised action with
probability $t$ and otherwise report $*$:
\[
 E_t(o_*,a'\mid a)=t\1_{a'=a},
 \qquad
 E_t(o_*,*\mid a)=1-t.
\]
Then $V(q,E_t)=\underline u+\lambda t\Sh(q)$.  At the uniform lottery on $n$ actions these values fill
$[\underline u,\underline u+\lambda\log n]$; the intervals are unbounded.

For part \textup{(i)}, let $W$ be a common-scale representation and define
$\varphi(v):=W(q,K)$ for any pair with $V(q,K)=v$.  This is well defined: equality of the two $V$-values gives
indifference in their two-block environment by Theorem~\ref{thm:main}, and the common-scale property gives equality
of the two $W$-values.  A strict inequality between $V$-values similarly gives a strict inequality between
$W$-values.  Thus $\varphi$ is strictly increasing and $W=\varphi\circ V$.  The converse is immediate from
Theorem~\ref{thm:main}.

For part \textup{(ii)}, the expected-utility and mutual-information chain rules give, for any two continuation
profiles,
\[
 V\bigl(q,P\triangleright\{K^y\}\bigr)
 -V\bigl(q,P\triangleright\{L^y\}\bigr)
 =\sum_{y:m(y)>0}m(y)
 \bigl[V(q_y,K^y)-V(q_y,L^y)\bigr].
\]
Hence every $\alpha V+\beta$, with $\alpha>0$, is common-scale and branch-additive.

Conversely, let the branch-additive common-scale representation be $W=\varphi\circ V$ and define
\[
 \psi(x):=\varphi(\underline u+x)-\varphi(\underline u),
 \qquad x\ge0.
\]
Fix $n\ge2$ and let $q$ be uniform on $n$ actions.  For $x\in[0,\lambda\log n]$, choose a constant-$o_*$ erasure
channel whose value is $\underline u+x$.  In the first continuation profile, use this channel after the first branch
of a public coin with probabilities $\theta$ and $1-\theta$, and use the uninformative constant-$o_*$ channel after
the second.  Let the comparison profile use the uninformative channel after both branches.  The two compound values
are $\underline u+\theta x$ and $\underline u$, so \eqref{eq:branch-additive-index} gives
\[
 \psi(\theta x)=\theta\psi(x)
 \qquad(0\le\theta\le1).
\]
Taking $x=\lambda\log n$ shows that $\psi$ is linear on $[0,\lambda\log n]$.  These intervals are nested and
unbounded, so $\psi(x)=\alpha x$ on $[0,\infty)$ for one $\alpha>0$.  Therefore
$\varphi(v)=\alpha v+\beta$ on the range of $V$.

Finally, any second trace-tempered index $\widetilde V$ is itself branch-additive.  Applying part \textup{(ii)} and
then restricting to deterministic payoff outcomes and to constant-payoff record channels gives
$\widetilde u=\alpha u+\beta$ and $\widetilde\lambda=\alpha\lambda$.
\end{proof}

\begin{proof}[Proof of Corollary~\ref{cor:matching}]
The result follows from Theorem~\ref{thm:main} and
$\I_{qK}(A;O,R)=\Sh(q)-\int\Sh(\pi)d\mu_{qK}(\pi)$.
\end{proof}

\subsection{Independence of the axioms}

\begin{proof}[Proof of Proposition~\ref{prop:independence}]
Fix a nonconstant $u$ and write $U(q,K):=\E_{q,K}[u(O)]$, $Z:=(O,R)$, and $I:=\I(A;Z)$.  All scores below are
extended to block and compound channels by the same formula.

\emph{Failure of Axiom \textup{(A1)}.}
Use the Pareto relation on the two coordinates $(U,I)$:
\[
 (q,K)\succeq(p,L)
 \quad\Longleftrightarrow\quad
 U(q,K)\ge U(p,L)\ \text{and}\ I(q,K)\ge I(p,L).
\]
It is transitive but incomplete whenever one alternative has more payoff and the other more information.  Its graph is
closed.  Both data-processing axioms are coordinatewise monotonic, and both chain rules make branchwise weak and strict
improvements aggregate.  The two relevance axioms and block coherence hold.

\emph{Failure of Axiom \textup{(A2)}.}
Order alternatives lexicographically by $(U,I)$, with payoff first.  This is a weak order and satisfies all structural,
processing, continuation, and relevance requirements.  It is not continuous even within one fixed channel.  Let a
three-action channel reveal the action fully and let only action $2$ pay the high outcome.  Put
\[
 p=(1/4,1/2,1/4),\qquad q=(0.49,0.50,0.01),\qquad
 q_n=(0.49-\varepsilon_n,0.50+\varepsilon_n,0.01),
\]
where $\varepsilon_n\downarrow0$.  The lotteries $p$ and $q$ have the same expected utility but
$\Sh(q)<\Sh(p)$, so $p\succ_Kq$.  Every $q_n$ has strictly higher expected utility than $p$, so
$q_n\succ_Kp$.  Thus the graph in Axiom~\textup{(A2)} is not closed.

\emph{Failure of Axiom \textup{(A3)}.}
Treat block labels as syntactic tags, with the outermost tag taking precedence.  Put $w(a^0)=1$, $w(a^i)=0$ for
$i\ne0$, and $w(a)=0$ on untagged actions.  For $\varepsilon>0$, use
\[
 V_\varepsilon(q,K):=U(q,K)+I(q,K)+\varepsilon\sum_aq(a)w(a).
\]
Every record- or action-processing comparison places the original alternative in block $0$, so the data-processing
inequality receives the additional $\varepsilon$.  If $\Delta_y$ is the base-score difference in continuation branch
$y$, its block comparison is weak exactly when $\Delta_y+\varepsilon\ge0$, while the aggregate block difference is
\[
 \sum_y m(y)\Delta_y+\varepsilon
 =\sum_y m(y)(\Delta_y+\varepsilon).
\]
Thus Axiom~\textup{(A6)}, including strictness, holds; the remaining axioms are immediate.  Duplication coherence fails.
Take constant payoff, $K$ revealing a binary action, a degenerate $q$, and a nearby nondegenerate $p$ with
$0<\Sh(p)<\varepsilon$.  Then $p\succ_Kq$, but $q^0\succ_{K\sqcup K}p^1$.

\emph{Failure of Axiom \textup{(A4)}.}
For $\varepsilon>0$, use
\begin{equation}
 V(q,K):=U(q,K)+I(A;Z)-\varepsilon\Sh(Z\mid A).
 \label{eq:no-record-dp}
\end{equation}
This score is continuous and branch-additive, since both mutual information and conditional entropy obey the chain
rule.  Under action processing $A\to B$, the score difference between the original and processed alternatives is
\[
 \bigl[I(A;Z)-I(B;Z)\bigr]
 +\varepsilon\bigl[\Sh(Z\mid B)-\Sh(Z\mid A)\bigr]
 =(1+\varepsilon)I(A;Z\mid B)\ge0,
\]
so \eqref{eq:no-record-dp} satisfies Axiom~\textup{(A5)}.  It also satisfies both relevance axioms.  Record data
processing fails: with constant payoff and a fair record independent of $A$, the original value is $-\varepsilon$;
deleting that noise gives value zero.

\emph{Failure of Axiom \textup{(A5)}.}
Let $G(q):=1-\sum_aq(a)^2$ be Gini entropy and define its posterior reduction
\[
 J_G(q,K):=G(q)-\int G(\pi)\,d\mu_{qK}(\pi),
 \qquad V:=U+J_G.
\]
Concavity of $G$ gives record data processing, and the potential-difference form gives the exact branch chain rule.
Full revelation is strictly valuable.  Action processing can nevertheless increase $J_G$.  Take
$q=(1/8,1/8,3/4)$, a binary record with
\[
 \Pr(0\mid a_1)=\Pr(0\mid a_2)=0,
 \qquad \Pr(0\mid a_3)=1/4,
\]
and merge $a_1,a_2$ into one reported action.  Direct calculation gives $J_G=9/416$ before the merge and
$J_G=3/104$ afterwards.  Hence Axiom~\textup{(A5)} fails.

\emph{Failure of Axiom \textup{(A6)}.}
Use $V:=U+\min\{I,1\}$, with information measured in bits.  It is continuous, respects both data-processing axioms,
and has both strict orientations.  Let $q$ be uniform on four actions.  A first-stage record reveals which of two pairs
contains the action.  In every branch, revealing the action within the pair strictly improves the continuation from
zero to one bit.  The two compounds have respectively one and two bits before capping and are therefore indifferent.
The strict clause of Axiom~\textup{(A6)} fails.

\emph{Failure of Axiom \textup{(A7)}.}
Use $V:=I$.  All other axioms hold, but all payoff lotteries are indifferent.

\emph{Failure of Axiom \textup{(A8)}.}
Use $V:=U$.  All other axioms hold, but revealing the action at a fixed payoff is valueless.
\end{proof}

\section{Proofs of the choice implications}\label{app:choice-proofs}

Results stated in the text only in summary form (Propositions~\ref{prop:binary}
and~\ref{prop:trace-demand}, Corollary~\ref{cor:lambda-limits}, and Remark~\ref{rem:elicitation}) are restated here
in full.  Proofs appear in an order that front-loads shared machinery.

\begin{proof}[Proof of Proposition~\ref{prop:signature}]
(i) Under either lottery, each visible pair $(o_0,r_j)$, $j=1,2$, has probability $\tfrac12$: under $q'$ because
each $a_j$ produces $(o_0,r_j)$ for sure, under $q''$ because each of $a_3,a_4$ produces either pair with equal
probability.  Hence $p_{q'K}=p_{q''K}$.  On $\supp(q')$ the rows are degenerate and distinct, so the visible pair
reveals the action and $\I_{q'K}=\Sh(q')=\log2$; on $\supp(q'')$ the rows are identical, so the visible pair is
independent of the action and $\I_{q''K}=0$.  The payoff terms agree because the payoff marginals do, giving the
displayed gap.
(ii) All rows of $K^\circ$ are equal, so $\I_{q,K^\circ}=0$ for every $q$; since the payoff is constantly $o_0$,
the representation gives indifference.
\end{proof}

\begin{proof}[Proof of Proposition~\ref{prop:optimal-choice}]
Expected utility is linear in $q$, while mutual information is concave in the input distribution of a fixed channel
\cite{CoverThomas2006}.  Continuity on the compact simplex gives existence.  Using
\[
 I_{q,K}(A;Z)=\sum_aq(a)\KL\!\left(K_Z(\cdot\mid a)\,\middle\|\,m_{q,K}\right),
\]
its directional derivative in the mass on action $a$ equals the displayed divergence up to a term common to every
action.  The result is the Kuhn--Tucker condition for a concave programme.  At constant $u$ it reduces to the classical
capacity condition.
\end{proof}

\begin{proof}[Proof of Proposition~\ref{prop:optimal-marginal}]
The entropy identity gives
\[
 V(q,K)
 =\lambda\Sh(m_{q,K})
  +\sum_aq(a)\bigl[\bar u(a)-\lambda\Sh(P_a)\bigr].
\]
If $q_0$ and $q_1$ are optimal with $m_{q_0,K}\ne m_{q_1,K}$, strict concavity of Shannon entropy and linearity of
$q\mapsto m_{q,K}$ and of the remaining sum give $V(tq_0+(1-t)q_1,K)>C^*(K)$ for $t\in(0,1)$, a contradiction.
Hence every optimum induces the same $m^*$.  By Proposition~\ref{prop:optimal-choice}, an optimum $q^*$ satisfies
the displayed equalities at some constant $c$; averaging them under $q^*$ gives $c=V(q^*,K)=C^*(K)$, so
$\supp(q^*)\subseteq\mathcal A^*$.  Conversely, if $m_{q,K}=m^*$ and $\supp(q)\subseteq\mathcal A^*$, then
\[
 V(q,K)
 =\sum_aq(a)\bigl[\bar u(a)+\lambda\KL(P_a\|m^*)\bigr]
 =C^*(K).
\]
Affine independence makes the representation $m^*=\sum_aq(a)P_a$ with $\supp(q)\subseteq\mathcal A^*$ unique.
Identical rows share $\bar u(a)$ and $\Sh(P_a)$, so both $m_{q,K}$ and the displayed expression for $V$ are
unchanged by redistribution within a class; collapsing the classes and lifting aggregate masses proves the last
claim.
\end{proof}

\begin{proof}[Proof of Proposition~\ref{cor:support-pure}]
Throughout, $m^*$ denotes the common optimal marginal of Proposition~\ref{prop:optimal-marginal}.
(i) Suppose $m^*(z)=0$ for some $z\in Z_K$ and choose $a$ with $K_Z(z\mid a)>0$.  Then
$\KL(K_Z(\cdot\mid a)\|m^*)=+\infty$, which is incompatible with the finite on-support equality if some optimum
uses $a$ and with the off-support inequality of Proposition~\ref{prop:optimal-choice} if none does.  Under the
private-consequence condition, an optimum with $q^*(a)=0$ would force $m^*(z_a)=0$.

(ii) For $q=\delta_a$ the induced marginal is $K_Z(\cdot\mid a)$ and the on-support constant is $\bar u(a)$; the
displayed system is the off-support condition of Proposition~\ref{prop:optimal-choice}, necessary and sufficient by
concavity.  Under strict inequalities, $\mathcal A^*=\{a\}$ in Proposition~\ref{prop:optimal-marginal}, so every
optimum is $\delta_a$.
\end{proof}

\begin{proposition}[Optimal mixing]\label{prop:binary}
Let $A=\{a,b\}$ with $K_a:=K_Z(\cdot\mid a)\ne K_b:=K_Z(\cdot\mid b)$ and
$\Delta:=\bar u(a)-\bar u(b)\ge0$.  For $x\in[0,1]$ let $q_x$ put mass $x$ on $b$, so that
$m_{q_x,K}=(1-x)K_a+xK_b=:m_x$, and for $x\in(0,1)$ define
\[
 g(x):=\KL\!\left(K_b\,\middle\|\,m_x\right)-\KL\!\left(K_a\,\middle\|\,m_x\right).
\]
\begin{enumerate}[label=\textup{(\roman*)},leftmargin=*]
\item $g$ is continuous and strictly decreasing, with $g(0^+)=\KL(K_b\,\|\,K_a)$ and
$g(1^-)=-\KL(K_a\,\|\,K_b)$ as extended-real limits.  If $\Delta<\lambda\,\KL(K_b\,\|\,K_a)$, the optimum is
unique and interior, $q^*=q_{x^*}$ with $x^*$ the unique solution of $\Delta=\lambda g(x)$.
\item The optimal mass $x^*(\Delta,\lambda)$ on the inferior action is nonincreasing in $\Delta$ and
nondecreasing in $\lambda$, strictly so on the mixing region; at $\Delta=0$ it is the capacity-achieving input
of the two-row channel, independently of $\lambda$.
\end{enumerate}
\end{proposition}

\begin{proof}
Write $f(x):=V(q_x,K)=\bar u(a)-x\Delta+\lambda I(x)$ with
$I(x)=\Sh(m_x)-(1-x)\Sh(K_a)-x\Sh(K_b)$.  Since $K_a\ne K_b$,
$x\mapsto m_x$ is affine and injective, so strict concavity of $\Sh$ \cite{CoverThomas2006} makes $f$ strictly
concave on $[0,1]$, and a unique maximiser exists.  On $(0,1)$ every $z\in\supp K_a\cup\supp K_b$ has $m_x(z)>0$
and, using $\sum_z(K_b(z)-K_a(z))=0$,
\[
 I'(x)=\Sh(K_a)-\Sh(K_b)-\sum_z\bigl(K_b(z)-K_a(z)\bigr)\log m_x(z)=g(x),
\]
so $f'(x)=-\Delta+\lambda g(x)$, with $g$ continuous and strictly decreasing as the derivative of a strictly
concave function.  As $x\downarrow0$, $m_x(z)\to K_a(z)$; if $\supp K_b\subseteq\supp K_a$ every summand converges
and $g(0^+)=\KL(K_b\|K_a)$, while if some $\bar z$ has $K_b(\bar z)>0=K_a(\bar z)$ the summand
$-K_b(\bar z)\log\bigl(xK_b(\bar z)\bigr)$ diverges to $+\infty$, matching the extended-sense divergence.
Exchanging the roles of $a$ and $b$ gives the limit at $1$.  By concavity and the mean value theorem,
$\delta_a$ is optimal if and only if $\lim_{x\downarrow0}f'(x)\le0$, i.e.\ $\Delta\ge\lambda\,\KL(K_b\|K_a)$;
this is the threshold stated in the text, and Proposition~\ref{cor:support-pure}\textup{(ii)} with two actions.
If $\Delta<\lambda\,\KL(K_b\|K_a)$, the maximiser is not $x=0$ by the threshold and not $x=1$
because $f'(1^-)=-\Delta-\lambda\KL(K_a\|K_b)<0$; it is therefore the unique root of $\Delta=\lambda g(x)$, which
exists because $g$ decreases strictly between its limits.  This proves \textup{(i)}.  The comparative statics of
\textup{(ii)} follow by inverting the
strictly decreasing $g$ at $\Delta/\lambda$: raising $\Delta$ raises the target, raising $\lambda$ lowers it; at
$\Delta=0$ the condition $g(x)=0$ is the capacity condition of Proposition~\ref{prop:optimal-choice} at constant
payoff.
\end{proof}

\begin{proposition}[Trace demand]\label{prop:trace-demand}
Fix $K$ and write $e(q):=\E_{q,K}[u(O)]$, $i(q):=\I_{qK}(A;O,R)$,
\[
 V^*(\lambda):=\max_{q\in\D(A)}\{e(q)+\lambda i(q)\},
 \qquad
 Q(\lambda):=\arg\max_{q\in\D(A)}\{e(q)+\lambda i(q)\}.
\]
\begin{enumerate}[label=\textup{(\roman*)},leftmargin=*]
\item $V^*$ is convex on $(0,\infty)$, $i$ is affine on the convex compact set $Q(\lambda)$, and
\[
 (V^*)'_-(\lambda)=\min_{q\in Q(\lambda)}i(q),
 \qquad
 (V^*)'_+(\lambda)=\max_{q\in Q(\lambda)}i(q);
\]
at every point of differentiability, $(V^*)'(\lambda)=i(q)$ for every $q\in Q(\lambda)$.
\item Let $0<\lambda_1<\lambda_2$ and $q_j\in Q(\lambda_j)$, writing $e_j:=e(q_j)$, $i_j:=i(q_j)$.  Then
$i_1\le i_2$, $e_1\ge e_2$, and
\[
 \lambda_1(i_2-i_1)\le e_1-e_2\le\lambda_2(i_2-i_1).
\]
Hence either $(e_1,i_1)=(e_2,i_2)$, or $i_1<i_2$ and $e_1>e_2$ and
\[
 \lambda_1\le\frac{e_1-e_2}{i_2-i_1}\le\lambda_2.
\]
\end{enumerate}
\end{proposition}

\begin{proof}[Proof of Proposition~\ref{prop:trace-demand}]
(i) $V^*$ is a pointwise maximum of functions affine in $\lambda$, hence convex and, since
$0\le i\le\log|A|$, finite.  On the convex compact set $Q(\lambda)$ the objective is constant and $e$ is linear, so
$i=(V^*-e)/\lambda$ is affine there.  For $q\in Q(\lambda)$ and $\mu>0$,
$V^*(\mu)\ge V^*(\lambda)+(\mu-\lambda)i(q)$, so each such $i(q)$ is a subgradient.  For the right derivative,
choose $q_h\in Q(\lambda+h)$; then
\[
 \max_{q\in Q(\lambda)}i(q)
 \le\frac{V^*(\lambda+h)-V^*(\lambda)}{h}
 \le i(q_h),
\]
the first inequality by the subgradient bound, the second by optimality of $Q(\lambda)$ at $\lambda$.  By
compactness and continuity, accumulation points of $(q_h)_{h\downarrow0}$ lie in $Q(\lambda)$, so the right-hand
side accumulates in $i(Q(\lambda))$ and the right derivative equals the maximum.  The left derivative is symmetric,
and a convex function on an interval is differentiable off a countable set.

(ii) Optimality at the two weights gives
$e_1+\lambda_1i_1\ge e_2+\lambda_1i_2$ and $e_2+\lambda_2i_2\ge e_1+\lambda_2i_1$.  Adding yields
$(\lambda_2-\lambda_1)(i_2-i_1)\ge0$, hence $i_2\ge i_1$; the two inequalities rearrange to the displayed bracket,
which gives $e_1\ge e_2$ and, when the pairs differ, the final bound.
\end{proof}

\begin{corollary}[Extreme trace weights]\label{cor:lambda-limits}
Let $e^0:=\max_qe(q)$, $i^0:=\max\{i(q):e(q)=e^0\}$, $C:=\max_qi(q)$, and $e^C:=\max\{e(q):i(q)=C\}$.  If
$\lambda_n\downarrow0$ and $q_n\in Q(\lambda_n)$, every accumulation point of $(q_n)$ maximises $i$ among the
maximisers of $e$; if $\lambda_n\uparrow\infty$, every accumulation point maximises $e$ among the maximisers of
$i$.  Moreover
\[
 V^*(\lambda)=e^0+\lambda i^0+o(\lambda)
 \quad\text{as }\lambda\downarrow0,
 \qquad
 V^*(\lambda)=\lambda C+e^C+o(1)
 \quad\text{as }\lambda\uparrow\infty.
\]
\end{corollary}

\begin{proof}[Proof of Corollary~\ref{cor:lambda-limits}]
For $\lambda_n\downarrow0$, comparing $q_n$ with an $e$-maximiser attaining $i^0$ gives
\[
 0\le e^0-e(q_n)\le\lambda_n\bigl(i(q_n)-i^0\bigr).
\]
Boundedness of $i$ forces $e(q_n)\to e^0$, and the same inequality gives $i(q_n)\ge i^0$; continuity delivers the
accumulation claim.  Dividing by $\lambda_n$,
\[
 \frac{V^*(\lambda_n)-e^0}{\lambda_n}=i(q_n)-\frac{e^0-e(q_n)}{\lambda_n},
\]
where the last fraction lies in $[0,\,i(q_n)-i^0]$ and vanishes, proving the first expansion.  For
$\lambda_n\uparrow\infty$, comparing $q_n$ with an $i$-maximiser attaining $e^C$ gives
$0\le\lambda_n(C-i(q_n))\le e(q_n)-e^C$, and the symmetric argument completes the proof.
\end{proof}

\begin{proof}[Proof of Proposition~\ref{prop:garbling-price}]
The processor leaves the payoff coordinate unchanged, so the expected-utility terms coincide and the difference is
$\lambda[\I_{qK}(A;O,R)-\I_{qL}(A;O,R')]$.  Write $Z:=(O,R)$ and $Z':=(O,R')$.  By construction $A\to Z\to Z'$ is
a Markov chain, so expanding $\I(A;Z,Z')$ in both orders gives
\[
 \I(A;Z)-\I(A;Z')=\I(A;Z\mid Z')=\I(A;R\mid O,R'),
\]
the last equality because $O$ is part of $Z'$.  Conditional mutual information is nonnegative and vanishes exactly
under the conditional independence in \textup{(b)} \cite{CoverThomas2006}, proving the identity and
\textup{(a)}$\Leftrightarrow$\textup{(b)}.  For \textup{(b)}$\Rightarrow$\textup{(c)}: the Markov property gives
$\Pr(A=\cdot\mid Z,Z')=\pi^{qK}_{Z}$ almost surely, while \textup{(b)} gives
$\Pr(A=\cdot\mid Z,Z')=\pi^{qL}_{Z'}$; the two random posteriors coincide almost surely, so their laws are equal.
For \textup{(c)}$\Rightarrow$\textup{(a)}: by the identity displayed after Corollary~\ref{cor:matching}, the
information term depends on $(q,K)$ only through $\mu_{qK}$, and the payoff terms already agree.
\end{proof}

\begin{proof}[Proof of Proposition~\ref{prop:blackwell}]
(i) If $L=KT$, Proposition~\ref{prop:garbling-price} gives $\I_{qK}\ge\I_{qL}$ at every $q$.  Because the payoff
kernels coincide, $V(q,K)-V(q,L)=\lambda[\I_{qK}-\I_{qL}]$ with $\lambda>0$, so uniform value dominance and
$K\succeq_cL$ are the same statement.

(ii) Let $A=\{0,1\}$, fix $o^\ast\in O$, and let both channels pay $o^\ast$ surely.  Let $K$ carry the binary
erasure record with erasure probability $\varepsilon=\tfrac12$ and $L$ the binary symmetric record with crossover
$p=\tfrac15$:
\[
 K(o^\ast,a\mid a)=1-\varepsilon,\quad K(o^\ast,\mathrm e\mid a)=\varepsilon,
 \qquad
 L(o^\ast,a\mid a)=1-p,\quad L(o^\ast,1-a\mid a)=p.
\]
The payoff kernels agree row by row, so it suffices to compare informations, and the comparison is invariant to the
base; compute in bits with $h$ the binary entropy.  For $q=(1-t,t)$, grouping records gives
\[
 \I_{qK}=(1-\varepsilon)\,h(t),
 \qquad
 \I_{qL}=h(p\ast t)-h(p),
 \qquad p\ast t:=p(1-t)+(1-p)t.
\]
By Mrs Gerber's Lemma \cite{WynerZiv1973}, $v\mapsto h\bigl(p\ast h^{-1}(v)\bigr)$ is convex on $[0,1]$, with
$h^{-1}$ valued in $[0,\tfrac12]$; since $h(p\ast t)$ is invariant under $t\mapsto1-t$, take $t\le\tfrac12$ and
$v=h(t)$.  The chord between $v=0$ and $v=1$ gives
\[
 h(p\ast t)\le(1-v)\,h(p)+v,
 \qquad\text{hence}\qquad
 \I_{qL}\le\bigl(1-h(p)\bigr)h(t)\le(1-\varepsilon)\,h(t)=\I_{qK},
\]
the last step because $\varepsilon=\tfrac12\le h(\tfrac15)$.  Thus $K\succeq_cL$.  Now suppose $L=KT$ for a
processor $T$, and write $t_r:=T(0\mid o^\ast,r)$ for $r\in\{0,1,\mathrm e\}$.  Matching the probability of the
processed record $0$ in the two rows,
\[
 (1-\varepsilon)t_0+\varepsilon t_{\mathrm e}=1-p,
 \qquad
 (1-\varepsilon)t_1+\varepsilon t_{\mathrm e}=p,
\]
and subtracting gives $t_0-t_1=(1-2p)/(1-\varepsilon)=\tfrac65>1$, impossible.  Compositions of record processors
are record processors, so no chain of garblings produces $L$ from $K$ either.
\end{proof}

\begin{proof}[Proof of Proposition~\ref{prop:wtp}]
(i) The material part of the compound is $(1-\varepsilon)\E_{q,K}[u(O)]+\varepsilon w(\beta)$.  For the trace part,
the coin is action-independent, so both branch posteriors equal $q$ and the chain rule for the compound gives an
information term of $(1-\varepsilon)\I_{qK}(A;O,R)$, the sweetener branch contributing zero.
(ii) By the moreover clause of Theorem~\ref{thm:main}, the indifference holds if and only if
$(1-\varepsilon)V(q,K)+\varepsilon w(\beta)=(1-\varepsilon)V(p,L)+\varepsilon w(\beta')$; rearrange, and
substitute the representation when the material parts cancel.
(iii) $w(\beta)-w(\beta')$ is continuous on $[0,1]^2$ with range $[-\Delta u,\Delta u]$; apply the intermediate
value theorem to the equation in \textup{(ii)}.
\end{proof}

\begin{proof}[Proof of Proposition~\ref{prop:single-price}]
(i) Inclusion is immediate.  Conversely, fix $(t,w)$ in the strip, choose $n$ with $\log n\ge t$, and on
$A=\{1,\dots,n\}$ use continuity of entropy along a segment from a point mass to the uniform lottery to find
$q_t$ with $\Sh(q_t)=t$.  Let $K(o^+,a\mid a)=\beta$ and $K(o^-,a\mid a)=1-\beta$ with $\beta$ chosen so that
$w(\beta)=w$.  The payoff lottery is action-independent with mean $w$, and the record reveals the action, so
$\phi(q_t,K)=(t,w)$.  The description of the indifference classes then reads the representation through $\phi$,
using the moreover clause of Theorem~\ref{thm:main}.
(ii) The identity rearranges $E_{qK}+\lambda\I_{qK}=E_{pL}+\lambda\I_{pL}$.  For uniqueness, let
$\lambda'\ne\lambda$, put $t:=\tfrac12\min\{1,\,(\overline u-\underline u)/\lambda\}>0$, and use \textup{(i)} to
realise pairs with $\phi$-coordinates $(t,\overline u-\lambda t)$ and $(0,\overline u)$; they are
indifferent, while $E+\lambda'\I$ differs across them by $(\lambda-\lambda')t\ne0$.
\end{proof}

\begin{remark}[Eliciting $\lambda$]\label{rem:elicitation}
Normalise $u(o^+)=1$ and $u(o^-)=0$.  All comparisons below are block comparisons, hence within the primitive
domain, and no information quantity is estimated: the information gap is an accounting identity of the design.
\emph{Step 1.}  For each $o\in O$, elicit $\beta_o$ with $(\delta_\ast,K_o)\sim(\delta_\ast,B_{\beta_o})$, where
$K_o$ is the degenerate one-action channel of Axiom~\textup{(A7)}; one-action channels carry zero information, so
$u(o)=\beta_o$.
\emph{Step 2.}  Fix $o\in O$, $n\ge2$, the uniform $q$ on $\{1,\dots,n\}$, a dilution $\varepsilon\in(0,1)$, and
the constant-payoff channels $K^{\mathrm{id}}_o$ and $K^0_o$ of Axiom~\textup{(A8)}.  Elicit $\beta$ with
\[
 \bigl(q,\,D_{\varepsilon,\beta}K^{0}_o\bigr)\sim\bigl(q,\,D_{\varepsilon,0}K^{\mathrm{id}}_o\bigr);
 \qquad\text{then}\qquad
 \lambda=\frac{\varepsilon}{1-\varepsilon}\cdot\frac{\beta}{\log n},
\]
by Proposition~\ref{prop:wtp}\textup{(i)} applied to $V(q,K^{\mathrm{id}}_o)=u(o)+\lambda\log n$ and
$V(q,K^{0}_o)=u(o)$.  Such a $\beta$ exists if and only if $\lambda\log n\le\varepsilon/(1-\varepsilon)$; a
subject still strictly preferring the right block at $\beta=1$ reveals
$\lambda>\varepsilon/\bigl((1-\varepsilon)\log n\bigr)$, and $\varepsilon$ is raised.
\emph{Controls.}  Orientation: $(q,K^{\mathrm{id}}_o)\succ(q,K^{0}_o)$ is Axiom~\textup{(A8)} itself.  Garbling:
under the total record garbling of $K^{\mathrm{id}}_o$, both blocks carry zero information and indifference must
obtain; a surviving strict preference indicates that an uncontrolled feature of the display still traces the
action.  Transport: repeating Step 2 at another $n$, prior, or payoff level must return the same $\lambda$, by
Proposition~\ref{prop:single-price}\textup{(ii)}; an estimate that declines in $\log n$ is the signature of the
capped-information criterion.
\end{remark}

\clearpage

\part*{Self-contained foundation for the trace component}
\addcontentsline{toc}{part}{Self-contained foundation for the trace component}

This appendix contains the pure-record characterization on which
Lemma~\ref{lem:trace-engine} rests.  It is included in full so that the present
paper has no logical dependency on an earlier manuscript.  In this auxiliary
model, $X$ is a finite visible-record alphabet, $P:A\to\Delta(X)$ is a record
channel, and the primitive relation ranks action lotteries by the trace they
leave in $X$.  Conditions \textup{(P1)--(P5)} below are exactly the restrictions
induced on the constant-payoff subdomain by, respectively,
\textup{(A1),(A8)}, \textup{(A2)}, \textup{(A3)}, \textup{(A4),(A5)}, and
\textup{(A6)} of the main model.

%<*leanpurestatement>
\section{The pure-trace characterization}
\label{pt:sec:environment}\label{lf:app:pure-trace}
% ===============================================================

Let $A$ be a finite set of \emph{actions} and $X$ a finite set of \emph{records}.  A channel is a stochastic kernel
$P:A\to\Delta(X)$.  For each finite channel $P$, the primitive behavioural datum is a binary relation
\[
\succeq_P\quad\text{on}\quad \Delta(A).
\]
Condition \textup{(P1)} below requires this relation to be a weak order.  The interpretation is that
$q\succeq_P q'$ means: in the fixed environment $P$, lottery $q$ displays at least as much
traceability as $q'$.
The environment $P$ is not an object of choice.  A lottery is the procedure being evaluated; the channel is the fixed
observation technology through which the realised act may be recorded by records.

\paragraph{Comparison environments.}
To compare lotteries carried by different channels without abandoning channel-indexed preferences, place the channels in
labelled, disjoint blocks of one larger environment.  For channels $P:A\to\Delta(X)$ and $Q:B\to\Delta(Y)$, write
$P\sqcup Q$ for the block-diagonal channel with action alphabet
$(A\times\{0\})\cup(B\times\{1\})$ and record alphabet
$(X\times\{0\})\cup(Y\times\{1\})$.  It sends $(a,0)$ to $(x,0)$ with probability $P(x\mid a)$ and $(b,1)$ to $(y,1)$
with probability $Q(y\mid b)$; all cross-block probabilities are zero.  For $q\in\Delta(A)$ and
$p\in\Delta(B)$, define lotteries on the tagged action alphabet by
\[
\begin{aligned}
q^0(a,0)&:=q(a), & q^0(b,1)&:=0,\\
p^1(a,0)&:=0,    & p^1(b,1)&:=p(b),
\end{aligned}
\qquad (a\in A,\ b\in B).
\]
The superscript is a block label, not a power.  Accordingly, $q^0\succeq_{P\sqcup Q}p^1$ is an ordinary preference
comparison inside one fixed environment.\footnote{For a complete numerical example, let
$A=\{a_1,a_2\}$, $B=\{b_1,b_2\}$, $X=\{x_0,x_1\}$, and $Y=\{y_0,y_1\}$.  With rows indexed by actions and columns
by records, take
$P=\bigl(\begin{smallmatrix}0.1&0.9\\0.8&0.2\end{smallmatrix}\bigr)$ and
$Q=\bigl(\begin{smallmatrix}0.4&0.6\\0.6&0.4\end{smallmatrix}\bigr)$.
Order the tagged actions as $(a_1,0),(a_2,0),(b_1,1),(b_2,1)$ and the tagged records as
$(x_0,0),(x_1,0),(y_0,1),(y_1,1)$.  Then $P\sqcup Q=\operatorname{diag}(P,Q)$: the first two actions generate only
the first two records, and the last two actions only the last two.  If $q=(0.3,0.7)$ and $p=(0.8,0.2)$, their
embeddings are $q^0=(0.3,0.7\,;\,0,0)$ and $p^1=(0,0\,;\,0.8,0.2)$, where the semicolon separates the two blocks.  Thus
$q^0\succeq_{P\sqcup Q}p^1$ says precisely that lottery $q$, observed through $P$, is at least as traceable as lottery
$p$, observed through $Q$; it does not treat either channel as an object of choice.  Under the representation proved
below, the two sides yield approximately $0.330$ and $0.019$ bits of mutual information, respectively, and hence the
comparison is strict in this example.}
When the two blocks have the same action alphabet, $q^0$ and $q^1$ are the two tagged copies of the same lottery $q$. For a finite family
$\{P_i:A_i\to\Delta(X_i)\}_{i\in K}$, write $\bigsqcup_{i\in K}P_i$ for the analogous finite block environment and
$q_i^i$ for the lottery that assigns $q_i(a)$ to each tagged action $(a,i)$ and zero to every other block.

\paragraph{Pure-record conditions.}
The conditions are imposed directly on the family $\{\succeq_P\}_P$.
Five conditions are imposed on the family. The first three are structural:
      each environment carries a continuous weak order with a fixed orientation
      (P1, P2), and labelled blocks provide the comparison framework (P3).
      The substantive restrictions are two: garbling cannot increase
      traceability (P4), and sequential observation is evaluated branch by
      branch (P5).


\begin{description}[leftmargin=*,style=sameline]

\item[(P1) Weak order and local non-triviality.]
For every finite channel $P:A\to\Delta(X)$, $\succeq_P$ is complete and transitive on $\Delta(A)$.  The family is
locally non-trivial: for every finite action set $A$ with $|A|\ge 2$ and every full-support $q\in\Delta(A)$,
\[
q^0\succ_{\Id_A\sqcup U_A}q^1,
\]
where $\Id_A:A\to\Delta(A)$ is full revelation and $U_A:A\to\Delta(\{\ast\})$ is the one-record uninformative channel.

This condition gives each observation environment an ordinary weak order over act lotteries and fixes the orientation of the
comparison.  When the realised act is genuinely uncertain, observing it directly is strictly more traceable than observing
nothing about it.

\item[(P2) Continuity.]
For every pair of finite alphabets $(A,X)$, the set
\[
\{(P,q,q')\in \Delta(X)^A\times\Delta(A)\times\Delta(A):q\succeq_P q'\}
\]
is closed in the usual Euclidean product topology.  Thus within-environment rankings vary continuously with both the
lotteries and the fixed channel.

Continuity says that traceability comparisons are stable under small changes in the procedure and the environment.  Small
modelling perturbations should not create abrupt reversals.  The condition is stated entirely in the primitives: it quantifies
over channels and lotteries, not over derived Bayesian objects.  Continuity of block comparisons in posterior-law
convergence at a fixed full-support prior is not assumed; it is derived below
(Lemma~\ref{pt:lem:plcont}).

\item[(P3) Block-comparison coherence.]
\emph{Duplication:} for every channel $P:A\to\Delta(X)$ and every $q,q'\in\Delta(A)$,
\[
q\succeq_P q'
\quad\Longleftrightarrow\quad
q^0\succeq_{P\sqcup P}(q')^1.
\]
\emph{Irrelevant blocks:} for every finite block environment
$\bigsqcup_{k\in K}P_k$, every ordered pair of distinct blocks $i,j\in K$, and every
$q_i\in\Delta(A_i)$ and $q_j\in\Delta(A_j)$,
\[
q_i^i\succeq_{\bigsqcup_{k\in K}P_k}q_j^j
\quad\Longleftrightarrow\quad
q_i^0\succeq_{P_i\sqcup P_j}q_j^1.
\]
Applied at the ordered pair $(1,0)$ of a two-block environment, the second clause exchanges block labels; this is
how reversed block comparisons are brought to canonical form throughout
(Lemma~\ref{pt:lem:blockcoh}\textup{(a)}).

Labelled blocks are a way to place procedures carried by different channels in a common comparison problem.  The labels
should not themselves create traceability, alter the ranking inside a duplicated environment, or let unused blocks affect
the comparison between active blocks.

\item[(P4) Split data-processing aversion.]
Condition \textup{(P4)} is the conjunction of the following two clauses.  This
is the formulation used by the formal theorem: the record and action
coordinates are processed separately, and no convention is imposed on null
reported actions.

\emph{\textup{(DP-rec)} Record post-processing.}  For every channel
$P:A\to\Delta(X)$, every stochastic $T:X\to\Delta(X')$, and every
$q\in\Delta(A)$, define
\[
(PT)(x'\mid a):=\sum_{x\in X}P(x\mid a)T(x'\mid x).
\]
Then
\[
q^0\succeq_{P\sqcup PT}q^1.
\]

\emph{\textup{(DP-act)} Bayesian action coarsening.}  For every
$P:A\to\Delta(X)$, every stochastic $S:A\to\Delta(A')$, every
$q\in\Delta(A)$, write
\[
(qS)(a'):=\sum_{a\in A}q(a)S(a'\mid a).
\]
For every channel $\widehat P:A'\to\Delta(X)$ satisfying
\begin{equation}
(qS)(a')\widehat P(x\mid a')
=\sum_{a\in A}q(a)P(x\mid a)S(a'\mid a)
\quad\text{for every }(a',x)\in A'\times X,
\tag{DP-act-law}\label{pt:eq:dpact}
\end{equation}
we have
\[
q^0\succeq_{P\sqcup \widehat P}(qS)^1.
\]
When $qS$ has full support, \eqref{pt:eq:dpact} determines $\widehat P$
uniquely; otherwise the clause applies to every channel consistent with it.
In particular, bijective relabellings are neutral by applying
\textup{(DP-act)} in both directions.

The act whose trace is evaluated may be reported through a possibly noisy description, and the visible record may be
garbled after it is generated.  Neither operation can create observable trace: the original procedure is weakly more
traceable than any procedure obtained by either operation.

\begin{remark}\label{pt:rem:dpform}
The two primitive clauses imply the convenient simultaneous form used in a
few calculations.  Let $S:A\to\Delta(A')$ and $T:X\to\Delta(X')$, and let
$P':A'\to\Delta(X')$ satisfy
\begin{equation}
(qS)(a')\,P'(x'\mid a')
=
\sum_{a\in A}\sum_{x\in X}q(a)P(x\mid a)S(a'\mid a)T(x'\mid x)
\quad\text{for every }(a',x')\in A'\times X'.
\tag{DP}\label{pt:eq:dpjoint}
\end{equation}
First apply \textup{(DP-rec)} to form $PT$, and then apply
\textup{(DP-act)} to $(q,PT)$ through $S$.  Equation~\eqref{pt:eq:dpjoint}
says exactly that the supplied $P'$ is an allowed Bayesian completion of this
second step, including on null reported actions.  Chain the two comparisons in
a common three-block environment using Conditions \textup{(P1)} and
\textup{(P3)}.  This yields
$q^0\succeq_{P\sqcup P'}(qS)^1$.  Conversely, the simultaneous form
specialises to each named clause.  Thus the two presentations are equivalent
under \textup{(P1)} and \textup{(P3)}, but the proposition below takes the two
named clauses as its literal premise.  If $(qS)(a')=0$, both
\eqref{pt:eq:dpact} and \eqref{pt:eq:dpjoint} leave the corresponding channel
row unrestricted; null reports require no convention.
\end{remark}

\end{description}

\paragraph{Reached records and posteriors.}
For a lottery $q\in\Delta(A)$ and a channel $P:A\to\Delta(X)$, write
\[
m_{q,P}(x):=\sum_{a\in A}q(a)P(x\mid a),
\qquad
r_x^{\,q,P}(a):=\frac{q(a)P(x\mid a)}{m_{q,P}(x)}
\quad\text{if }m_{q,P}(x)>0
\]
for, respectively, the probability of record $x$ and the posterior belief about the realised action.  Call $x$ reached
when $m_{q,P}(x)>0$; all branch comparisons below are restricted to reached records.

\begin{description}[leftmargin=*,style=sameline]

\item[(P5) Branchwise continuation monotonicity.]
Let $P_1:A\to\Delta(X_1)$ be a first-stage channel and, after each first-stage record $x$, let
$Q^x,R^x:A\to\Delta(X_2^x)$ be two alternative continuation channels.  Write
$P_1\triangleright\{Q^x\}$, and analogously $P_1\triangleright\{R^x\}$, for the compound channel with record alphabet
$\bigcup_{x\in X_1}(\{x\}\times X_2^x)$ that reports both stages:
\[
(P_1\triangleright\{Q^x\})(x,z\mid a):=P_1(x\mid a)Q^x(z\mid a).
\]
Fix $q\in\Delta(A)$ and let
$m(x):=m_{q,P_1}(x)$ and $r_x:=r_x^{\,q,P_1}$ for all branches with $m(x)>0$.

If
\[
r_x^0\succeq_{Q^x\sqcup R^x}r_x^1
\qquad\text{for every }x\text{ with }m(x)>0,
\]
then
\[
q^0
\succeq_{(P_1\triangleright\{Q^x\})\sqcup(P_1\triangleright\{R^x\})}
q^1.
\]
If at least one reached branch comparison is strict, then the aggregate comparison is strict.

Sequential observation is evaluated branch by branch.  After a first-stage signal, if one continuation leaves at least as
much trace as another at every posterior that can be reached, then replacing the continuation branch by branch should not
reduce overall traceability.  If the improvement is strict on a reached branch, the aggregate comparison is strict.

\end{description}

\begin{remark}\label{pt:rem:A5role}
Condition~\textup{(P5)} plays in this system the role the sure-thing principle plays in Savage's: it is a dominance
condition across the branches of a resolving observation, with zero-probability branches unconstrained exactly as
Savage's null events are.  It is also where all cardinal content resides.  The other conditions are ordinal,
environment-by-environment conditions; it is \textup{(P5)}, through its derived consequences---public-coin
independence (Lemma~\ref{pt:lem:publiccoin}), background inertness (Lemma~\ref{pt:lem:background}), and the chain rule
(Lemma~\ref{pt:lem:chain})---that aligns the initially separate scales across priors and forces branch-additive
aggregation.  The strictness clause is not decorative: it is what allows the converse directions of
Lemmas~\ref{pt:lem:publiccoin} and~\ref{pt:lem:background} to be extracted from a dominance condition.
For example, a mutual-information index capped at a fixed level satisfies
\textup{(P1)}--\textup{(P4)} but fails exactly \textup{(P5)}.  Compound observation environments are
accordingly the natural locus of experimental testing: they are where branch-additive and non-additive criteria come
apart.
\end{remark}

% ===============================================================
\subsection{Auxiliary result}
% ===============================================================

Throughout, logarithms are base two, so information is measured in bits; changing the base merely rescales all
information values.
For a finite distribution $x$, write
\[
\Sh(x):=-\sum_z x(z)\log x(z),
\qquad 0\log 0:=0,
\]
and, for a prior--channel pair $(q,P)$, define
\[
I_{q,P}(A;X)
:=
\Sh(m_{q,P})-\sum_{a\in A}q(a)\Sh\bigl(P(\cdot\mid a)\bigr).
\]
When the alphabets are clear, write $I_{q,P}$; write $I$ for the map $(q,P)\mapsto I_{q,P}$.

\begin{proposition}[Pure-trace characterisation]\label{pt:prop:main}
For a family $\{\succeq_P\}_P$ as described above, the following are equivalent.
\begin{enumerate}[label=\textup{(\roman*)},leftmargin=*]
\item The family satisfies Conditions~\textup{(P1)--(P5)}.
\item For every finite channel $P:A\to\Delta(X)$ and every $q,q'\in\Delta(A)$,
\[
q\succeq_P q'
\quad\Longleftrightarrow\quad
I_{q,P}(A;X)\ge I_{q',P}(A;X).
\]
\end{enumerate}
Moreover, under these equivalent conditions, block-supported cross-channel comparisons are represented on the same
scale: for every finite block environment $\bigsqcup_{k\in K}P_k$, every two distinct blocks $i,j\in K$, and every
$q_i\in\Delta(A_i)$, $q_j\in\Delta(A_j)$,
\[
q_i^i\succeq_{\bigsqcup_{k\in K}P_k}q_j^j
\quad\Longleftrightarrow\quad
I_{q_i,P_i}(A_i;X_i)\ge I_{q_j,P_j}(A_j;X_j).
\]
\end{proposition}
%</leanpurestatement>

The proof of the equivalence of \textup{(i)} and \textup{(ii)} is in the Appendix; the moreover clause is immediate
from \textup{(ii)}, since a block-supported lottery has constant block label.  The sufficiency argument is the hard part and has three
movements.  First, \emph{quotient}: block coherence, record
post-processing, and a prior-specific Blackwell argument show that at each full-support prior only the induced law of
posterior beliefs about the act matters, and that this quotient inherits closedness from Condition~\textup{(P2)}
(Lemmas~\ref{pt:lem:blockcoh}--\ref{pt:lem:plcont}).  Second, \emph{cardinalisation}: public-coin independence---derived from
Condition~\textup{(P5)}, not assumed---delivers a continuous affine representative on each posterior-law fibre by the
mixture-space theorem; derived background inertness of independent components yields a quasi-additive product form; and
a tangent-space argument makes branch coefficients path-independent, giving a cocycle, a chain rule, and coherent
scales across support faces (Lemmas~\ref{pt:lem:postsep}--\ref{pt:lem:facescales}).  Third, \emph{identification}: the key
step is scale alignment.  The conditions derive it through data processing and branchwise continuation, rather than
through an assumption on product environments.  The initially fibrewise affine indices are thereby placed on one
cardinal scale; a two-grouping argument eliminates the possible product interaction, and Faddeev's recursion then
identifies the value of full revelation with Shannon entropy, hence the index with mutual information
(Lemmas~\ref{pt:lem:scalecoherence}--\ref{pt:lem:globalsketch}).

To state uniqueness, it is useful to make the common scale explicit.  A real-valued index $V$ on all finite
prior--channel pairs is a \emph{common-scale representation} if, for every channel $P:A\to\Delta(X)$ and every
$q,q'\in\Delta(A)$,
\[
q\succeq_Pq'
\quad\Longleftrightarrow\quad
V(q,P)\ge V(q',P),
\]
and, for every finite block environment $\bigsqcup_{k\in K}P_k$, every two distinct blocks $i,j\in K$, and all
$q_i\in\Delta(A_i)$ and $q_j\in\Delta(A_j)$,
\[
q_i^i\succeq_{\bigsqcup_{k\in K}P_k}q_j^j
\quad\Longleftrightarrow\quad
V(q_i,P_i)\ge V(q_j,P_j).
\]
Call a common-scale representation \emph{branch-additive} if, for every finite $A$, every $q\in\Delta(A)$, every
first-stage channel $P_1:A\to\Delta(X_1)$, and every continuation profile
$\{Q^x:A\to\Delta(X_2^x)\}_{x\in X_1}$,
\[
V\bigl(q,P_1\triangleright\{Q^x\}\bigr)
=
V(q,P_1)+\sum_{x:m(x)>0}m(x)\,V(r_x,Q^x),
\]
where $m(x)=m_{q,P_1}(x)$ and $r_x=r_x^{q,P_1}$ on every reached branch.

\begin{proposition}[Uniqueness]\label{pt:prop:uniqueness}
Let $\{\succeq_P\}_P$ satisfy the equivalent conditions of Proposition~\ref{pt:prop:main}.  For any real-valued index $V$ on
all finite prior--channel pairs:
\begin{enumerate}[label=\textup{(\roman*)},leftmargin=*]
\item $V$ is a common-scale representation if and only if $V=\varphi\circ I$ for a single strictly increasing
$\varphi:[0,\infty)\to\R$.
\item A common-scale representation $V$ is branch-additive if and only if $V=aI$ for some $a>0$.
\end{enumerate}
\end{proposition}

The primitive remains ordinal: part \textup{(i)} permits any common strictly increasing transformation.  Once
block-supported procedures are required to be scored by their component-pair values, however, the block comparisons
permit only one transformation across environments.  Branch additivity is not another behavioural condition; it is an
exact accounting requirement on a numerical index.  Part \textup{(ii)} shows that this requirement selects the linear
gauge, just as mixture affinity selects the expected-utility index within its ordinal class.  It fixes both the zero
and the curvature of the information scale, leaving only a positive choice of unit, with bits and nats the standard
base-two and base-$e$ conventions.

This normalization is what makes the additive combination in the main model meaningful.  If the trace
component were represented by a nonlinear transformation $\varphi\circ I$, a coefficient multiplying it would have no
stable interpretation.  Proposition~\ref{pt:prop:uniqueness} shows that the branch-additive gauge is $V=aI$; the factor
$a$ is absorbed into the main coefficient $\lambda$.  Once the payoff-utility scale and the information unit are fixed,
one effective trade-off parameter can therefore be carried across environments.  The proof is at the end of this
appendix.

Uniqueness concerns the scale; a complementary observation concerns the sign.  The condition
system splits into a structural part and an oriented part.  Conditions~\textup{(P2)} and \textup{(P3)} are equivalences or
closedness conditions and are invariant under reversal of every preference in the family; Condition~\textup{(P5)} is
self-dual, since branchwise monotonicity for the reversed family is the original condition with the two continuation
profiles exchanged, and strictness survives the exchange.  Only the orientation clause of Condition~\textup{(P1)} and the
data-processing condition \textup{(P4)} carry a sign.

Call an \emph{orientation} of the system a choice, for each of the two signed clauses---the non-triviality clause of
\textup{(P1)} and Condition~\textup{(P4)}---of either its original or its reversed form, or the replacement of both by
total indifference on the comparisons of the non-triviality clause.

\begin{corollary}[Sign trichotomy]\label{pt:cor:signduality}
Let $\{\succeq_P\}_P$ satisfy the weak-order clause of \textup{(P1)} and Conditions~\textup{(P2)}, \textup{(P3)},
\textup{(P5)}.  Then:
\begin{enumerate}[label=\textup{(\roman*)},leftmargin=*]
\item under the original orientation of both signed clauses, the family is represented by $I_{q,P}$
\textup{(}Proposition~\ref{pt:prop:main}\textup{)};
\item under the reversal of both signed clauses, the family is represented by $-I_{q,P}$;
\item the two mixed orientations are inconsistent: no family satisfies the original \textup{(P4)} together with
reversed non-triviality, or reversed \textup{(P4)} together with original non-triviality.
\end{enumerate}
The totally indifferent family satisfies the structural conditions with both signed clauses replaced by indifference.
Hence the structural conditions admit exactly three coherent orientations, represented by $I$, $-I$, and $0$.
\end{corollary}

\begin{proof}
For (ii), the reversed family $\{\preceq_P\}_P$ satisfies \textup{(P1)}--\textup{(P5)}---the structural conditions are
reversal-invariant and \textup{(P5)} is self-dual as noted above---so Proposition~\ref{pt:prop:main} represents it by
$I_{q,P}$, hence the original family by $-I_{q,P}$.  For (iii), fix a full-support $q$ on a non-singleton $A$ and
apply \textup{(P4)} to $P=\Id_A$, $S=\Id_A$, and a constant record kernel
$T:A\to\Delta(\{\ast\})$.  Then $qS=q$ and the processed channel is $U_A$, so the original orientation gives
\[
q^0\succeq_{\Id_A\sqcup U_A}q^1,
\]
contradicting reversed non-triviality.  The reversed orientation of \textup{(P4)} gives the opposite weak comparison,
contradicting original non-triviality.  The final claim is the necessity verification in the Appendix for (i), its reversal for (ii),
and direct verification for the indifferent family.
\end{proof}

\begin{proposition}[Independence of the pure-record conditions]\label{pt:prop:independence}
None of Conditions~\textup{(P1)}--\textup{(P5)} is implied by the other four: for each condition there is a family of weak
orders that satisfies the remaining four and violates it.
\end{proposition}

The five families are constructed in the Appendix.  Two of them carry behavioural content of their own and mark the
substantive frontier of the system: the Shannon-entropy-scaled criterion $\bigl(1+\Sh(q)\bigr)I_{q,P}$ satisfies every condition
except data processing, valuing the uncertainty of a procedure beyond its trace, and the satiation criterion
$\min\{I_{q,P},c\}$, for any fixed $c>0$, satisfies every condition except branchwise continuation monotonicity, ceasing to value revelation
once $c$ bits are secured.

Taken together, the results measure the freedom the conditions leave, and it amounts to a sign and a unit.  Behaviour
selects the sign: a family satisfying the structural conditions together with an orientation of the two signed clauses is
represented by $I$, by $-I$, or by total indifference (Corollary~\ref{pt:cor:signduality}), and none of the five conditions
securing this rigidity is redundant (Proposition~\ref{pt:prop:independence}).  Convention selects the unit: within an
orientation, the branch-additive representations are exactly the positive multiples of $I$
(Proposition~\ref{pt:prop:uniqueness}).  In this sense the formula is less the content of Proposition~\ref{pt:prop:main} than
the trichotomy is: the conditions leave open not which functional measures traceable agency, but whether the agent seeks
it, shuns it, or is indifferent to it.  The two oriented arms are not symmetric in behavioural content, however; the
Discussion takes up what the negative orientation does and does not deliver.



\paragraph{Embedding into the payoff--record model.}
Fix any $o_0\in O$ and identify an auxiliary record channel
$P:A\to\Delta(X)$ with the joint channel
$K_P(o,x\mid a)=\mathbf 1_{\{o=o_0\}}P(x\mid a)$.  Main Axiom~\textup{(A4)}
is \textup{(DP-rec)}, main Axiom~\textup{(A5)} is \textup{(DP-act)}, and
main Axiom~\textup{(A6)} is \textup{(P5)} because every payoff is constant.
Thus Proposition~\ref{pt:prop:main} applies directly to the restriction used in
Lemma~\ref{lem:trace-engine}.

%<*leanpureproof>
\subsection{Proof of the pure-trace characterization}\label{pt:app:sufficiency}
% ===============================================================

\subsubsection{Necessity}

For every finite channel $P:A\to\Delta(X)$, define
\[
q\succeq^I_P q'
\quad\Longleftrightarrow\quad
I_{q,P}(A;X)\ge I_{q',P}(A;X).
\]
We verify that the family $\{\succeq^I_P\}_P$ satisfies Conditions~\textup{(P1)--(P5)}: this is the implication
\textup{(ii)$\Rightarrow$(i)} of Proposition~\ref{pt:prop:main}, and it establishes the consistency of the condition
system.

The \emph{posterior law} induced by $(q,P)$ is the finite distribution
\[
\mu_{q,P}:=
\sum_{x:m_{q,P}(x)>0}m_{q,P}(x)\,\delta_{r_x^{\,q,P}}
\]
over posterior beliefs.  Bayes plausibility gives $\int r\,d\mu_{q,P}(r)=q$.  We repeatedly use the elementary identities
\[
I_{q,P}
=
\Sh(m_{q,P})-\sum_{a\in A}q(a)\Sh(P(\cdot\mid a)),
\tag{N1}
\]
and
\[
I_{q,P}
=
\Sh(q)-\int_{\Delta(A)}\Sh(r)\,d\mu_{q,P}(r).
\tag{N2}
\]
Both formulae are read on the support of the relevant distributions; zero-probability rows and records do not affect the
value.

\emph{Condition \textup{(P1)}.}
For each fixed $P$, the relation $\succeq^I_P$ is complete and transitive because it is represented by the real-valued
function $q\mapsto I_{q,P}$.  If $|A|\ge2$ and $q$ has full support, then full revelation gives
\[
I_{q,\Id_A}=\Sh(q)>0,
\]
whereas the one-record uninformative channel $U_A$ gives $I_{q,U_A}=0$.  Hence
\[
q^0\succ^I_{\Id_A\sqcup U_A}q^1.
\]

\emph{Condition \textup{(P2)}.}
For fixed finite alphabets, (N1) shows that $(P,q)\mapsto I_{q,P}$ is continuous, since Shannon entropy is continuous on
every finite simplex.  Therefore the set
\[
\{(P,q,q'):q\succeq^I_P q'\}
=
\{(P,q,q'):I_{q,P}\ge I_{q',P}\}
\]
is closed.  Posterior-law continuity of block comparisons is not an condition; for the mutual-information family it also
holds directly, by (N2) and continuity of $\Sh$ on the compact simplex $\Delta(A)$, consistently with its derivation
from the conditions in Lemma~\ref{pt:lem:plcont} below.

\emph{Condition \textup{(P3)}.}
For every channel $P$ and lotteries $q,q'\in\Delta(A)$,
\[
I_{q^0,P\sqcup P}=I_{q,P},
\qquad
I_{(q')^1,P\sqcup P}=I_{q',P},
\]
because a block-supported lottery puts probability one on its own block and the block label is constant under that
lottery.  Hence
\[
q\succeq^I_P q'
\quad\Longleftrightarrow\quad
q^0\succeq^I_{P\sqcup P}(q')^1.
\]
The same observation proves invariance to adding or deleting irrelevant blocks.  If $q_i^i$ and $q_j^j$ are supported on
blocks $i$ and $j$, their mutual-information values in $\bigsqcup_{k\in K}P_k$ are exactly
\[
I_{q_i,P_i}
\quad\text{and}\quad
I_{q_j,P_j},
\]
independently of all other blocks.

\emph{Condition \textup{(P4)}.}
It is useful to verify the stronger simultaneous inequality from
Remark~\ref{pt:rem:dpform}, which immediately supplies both named clauses.
Let $S:A\to\Delta(A')$ and $T:X\to\Delta(X')$ be stochastic, and let $P'$ satisfy \eqref{pt:eq:dpjoint}.  Generate
$A\sim q$; then, conditionally independently given $A$, generate $A'$ from $S(\cdot\mid A)$ and $X$ from
$P(\cdot\mid A)$; finally generate $X'$ from $T(\cdot\mid X)$.  The joint law of $(A,A',X,X')$ is
\[
q(a)S(a'\mid a)P(x\mid a)T(x'\mid x),
\]
so $A'$ is conditionally independent of $(X,X')$ given $A$, and $X'$ is conditionally independent of $(A,A')$ given $X$.
Hence
\[
A'\longleftarrow A\longrightarrow X\longrightarrow X'
\]
is a Markov chain in the sense that $A'\to A\to X\to X'$ satisfies the required conditional independences.  Data
processing applied to the fork $A'\leftarrow A\to X$ gives $I(A;X)\ge I(A';X)$, and data processing applied to
$A'\to X\to X'$ gives $I(A';X)\ge I(A';X')$.  Therefore
\[
I_{q,P}(A;X)\ge I(A';X').
\]
By construction the joint law of $(A',X')$ is exactly $(qS)(a')P'(x'\mid a')$, the left-hand side of
\eqref{pt:eq:dpjoint}, so $I(A';X')=I_{qS,P'}(A';X')$.  Rows $a'$ with $(qS)(a')=0$ carry no mass under $qS$, so this value
is the same for every $P'$ satisfying \eqref{pt:eq:dpjoint}, however the null rows are completed.  Using the block
computation of \textup{(P3)} above to evaluate the two blocks separately,
\[
q^0\succeq^I_{P\sqcup P'}(qS)^1.
\]
The first named clause is the case $S=\Id_A$ and gives
$I_{q,P}(A;X)\ge I_{q,PT}(A;X')$.  For later use, in the action-only case $T=\Id_X$ write $S^qP$ for the conditional
channel from the reported action to the record on every reached row:
\[
S^qP(x\mid a')
:=
\frac{\sum_{a\in A}q(a)S(a'\mid a)P(x\mid a)}{(qS)(a')}
\qquad\text{when }(qS)(a')>0.
\]
A \emph{completion} of $S^qP$ is any channel $\widehat P:A'\to\Delta(X)$ that agrees with this formula on reached rows;
its values on null rows are unrestricted.  The second named clause gives
$I_{q,P}(A;X)\ge I_{qS,\widehat P}(A';X)$ for every completion $\widehat P$ of $S^qP$.

\emph{Condition \textup{(P5)}.}
Fix a first-stage channel $P_1:A\to\Delta(X_1)$.  Let
\[
m(x)=m_{q,P_1}(x),
\qquad
r_x=r_x^{\,q,P_1}.
\]
For a continuation profile $\mathcal Q=\{Q^x\}_x$, the chain rule for mutual information gives
\[
I_{q,P_1\triangleright\mathcal Q}(A;X_1,X_2)
=
I_{q,P_1}(A;X_1)
+
\sum_{x:m(x)>0}m(x)\,I_{r_x,Q^x}(A;X_2^x).
\tag{N3}
\]
Therefore, if
\[
r_x^0\succeq^I_{Q^x\sqcup R^x}r_x^1
\quad\text{for every positive-probability branch }x,
\]
then $I_{r_x,Q^x}\ge I_{r_x,R^x}$ for every such $x$.  Using (N3),
\[
I_{q,P_1\triangleright\{Q^x\}}
\ge
I_{q,P_1\triangleright\{R^x\}},
\]
so
\[
q^0
\succeq^I_{(P_1\triangleright\{Q^x\})\sqcup(P_1\triangleright\{R^x\})}
q^1.
\]
If one branch comparison is strict and that branch has positive probability, then the corresponding summand in (N3) is
strictly larger, so the aggregate comparison is strict.

No product condition is imposed, so the verification ends here.  For later use, write $q_1\otimes q_2$ for the product
lottery and define the product of channels by
\[
(P_1\otimes P_2)(x_1,x_2\mid a_1,a_2)
:=P_1(x_1\mid a_1)P_2(x_2\mid a_2).
\]
The mutual-information family satisfies the background-inertness property derived in Lemma~\ref{pt:lem:background} below,
since
\[
I_{q_1\otimes q_2,P_1\otimes R_2}(A_1,A_2;X_1,X_2)
=
I_{q_1,P_1}(A_1;X_1)+I_{q_2,R_2}(A_2;X_2),
\tag{N4}
\]
so that the comparison of $P_1\otimes R_2$ and $Q_1\otimes R_2$ at $q_1\otimes q_2$ reduces to the comparison of $P_1$
and $Q_1$ at $q_1$, whatever the common background $R_2$.  Consistently with Lemma~\ref{pt:lem:background}, this is a
consequence of the chain rule (N3) and not an independent requirement.

\subsubsection{Sufficiency}

Assume Conditions~\textup{(P1)--(P5)}.  We prove the implication \textup{(i)$\Rightarrow$(ii)} in
Proposition~\ref{pt:prop:main}.  The proof is written first for full-support lotteries.  Boundary cases are obtained by
exact restriction to the positive support and transport back to the original action set.  The key point is that posterior-law sufficiency is not assumed; it
is derived from block comparisons and record post-processing, and so is its continuity
(Lemma~\ref{pt:lem:plcont}).

\paragraph{How to read the proof.}
The argument has five interfaces.  Each stage ends with exactly the object
needed by the next one:
\begin{enumerate}[label=\textup{Stage \arabic*:},leftmargin=*]
\item replace channels by Bayes-plausible posterior laws and derive the needed
  quotient continuity from primitive closedness;
\item cardinalise each fixed-prior quotient and make support, neutrality, and
  relabelling behaviour explicit;
\item transport those affine values through independent products and into one
  cross-prior block scale, allowing provisionally a bilinear interaction;
\item identify reached-branch coefficients, prove their cocycle, and align
  the scales on nested support faces;
\item compare product and sequential decompositions, eliminate the
  interaction, and invoke Faddeev's recursion to obtain Shannon entropy and
  hence mutual information.
\end{enumerate}
The long posterior-grid, product-coefficient, and branch tangent-space
calculations are technical subarguments.  A reader interested in the logical
spine may use their lemma statements as interfaces.  The kernel proof chooses
its canonical representative before the branch and product calculations and
implements the branch interfaces using finite atomic signed posterior laws.
The paper keeps product calibration before branch calibration because this
makes the scalar algebra easier to read.  The early full-revelation
normalisation and exact relabelling lemma below supply the coherence needed by
that presentation; the kernel proof reaches the same later interfaces in the
opposite order.

\paragraph{Stage 1: posterior-law quotient and derived continuity.}
\addcontentsline{toc}{subsubsection}{Stage 1: posterior-law quotient and derived continuity}
The first five lemmas establish a coherent order on posterior laws.  In
particular, posterior-law sufficiency and continuity are conclusions, not
extra axioms; the barycentric spread--merge construction in
Lemma~\ref{pt:lem:plcont} is the technical core of this stage.

\begin{lemma}[Coherent block comparisons]\label{pt:lem:blockcoh}
Assume Conditions~\textup{(P1)} and \textup{(P3)}.

\textup{(a)}  Let $(q_1,P_1)$, $(q_2,P_2)$, $(q_3,P_3)$ be pairs of a lottery $q_k\in\Delta(A_k)$ and a channel
$P_k:A_k\to\Delta(X_k)$, with action sets, record sets, and lotteries not necessarily equal and the lotteries not
necessarily of full support.  Then the two-block comparisons between such pairs are complete,
\[
q_1^0\succeq_{P_1\sqcup P_2}q_2^1
\quad\text{or}\quad
q_2^0\succeq_{P_2\sqcup P_1}q_1^1,
\]
satisfy the reverse-block equivalence
\[
q_2^1\succeq_{P_1\sqcup P_2}q_1^0
\quad\Longleftrightarrow\quad
q_2^0\succeq_{P_2\sqcup P_1}q_1^1,
\]
and are transitive:
\[
q_1^0\succeq_{P_1\sqcup P_2}q_2^1
\quad\text{and}\quad
q_2^0\succeq_{P_2\sqcup P_3}q_3^1
\quad\Longrightarrow\quad
q_1^0\succeq_{P_1\sqcup P_3}q_3^1.
\]
In particular, for a fixed full-support $q\in\Delta(A)$ the comparisons $q^0\succeq_{P\sqcup Q}q^1$ define a weak order
over channels on $A$ at prior $q$, and pairwise comparisons between block-supported lotteries are independent of the
block environment in which the two blocks are embedded.

\textup{(b)}  Under \textup{(P2)}, the comparisons in \textup{(a)} are also closed in the channels at fixed alphabets
and fixed lotteries: if $P_n\to P$ in $\Delta(X_1)^{A_1}$, $Q_n\to Q$ in $\Delta(X_2)^{A_2}$, $q_1\in\Delta(A_1)$ and
$q_2\in\Delta(A_2)$ are fixed, and
\[
q_1^0\succeq_{P_n\sqcup Q_n}q_2^1
\quad\text{for all }n,
\]
then
\[
q_1^0\succeq_{P\sqcup Q}q_2^1.
\]
Fixedness of the alphabets is inherited from Condition~\textup{(P2)}, which is stated at a fixed pair of alphabets.
\end{lemma}

\begin{proof}
\textup{(a)}  Completeness follows from Condition~\textup{(P1)} applied to the two-block environment $P_1\sqcup P_2$: either
$q_1^0\succeq_{P_1\sqcup P_2}q_2^1$ or $q_2^1\succeq_{P_1\sqcup P_2}q_1^0$.  In the second case, applying
Condition~\textup{(P3)} to the two-block environment with blocks $P_1$ and $P_2$, using the ordered pair $(i,j)=(1,0)$,
gives the canonical reverse comparison
\[
q_2^1\succeq_{P_1\sqcup P_2}q_1^0
\quad\Longleftrightarrow\quad
q_2^0\succeq_{P_2\sqcup P_1}q_1^1,
\]
which is the reverse-block equivalence.

For transitivity, form the three-block environment $\bigsqcup_{k=1}^3 P_k$.  By the irrelevant-blocks clause of
Condition~\textup{(P3)} with block pair $(1,2)$, the first hypothesis is equivalent to
\[
q_1^1\succeq_{\bigsqcup_{k=1}^3 P_k}q_2^2,
\]
and by the same clause with block pair $(2,3)$, the second hypothesis is equivalent to
\[
q_2^2\succeq_{\bigsqcup_{k=1}^3 P_k}q_3^3.
\]
Transitivity of the weak order given by Condition~\textup{(P1)} in this common environment gives
\[
q_1^1\succeq_{\bigsqcup_{k=1}^3 P_k}q_3^3,
\]
and applying Condition~\textup{(P3)} again, now deleting the irrelevant middle block, yields
$q_1^0\succeq_{P_1\sqcup P_3}q_3^1$.  Nothing in these arguments requires the pairs to share alphabets or lotteries, or
the lotteries to have full support.  The specialisation to a fixed full-support $q$ and channels on a common $A$ is
immediate.

\textup{(b)}  The block channel $P_n\sqcup Q_n$ has fixed tagged copies of the alphabets $A_1,A_2,X_1,X_2$ and converges
to $P\sqcup Q$ in the finite-dimensional product simplex.  Since
the block-supported lotteries $q_1^0,q_2^1$ are fixed, Condition~\textup{(P2)} applied at this pair of alphabets gives
\[
q_1^0\succeq_{P\sqcup Q}q_2^1.
\qedhere
\]
\end{proof}

\begin{lemma}[Prior-specific Blackwell equivalence]\label{pt:lem:blackwell}
Let $q\in\Delta(A)$ have full support.  Let $P:A\to\Delta(X)$ and $Q:A\to\Delta(Y)$ be finite channels.  If
\[
\mu_{q,P}=\mu_{q,Q},
\]
then there are stochastic kernels $T:X\to\Delta(Y)$ and $T':Y\to\Delta(X)$ such that
\[
Q=PT,
\qquad
P=QT'.
\]
\end{lemma}

\begin{proof}
It is enough to construct $T$; the construction of $T'$ is symmetric.  Let
\[
m_P(x):=m_{q,P}(x),\qquad m_Q(y):=m_{q,Q}(y).
\]
If $m_P(x)=0$, then $P(x\mid a)=0$ for every $a\in A$, because $q$ has full support.  Such records do not affect
$PT$, so define the corresponding row of $T$ arbitrarily.  The same observation applies to null records of $Q$.

Let $\mathcal R$ be the finite set of posterior vectors that occur with positive probability under either channel.  For
$r\in\mathcal R$, define
\[
X_r:=\{x\in X:m_P(x)>0,\ r_x^{\,q,P}=r\},
\qquad
Y_r:=\{y\in Y:m_Q(y)>0,\ r_y^{\,q,Q}=r\}.
\]
The sets $\{X_r\}_{r\in\mathcal R}$ partition the positive-probability records of $P$, and the sets
$\{Y_r\}_{r\in\mathcal R}$ partition the positive-probability records of $Q$, by posterior vector.
Equality of posterior laws implies that the total probability attached to posterior $r$ is the same for the two
channels:
\[
\lambda_r:=\sum_{x\in X_r}m_P(x)=\sum_{y\in Y_r}m_Q(y)>0.
\]
For $x\in X_r$ and $y\in Y_r$, set
\[
T(y\mid x):=\frac{m_Q(y)}{\lambda_r},
\]
and set $T(y\mid x)=0$ when $x\in X_r$ and $y\notin Y_r$.  This defines a stochastic row for every positive-probability
record $x$, since the probabilities over $Y_r$ sum to one.  Rows corresponding to null records of $P$ were already
chosen arbitrarily.

Fix $a\in A$ and $y\in Y_r$.  Since all records in $X_r$ have posterior $r$, Bayes' rule gives
\[
q(a)P(x\mid a)=m_P(x)r(a)
\qquad(x\in X_r).
\]
Therefore
\[
\sum_{x\in X_r}P(x\mid a)
=
\frac{\lambda_r r(a)}{q(a)}.
\]
Using the definition of $T$,
\[
(PT)(y\mid a)
=
\sum_{x\in X_r}P(x\mid a)\frac{m_Q(y)}{\lambda_r}
=
\frac{m_Q(y)r(a)}{q(a)}.
\]
But $y\in Y_r$ also has posterior $r$, so Bayes' rule for $Q$ gives
\[
Q(y\mid a)=\frac{m_Q(y)r(a)}{q(a)}.
\]
Thus $(PT)(y\mid a)=Q(y\mid a)$ for every positive-probability $y$.  If $m_Q(y)=0$, then full support of $q$ implies
$Q(y\mid a)=0$ for every $a$.  No positive-probability record of $P$ sends mass to such a $y$ under the construction;
rows of $T$ indexed by null records of $P$ may be chosen arbitrarily, but those records have $P(x\mid a)=0$ for every
$a$, so they contribute nothing to $PT$.  Hence $(PT)(y\mid a)=0=Q(y\mid a)$ for every $a$, and therefore $Q=PT$.
\end{proof}

\begin{lemma}[Posterior-law sufficiency]\label{pt:lem:plsuff}
Fix a full-support $q\in\Delta(A)$.  If $P:A\to\Delta(X)$ and $Q:A\to\Delta(Y)$ generate the same posterior law at $q$,
\[
\mu_{q,P}=\mu_{q,Q},
\]
then
\[
q^0\sim_{P\sqcup Q}q^1.
\]
Consequently, by Lemma~\ref{pt:lem:blockcoh}, at a fixed full-support prior the block comparison of channels depends only
on the induced posterior law.
\end{lemma}

\begin{proof}
By Lemma~\ref{pt:lem:blackwell}, there are stochastic kernels $T$ and $T'$ such that $Q=PT$ and $P=QT'$.  Clause
\textup{(DP-rec)} of Condition~\textup{(P4)} gives
\[
q^0\succeq_{P\sqcup PT}q^1,
\qquad
q^0\succeq_{Q\sqcup QT'}q^1.
\]
Using $PT=Q$, the first comparison is
\[
q^0\succeq_{P\sqcup Q}q^1.
\]
Using $QT'=P$, the second comparison is
\[
q^0\succeq_{Q\sqcup P}q^1.
\]
By the reverse-block equivalence in Lemma~\ref{pt:lem:blockcoh} (equivalently, Condition~\textup{(P3)} applied to
$P\sqcup Q$ with ordered pair $(i,j)=(1,0)$), this is equivalent to
\[
q^1\succeq_{P\sqcup Q}q^0.
\]
Thus $q^0\sim_{P\sqcup Q}q^1$.
\end{proof}

For channels $P:A\to\Delta(X_P)$ and $R:A\to\Delta(X_R)$ and $\lambda\in[0,1]$, write
$\lambda P\oplus(1-\lambda)R$ for the channel that publicly selects $P$ with probability $\lambda$ and $R$ otherwise,
and reports the selected channel together with its record.  Thus
\[
\bigl(\lambda P\oplus(1-\lambda)R\bigr)((z,i)\mid a)
=
\begin{cases}
\lambda P(z\mid a),&i=0,\\
(1-\lambda)R(z\mid a),&i=1.
\end{cases}
\]
Because the coin is observed,
\[
\mu_{q,\lambda P\oplus(1-\lambda)R}
=\lambda\mu_{q,P}+(1-\lambda)\mu_{q,R}.
\]

\begin{lemma}[Derived public-coin independence]\label{pt:lem:publiccoin}
Assume Conditions~\textup{(P1)} and \textup{(P3)}, clause \textup{(DP-rec)} of Condition~\textup{(P4)}, and
Condition~\textup{(P5)}.  Fix a finite action set $A$, a
full-support $q\in\Delta(A)$, and channels $P,Q,R$ on $A$.  For every $\lambda\in(0,1)$,
\[
q^0\succeq_{P\sqcup Q}q^1
\quad\Longleftrightarrow\quad
q^0\succeq_{(\lambda P\oplus(1-\lambda)R)\sqcup(\lambda Q\oplus(1-\lambda)R)}q^1.
\]
\end{lemma}

\begin{proof}
Throughout, for channels $V,W$ on $A$ write $V\succeq_q W$ for the block comparison $q^0\succeq_{V\sqcup W}q^1$, and
$V\succ_q W$ for the corresponding strict comparison.  By Lemma~\ref{pt:lem:blockcoh} the comparison $\succeq_q$ is
complete and transitive, and by Lemma~\ref{pt:lem:plsuff} it depends only on the pair of posterior laws
$(\mu_{q,V},\mu_{q,W})$.  In particular, channels with equal posterior laws at $q$ may be substituted on either side of
any comparison; we refer to this as \emph{substitution}.  Substitution also preserves strict comparisons: if
$\mu_{q,V}=\mu_{q,V'}$ and $\mu_{q,W}=\mu_{q,W'}$, then $V\sim_q V'$ and $W\sim_q W'$ by Lemma~\ref{pt:lem:plsuff}, and
transitivity of the complete relation $\succeq_q$ gives $V\succ_q W$ if and only if $V'\succ_q W'$.

\emph{Step 1: the public coin is an uninformative first stage.}
Define the constant channel $C:A\to\Delta(\{0,1\})$ by
\[
C(0\mid a):=\lambda,
\qquad
C(1\mid a):=1-\lambda
\qquad(a\in A).
\]
Since $\lambda\in(0,1)$, both branches have positive mass: $m_{q,C}(0)=\lambda$ and $m_{q,C}(1)=1-\lambda$.  Because
the rows of $C$ are constant, Bayes' rule gives, for each $x\in\{0,1\}$,
\[
r_x^{\,q,C}(a)=\frac{q(a)C(x\mid a)}{m_{q,C}(x)}=q(a),
\]
so both branch posteriors equal the prior: $r_0=r_1=q$.

\emph{Step 2: padding the continuations to a common alphabet.}
Condition~\textup{(P5)} requires the two continuation channels on a given branch to share a record alphabet, whereas $P$
and $Q$ may not.  Let $X_{PQ}:=X_P\sqcup X_Q$ and define $\widetilde P,\widetilde Q:A\to\Delta(X_{PQ})$ by
\[
\widetilde P(z\mid a):=
\begin{cases}
P(z\mid a)&z\in X_P,\\
0&z\in X_Q,
\end{cases}
\qquad
\widetilde Q(z\mid a):=
\begin{cases}
0&z\in X_P,\\
Q(z\mid a)&z\in X_Q.
\end{cases}
\]
Padding adds only zero-probability records, so the positive-mass records and their posteriors are unchanged:
$\mu_{q,\widetilde P}=\mu_{q,P}$ and $\mu_{q,\widetilde Q}=\mu_{q,Q}$.  By substitution,
\[
q^0\succeq_{P\sqcup Q}q^1
\quad\Longleftrightarrow\quad
\widetilde P\succeq_q\widetilde Q,
\]
and the same equivalence holds for the strict comparisons.

\emph{Step 3: the compounds implement the public-coin mixtures.}
Consider the two-stage channels with first stage $C$, continuation pair $\widetilde P$ versus $\widetilde Q$ on branch
$0$, and the common continuation $R$ on branch $1$:
\[
C\triangleright\{\widetilde P,R\}
\qquad\text{and}\qquad
C\triangleright\{\widetilde Q,R\},
\]
where the profile notation lists the branch-$0$ and branch-$1$ continuations in order.  By the definition of two-stage
composition,
\[
(C\triangleright\{\widetilde P,R\})(0,z\mid a)=\lambda\,\widetilde P(z\mid a),
\qquad
(C\triangleright\{\widetilde P,R\})(1,z\mid a)=(1-\lambda)\,R(z\mid a).
\]
A record $(0,z)$ has marginal mass $\lambda\,m_{q,\widetilde P}(z)$, and when this is positive the constant factor
$\lambda$ cancels in Bayes' rule, so its posterior is $r_z^{\,q,\widetilde P}$; symmetrically for records $(1,z)$.
Hence
\[
\mu_{q,C\triangleright\{\widetilde P,R\}}
=
\lambda\mu_{q,\widetilde P}+(1-\lambda)\mu_{q,R}
=
\lambda\mu_{q,P}+(1-\lambda)\mu_{q,R}
=
\mu_{q,\lambda P\oplus(1-\lambda)R},
\]
using the public-coin mixing identity, and likewise
$\mu_{q,C\triangleright\{\widetilde Q,R\}}=\mu_{q,\lambda Q\oplus(1-\lambda)R}$.  By substitution, the right-hand
comparison of the lemma is therefore equivalent to
\[
C\triangleright\{\widetilde P,R\}
\succeq_q
C\triangleright\{\widetilde Q,R\}.
\tag{$\ast$}
\]

\emph{Step 4: forward implication.}
Assume $q^0\succeq_{P\sqcup Q}q^1$.  We verify the branch hypotheses of Condition~\textup{(P5)} for the first stage $C$ and
the continuation profiles $\{\widetilde P,R\}$ versus $\{\widetilde Q,R\}$.  Both branches have positive mass and both
branch posteriors equal $q$ (Step~1).  On branch $0$ the hypothesis is
$q^0\succeq_{\widetilde P\sqcup\widetilde Q}q^1$, which holds by Step~2.  On branch $1$ the two profiles share the
continuation $R$, and the hypothesis $q^0\succeq_{R\sqcup R}q^1$ holds because Condition~\textup{(P1)} gives
$q\succeq_R q$ and the duplication clause of Condition~\textup{(P3)} converts it into the two-block form.
Condition~\textup{(P5)} then yields $(\ast)$, hence the mixture comparison by Step~3.

\emph{Step 5: converse, via the strictness clause.}
Assume the mixture comparison holds, so that $(\ast)$ holds by Step~3, and suppose towards a contradiction that
$q^0\succeq_{P\sqcup Q}q^1$ fails.  By completeness and the reverse-block equivalence of Lemma~\ref{pt:lem:blockcoh},
$Q\succ_q P$, and Step~2 gives $\widetilde Q\succ_q\widetilde P$.  Now apply Condition~\textup{(P5)} with the roles of the
two profiles exchanged: the branch-$0$ comparison $\widetilde Q\succ_q\widetilde P$ is strict on a branch of mass
$\lambda>0$, and the branch-$1$ comparison $R\succeq_q R$ holds as in Step~4.  The strictness clause of
Condition~\textup{(P5)} yields
\[
C\triangleright\{\widetilde Q,R\}
\succ_q
C\triangleright\{\widetilde P,R\},
\]
which contradicts $(\ast)$.  Hence $q^0\succeq_{P\sqcup Q}q^1$.
\end{proof}

The lemma shows that public-coin independence is not an independent behavioural commitment: a public coin is exactly a
first-stage observation that reveals nothing about the act, so invariance under common public mixtures is branchwise
continuation monotonicity applied through trivial branch posteriors.  The statement is given at full-support priors,
which is the only case used below (Lemma~\ref{pt:lem:postsep}).

For a finite action set $A$, let $\mathsf X_A$ be the set of finite-support probability measures on $\Delta(A)$.  Write
\[
b(\mu):=\int_{\Delta(A)}r\,d\mu(r),
\qquad
\mathcal M_q:=\{\mu\in\mathsf X_A:b(\mu)=q\}.
\]
Thus $\mathcal M_q$ is the set of finite posterior laws with prior $q$ as their mean, and
$\mu_{q,P}\in\mathcal M_q$ for every channel $P$.

\begin{lemma}[Posterior-law continuity]\label{pt:lem:plcont}
Fix a full-support $q\in\Delta(A)$.  Let $P_n,Q_n,P,Q$ be finite channels on $A$ with
\[
\mu_{q,P_n}\Rightarrow \mu_{q,P},
\qquad
\mu_{q,Q_n}\Rightarrow \mu_{q,Q}.
\]
If
\[
q^0\succeq_{P_n\sqcup Q_n}q^1
\qquad\text{for all }n,
\]
then
\[
q^0\succeq_{P\sqcup Q}q^1.
\]
\end{lemma}

\begin{proof}
Throughout the proof, for channels $V,W$ on $A$ write $V\succeq_q W$ for the block comparison
$q^0\succeq_{V\sqcup W}q^1$.  By Lemma~\ref{pt:lem:blockcoh} this comparison is complete and transitive, and by
Lemma~\ref{pt:lem:plsuff} it depends only on the pair of posterior laws $(\mu_{q,V},\mu_{q,W})$; in particular, channels
with equal posterior laws at $q$ may be substituted on either side of any comparison.  We refer to this as
\emph{substitution}.

For a finite-support Bayes-plausible law
\[
\mu=\sum_{k\in K}\lambda_k\delta_{r_k}\in\mathcal M_q,
\]
with atoms listed so that $\lambda_k>0$, define the canonical implementing channel
$C_\mu:A\to\Delta(K)$ by
\[
C_\mu(k\mid a):=\frac{\lambda_k r_k(a)}{q(a)}.
\]
Full support of $q$ and $\sum_k\lambda_k r_k=q$ make $C_\mu$ a channel with $\mu_{q,C_\mu}=\mu$.

\emph{Step 1: merging posteriors moves weakly down.}
Let $\mu=\sum_{k\in K}\lambda_k\delta_{r_k}\in\mathcal M_q$ and let $\{K_1,\dots,K_m\}$ partition $K$ into nonempty
groups.  The merged law is
\[
\operatorname{mrg}(\mu):=\sum_{j=1}^m w_j\,\delta_{s_j},
\qquad
w_j:=\sum_{k\in K_j}\lambda_k,
\qquad
s_j:=\frac{1}{w_j}\sum_{k\in K_j}\lambda_k r_k.
\]
Barycentres are preserved, so $\operatorname{mrg}(\mu)\in\mathcal M_q$.  Let $T:K\to\Delta(\{1,\dots,m\})$ be the
deterministic kernel with $T(j\mid k)=1$ if and only if $k\in K_j$.  Then
\[
(C_\mu T)(j\mid a)
=
\sum_{k\in K_j}\frac{\lambda_k r_k(a)}{q(a)}
=
\frac{w_j s_j(a)}{q(a)}
=
C_{\operatorname{mrg}(\mu)}(j\mid a),
\]
so clause \textup{(DP-rec)} of Condition~\textup{(P4)} gives
\[
C_\mu\succeq_q C_{\operatorname{mrg}(\mu)}.
\]
Conversely, say that $\mu'\in\mathcal M_q$ is a \emph{spread} of $\mu$ if $\mu'$ is obtained by replacing each atom
$(\lambda_k,r_k)$ of $\mu$ by finitely many atoms $(\lambda_k w_{kj},v_{kj})_j$ with $w_{kj}\ge0$,
$\sum_j w_{kj}=1$, and $\sum_j w_{kj}v_{kj}=r_k$.  Then $\mu=\operatorname{mrg}(\mu')$ for the grouping by $k$, so
\[
C_{\mu'}\succeq_q C_\mu.
\]
Atoms with $w_{kj}=0$ are deleted from the list; and if the same posterior appears in several list entries, any two
listings of the same measure induce the same posterior law at $q$, so the corresponding canonical channels are
interchangeable by substitution.
Merging records can only obscure which act was realised; splitting an atom into a barycentric family is the reverse
operation.

\emph{Step 2: bounded supports.}
We claim the lemma holds under the additional hypothesis that all laws $\mu_{q,P_n}$ and $\mu_{q,Q_n}$ have support of
size at most $K$, for some fixed $K$.  Write
\[
\mu_{q,P_n}=\sum_{k=1}^{K}\lambda_{n,k}\,\delta_{r_{n,k}},
\]
padding with zero-mass atoms ($\lambda_{n,k}=0$, $r_{n,k}:=q$) so that exactly $K$ terms appear, and define the
$K$-record channel $C_n:A\to\Delta(\{1,\dots,K\})$ by
\[
C_n(k\mid a):=\frac{\lambda_{n,k}\,r_{n,k}(a)}{q(a)}.
\]
Each row sums to one because $\sum_k\lambda_{n,k}r_{n,k}=q$, zero-mass columns contribute nothing, and
$\mu_{q,C_n}=\mu_{q,P_n}$.  Define $D_n$ from $\mu_{q,Q_n}$ in the same way.  By substitution,
\[
C_n\succeq_q D_n
\qquad\text{for all }n.
\]
The parameter sequences of $C_n$ and $D_n$ jointly live in the compact metric space
$([0,1]\times\Delta(A))^{2K}$; pass to a single subsequence along which $\lambda_{n,k}\to\lambda_k^\ast$ and
$r_{n,k}\to r_k^\ast$ for every $k$, simultaneously for both channels.  Then $C_n\to C^\ast$ entrywise, where
$C^\ast(k\mid a)=\lambda_k^\ast r_k^\ast(a)/q(a)$; each row of $C^\ast$ sums to one as an entrywise limit of rows
summing to one, so $C^\ast$ is a channel.  For every continuous $f$ on $\Delta(A)$,
\[
\int f\,d\mu_{q,P_n}
=
\sum_{k=1}^{K}\lambda_{n,k}f(r_{n,k})
\longrightarrow
\sum_{k=1}^{K}\lambda_k^\ast f(r_k^\ast),
\]
so by uniqueness of weak limits
\[
\mu_{q,P}=\sum_{k=1}^{K}\lambda_k^\ast\,\delta_{r_k^\ast}
=\mu_{q,C^\ast},
\]
where the second equality holds because record $k$ of $C^\ast$ has probability $\lambda_k^\ast$ and, when
$\lambda_k^\ast>0$, posterior $r_k^\ast$.  Similarly $D_n\to D^\ast$ entrywise with
$\mu_{q,D^\ast}=\mu_{q,Q}$.  The closedness clause, Lemma~\ref{pt:lem:blockcoh}(b), applied along the subsequence gives
\[
C^\ast\succeq_q D^\ast,
\]
and substitution returns $q^0\succeq_{P\sqcup Q}q^1$.  This proves the bounded-support case; note that only the
primitive closedness of Condition~\textup{(P2)} has been used, through Lemma~\ref{pt:lem:blockcoh}.

\emph{Step 3: the general case, by an $\varepsilon$-grid sandwich.}
Write $\mu_n:=\mu_{q,P_n}$, $\nu_n:=\mu_{q,Q_n}$, $\mu:=\mu_{q,P}$, $\nu:=\mu_{q,Q}$.  The only remaining gap is that
the supports of $\mu_n,\nu_n$ need not be bounded.  Fix $\varepsilon>0$.  Two finite partitions of $\Delta(A)$ are
used, one for each surrogate; they play independent roles and need not be related.  For the spread, fix a finite
triangulation $\mathcal T$ of $\Delta(A)$ into simplices of diameter at most $\varepsilon$, together with a measurable
partition $\{\widetilde S:S\in\mathcal T\}$ of $\Delta(A)$ with $\widetilde S\subseteq S$, assigning each shared face
to exactly one incident cell.  For the merge, fix a finite Borel partition $\{E_1,\dots,E_M\}$ of $\Delta(A)$ into
sets of diameter at most $\varepsilon$ whose relative boundaries carry no $\nu$-mass: cover $\Delta(A)$ by finitely
many open balls of radius at most $\varepsilon/2$, choosing the radii so that no atom of $\nu$ lies on a bounding
sphere ($\nu$ has finitely many atoms, so all but finitely many radii work), and disjointify.

For finite-support $\rho\in\mathcal M_q$ define two bounded-support surrogates.  Both remain in
$\mathcal M_q$: the spread operation replaces each posterior by a barycentric decomposition of itself, while the merge
operation replaces a group of posteriors by their conditional mean, weighted by the total mass of the group.  Cells with
zero mass are ignored (or, equivalently, padded by an arbitrary posterior with zero weight when a fixed record alphabet is
convenient).
The \emph{grid spread} $\operatorname{sp}_\varepsilon(\rho)$ replaces each atom $(\lambda,\beta)$ with
$\beta\in\widetilde S$ by the atoms $\bigl(\lambda\,w_v(\beta),\,v\bigr)_{v}$, where $v$ ranges over the vertices of
$S$ and $(w_v(\beta))_v$ are the barycentric coordinates of $\beta$ in $S$.  It is a spread of $\rho$ in the sense of
Step~1, is supported on the fixed vertex set of $\mathcal T$, and satisfies
$C_{\operatorname{sp}_\varepsilon(\rho)}\succeq_q C_\rho$.
The \emph{cell merge} $\operatorname{mg}_\varepsilon(\rho)$ merges the atoms of $\rho$ within each class $E_j$ of the
Borel partition; it has at most one atom per class and satisfies
$C_\rho\succeq_q C_{\operatorname{mg}_\varepsilon(\rho)}$ by Step~1, since Step~1 allows any grouping of atoms.

By substitution and transitivity (Lemma~\ref{pt:lem:blockcoh}), the hypothesis $P_n\succeq_q Q_n$ gives the chain
\[
C_{\operatorname{sp}_\varepsilon(\mu_n)}
\succeq_q
C_{\mu_n}
\succeq_q
C_{\nu_n}
\succeq_q
C_{\operatorname{mg}_\varepsilon(\nu_n)}.
\]

Both surrogate sequences converge weakly on fixed alphabets.  For the spread side, let $f$ be continuous on
$\Delta(A)$.  The triangulation defines a continuous piecewise-affine interpolant $I_{\mathcal T}f$ by
\[
I_{\mathcal T}f(r):=\sum_{v\in\operatorname{vert}(S)}w_v(r)f(v)
\qquad\text{for }r\in S.
\]
This is well-defined: on a shared face the barycentric coordinates computed in the two incident simplices put zero
weight on vertices outside the face and coincide with the intrinsic barycentric coordinates on that face.  Hence
$I_{\mathcal T}f$ is continuous on all of $\Delta(A)$.  Moreover, for every finite-support $\rho$,
\[
\int f\,d\operatorname{sp}_\varepsilon(\rho)
= \int I_{\mathcal T}f\,d\rho,
\]
because each posterior is replaced by the barycentric lottery over the vertices of its simplex.  Thus
$\mu_n\Rightarrow\mu$ gives
$\operatorname{sp}_\varepsilon(\mu_n)\Rightarrow\operatorname{sp}_\varepsilon(\mu)$.  For the merge side, each class
$E_j$ has relative boundary contained in the chosen bounding spheres, which carry no $\nu$-mass, so $E_j$ is a
$\nu$-continuity set; therefore
$\nu_n(E_j)\to\nu(E_j)$ and, since the truncated map $r\mapsto r\,\1_{E_j}(r)$ is bounded and its discontinuity set
lies in the relative boundary of $E_j$, which is $\nu$-null,
$\int_{E_j}r\,d\nu_n\to\int_{E_j}r\,d\nu$ for every
class.  Classes with $\nu(E_j)>0$ thus have converging merged masses and merged locations.  If $\nu(E_j)=0$, then the
merged mass tends to zero, so the chosen location in that class is immaterial for weak convergence; padding such a class
by any fixed posterior only changes a zero-weight atom.  Hence
$\operatorname{mg}_\varepsilon(\nu_n)\Rightarrow\operatorname{mg}_\varepsilon(\nu)$.

The laws $\operatorname{sp}_\varepsilon(\mu_n)$ and $\operatorname{mg}_\varepsilon(\nu_n)$ have supports of size at
most $N(\varepsilon)$, the maximum of the number of vertices of $\mathcal T$ and the number $M$ of classes of the
Borel partition.  Step~2 applied
to the chain's endpoints therefore yields
\[
C_{\operatorname{sp}_\varepsilon(\mu)}
\succeq_q
C_{\operatorname{mg}_\varepsilon(\nu)}.
\]

Finally, send $\varepsilon\to0$ along a sequence $\varepsilon_m\downarrow0$ with triangulations and Borel partitions
chosen as above.  The
law $\operatorname{sp}_{\varepsilon_m}(\mu)$ has at most $k_P\,|A|$ atoms, where $k_P$ is the number of atoms of
$\mu$: each atom splits among the $|A|$ vertices of one cell.  Each new atom lies within $\varepsilon_m$ of the atom
it came from and group masses are preserved, so $\operatorname{sp}_{\varepsilon_m}(\mu)\Rightarrow\mu$.  The law
$\operatorname{mg}_{\varepsilon_m}(\nu)$ has at most $k_Q$ atoms, where $k_Q$ is the number of atoms of $\nu$: only
classes meeting the support of $\nu$ carry mass.  Once $\varepsilon_m$ is smaller than the minimal distance between
distinct atoms of $\nu$, each class contains at most one atom of $\nu$ and
$\operatorname{mg}_{\varepsilon_m}(\nu)=\nu$ exactly.  The comparisons
$C_{\operatorname{sp}_{\varepsilon_m}(\mu)}\succeq_q C_{\operatorname{mg}_{\varepsilon_m}(\nu)}$ hold for every $m$,
the supports are uniformly bounded, and the laws converge weakly to $\mu$ and $\nu$; one further application of
Step~2 gives $C_\mu\succeq_q C_\nu$.  Substitution returns
\[
q^0\succeq_{P\sqcup Q}q^1.
\qedhere
\]
\end{proof}

The lemma shows that continuity of block comparisons in posterior-law convergence is not an independent behavioural
commitment: it is the primitive closedness of Condition~\textup{(P2)} seen through the quotient that
Condition~\textup{(P3)} and clause \textup{(DP-rec)} of Condition~\textup{(P4)} impose, with the convex order supplying compactness where record alphabets
grow without bound.  Only weak preference is closed in the limit; no strictness claim is made or used.

\paragraph{Stage 2: affine fibres and canonical boundary behaviour.}
\addcontentsline{toc}{subsubsection}{Stage 2: affine fibres and canonical boundary behaviour}
Public-coin independence and quotient continuity now permit
Herstein--Milnor cardinalisation.  The next lemmas record the finite integral
form, normalize no information to zero, and make action coarsening, exact
support restriction, and undetectable distinctions available before values
are compared across priors.  In the formal development these facts are
packaged by a full-revelation-normalized selected representative.  The paper
now makes that same initial selection explicitly; a later coherent product
gauge is declared when it changes the representative.

\begin{lemma}[Posterior-separable representation]\label{pt:lem:postsep}
Fix a full-support $q\in\Delta(A)$.  There is a continuous affine functional
\[
F_q:\mathcal M_q\to\R
\]
such that, for any two channels $P,Q$ on $A$,
\[
q^0\succeq_{P\sqcup Q}q^1
\quad\Longleftrightarrow\quad
F_q(\mu_{q,P})\ge F_q(\mu_{q,Q}).
\]
Moreover, for some continuous $\phi_q:\Delta(A)\to\R$,
\[
F_q(\mu)=\int_{\Delta(A)}\phi_q(r)\,d\mu(r),
\]
and therefore
\[
F_q(\mu_{q,P})
=
\sum_{x:m_{q,P}(x)>0}m_{q,P}(x)\,\phi_q(r_x^{\,q,P}).
\]
The affine functional is unique up to positive affine transformations.  The integrand $\phi_q$ is not pointwise unique;
adding an affine function of $r$ leaves its integral on the fibre $\mathcal M_q$ changed only by a constant.
\end{lemma}

\begin{proof}
First note that every finite-support Bayes-plausible law is implemented by a finite channel.  If
\[
\mu=\sum_{k=1}^n\lambda_k\delta_{r_k}\in\mathcal M_q,
\]
define
\[
P(k\mid a)=\frac{\lambda_k r_k(a)}{q(a)}.
\]
Since $q$ has full support and $\sum_k\lambda_k r_k=q$, this is a channel.  Its record probability at $k$ is
$\lambda_k$, and the posterior after $k$ is $r_k$; hence $\mu_{q,P}=\mu$.

Define a comparison on $\mathcal M_q$ as follows: for $\mu,\nu\in\mathcal M_q$, choose any channels $P,Q$ implementing
them and compare the two block-supported copies of $q$ in $P\sqcup Q$.  This comparison is independent of the chosen
implementations.  Indeed, if $P,P'$ both implement $\mu$ and $Q,Q'$ both implement $\nu$, then
Lemma~\ref{pt:lem:plsuff} gives
\[
q^0\sim_{P\sqcup P'}q^1,
\qquad
q^0\sim_{Q\sqcup Q'}q^1.
\]
Placing the relevant blocks in common block environments and using Lemma~\ref{pt:lem:blockcoh}, transitivity gives
\[
q^0\succeq_{P\sqcup Q}q^1
\quad\Longleftrightarrow\quad
q^0\succeq_{P'\sqcup Q'}q^1.
\]
Thus the quotient comparison on $\mathcal M_q$ is well defined; denote it by $\succeq_q$.  This construction is only at
the fixed full-support prior $q$, and every law in $\mathcal M_q$ is finite-support and, by the preceding construction,
implemented by a finite channel.  Lemma~\ref{pt:lem:blockcoh} also gives completeness and transitivity of this quotient
comparison.

The set $\mathcal M_q$ is a mixture space under ordinary convex combination.  It is convex because barycentres are
preserved by convex combination.  The quotient relation $\succeq_q$ is a weak order by Lemma~\ref{pt:lem:blockcoh}.
Public-coin mixtures of implementing channels correspond exactly to convex mixtures of posterior laws:
\[
\mu_{q,\lambda P\oplus(1-\lambda)R}
=
\lambda\mu_{q,P}+(1-\lambda)\mu_{q,R}.
\]
Therefore the derived public-coin independence property (Lemma~\ref{pt:lem:publiccoin}) gives mixture independence: for every $\mu,\nu,\rho\in\mathcal M_q$ and
$\lambda\in(0,1)$,
\[
\mu\succeq_q\nu
\quad\Longleftrightarrow\quad
\lambda\mu+(1-\lambda)\rho
\succeq_q
\lambda\nu+(1-\lambda)\rho.
\]

The relative weak topology on $\mathcal M_q$ is inherited from the weak topology on probability measures over the compact
metric simplex $\Delta(A)$, and convex-mixture paths are continuous in this topology.  Lemma~\ref{pt:lem:plcont} gives
relatively closed upper and lower contour sets.  Indeed, if
$\mu_n\Rightarrow\mu$ with $\mu_n,\mu,z\in\mathcal M_q$ and $\mu_n\succeq_q z$, choose finite implementing channels for
these laws and apply Lemma~\ref{pt:lem:plcont} (with the constant sequence implementing $z$) to get $\mu\succeq_q z$; the
argument for $\{\,\mu:z\succeq_q\mu\,\}$ is the
same.  Therefore, for fixed $\mu,\nu,\rho\in\mathcal M_q$, the sets
\[
\{\lambda\in[0,1]:\lambda\mu+(1-\lambda)\nu\succeq_q\rho\}
\quad\text{and}\quad
\{\lambda\in[0,1]:\rho\succeq_q\lambda\mu+(1-\lambda)\nu\}
\]
are closed.

By the Herstein--Milnor mixture-space theorem \cite{HersteinMilnor1953}, there is a mixture-preserving representation
$F_q:\mathcal M_q\to\R$:
\[
\mu\succeq_q\nu
\quad\Longleftrightarrow\quad
F_q(\mu)\ge F_q(\nu),
\]
and
\[
F_q(\lambda\mu+(1-\lambda)\nu)
=
\lambda F_q(\mu)+(1-\lambda)F_q(\nu).
\]
If the quotient order is nonconstant, this representative is unique up to positive affine transformations.  If the
quotient order is constant, take $F_q\equiv0$.

Moreover, $F_q$ is continuous in the relative weak topology.  For any $z\in\mathcal M_q$, the sets
\[
\{\mu:F_q(\mu)\ge F_q(z)\}
=
\{\mu:\mu\succeq_q z\}
\quad\text{and}\quad
\{\mu:F_q(\mu)\le F_q(z)\}
=
\{\mu:z\succeq_q\mu\}
\]
are relatively closed by the preceding Lemma~\ref{pt:lem:plcont} argument.  Since $F_q$ is affine, its range on the convex set
$\mathcal M_q$ is an interval.  Hence every upper and lower level set of $F_q$ is relatively closed: levels outside the
range are trivial, and levels inside the range equal $F_q(z)$ for some $z\in\mathcal M_q$.  Thus $F_q$ is both upper and
lower semicontinuous, and therefore continuous.  For any channels $P,Q$ on $A$,
\[
q^0\succeq_{P\sqcup Q}q^1
\quad\Longleftrightarrow\quad
F_q(\mu_{q,P})\ge F_q(\mu_{q,Q}).
\]

The no-information posterior law at $q$ is $\delta_q$.  Write
\[
\chi_q:=\sum_{a\in A}q(a)\delta_{e_a}
\]
for the full-revelation posterior law, where $e_a$ assigns probability one to action $a$.

It remains to write the affine functional in posterior-separable form, without leaving the finite-support domain.  Let
$X:=\Delta(A)$.  Choose
\[
0<\varepsilon<\min_{a\in A}q(a).
\]
For each $r\in X$, define
\[
s_r:=\frac{q-\varepsilon r}{1-\varepsilon}\in\Delta(A)
\]
and
\[
\eta_r
:=
\varepsilon\delta_r+(1-\varepsilon)\sum_{a\in A}s_r(a)\delta_{e_a}.
\]
Then $\eta_r\in\mathcal M_q$, because its barycentre is
\[
\varepsilon r+(1-\varepsilon)s_r=q.
\]
Choose constants $c_a$ such that
\[
\sum_{a\in A}q(a)c_a=F_q(\chi_q),
\]
for instance $c_a=F_q(\chi_q)$ for every $a$.  Define
\[
\phi_q(r):=
\frac{
F_q(\eta_r)
-
(1-\varepsilon)\sum_{a\in A}s_r(a)c_a
}{\varepsilon}.
\]
The map $r\mapsto\eta_r$ is weakly continuous and $F_q$ is continuous, so $\phi_q$ is continuous.

Now take any finite-support law
\[
\mu=\sum_{k=1}^n\lambda_k\delta_{r_k}\in\mathcal M_q.
\]
Since $\sum_k\lambda_k r_k=q$, we have
\[
\sum_{k=1}^n\lambda_k s_{r_k}=q.
\]
Therefore
\[
\sum_{k=1}^n\lambda_k\eta_{r_k}
=
\varepsilon\mu+(1-\varepsilon)\chi_q.
\]
By affinity of $F_q$,
\[
\sum_{k=1}^n\lambda_k F_q(\eta_{r_k})
=
F_q\!\left(\varepsilon\mu+(1-\varepsilon)\chi_q\right)
=
\varepsilon F_q(\mu)+(1-\varepsilon)F_q(\chi_q).
\]
Using $\sum_k\lambda_k s_{r_k}=q$ and $\sum_a q(a)c_a=F_q(\chi_q)$ gives
\[
\sum_{k=1}^n\lambda_k\phi_q(r_k)=F_q(\mu).
\]
This is exactly
\[
F_q(\mu)=\int_X\phi_q(r)\,d\mu(r)
\]
for every finite-support $\mu\in\mathcal M_q$.  Applying this to $\mu_{q,P}$ gives the stated finite-sum formula.  The
affine representative $F_q$ is unique up to positive affine transformations by Herstein--Milnor.  The integrand is not
pointwise unique on a fixed barycentre fibre: if $\ell$ is affine in $r$, then $\int\ell(r)\,d\mu(r)=\ell(q)$ for every
$\mu\in\mathcal M_q$.
\end{proof}

\begin{lemma}[Refinement monotonicity and canonical normalisation]\label{pt:lem:convnorm}
Fix a full-support $q\in\Delta(A)$ and let $F_q,\phi_q$ be as in Lemma~\ref{pt:lem:postsep}.  If a finite law
$\mu\in\mathcal M_q$ is obtained from another finite law $\nu\in\mathcal M_q$ by replacing each posterior of $\nu$ by a
finite Bayes-plausible continuation experiment, then
\[
F_q(\mu)\ge F_q(\nu).
\]
Consequently, $\phi_q$ may be chosen convex.  After replacing $\phi_q$ by an observationally equivalent representative
and subtracting a constant from $F_q$, one may impose
\[
\phi_q(q)=F_q(\delta_q)=0,
\qquad
\phi_q(r)\ge 0\quad\text{for every }r\in\Delta(A),
\]
and therefore
\[
F_q(\mu)\ge 0\quad\text{for every }\mu\in\mathcal M_q.
\]
If $|A|\ge2$, the representative can and will be selected uniquely by the
additional full-revelation normalisation
\[
F_q(\chi_q)=1.
\]
If $|A|=1$, then $\mathcal M_q=\{\delta_q\}$ and we select $F_q\equiv0$.
\end{lemma}

\begin{proof}
Write
\[
\nu=\sum_{i=1}^n m_i\delta_{r_i},
\qquad
\mu=\sum_{i=1}^n m_i\nu_i,
\qquad
\nu_i\in\mathcal M_{r_i}.
\]
Choose a channel $P:A\to\Delta(\{1,\dots,n\})$ implementing $\nu$ at prior $q$.  For each $i$, choose a finite
continuation channel $Q^i:A\to\Delta(Z_i)$ implementing $\nu_i$ at prior $r_i$; on actions outside the support of
$r_i$, define the rows arbitrarily.  Those rows are never reached under prior $r_i$ and do not affect $\nu_i$; moreover,
if $r_i(a)=0$, then $q(a)P(i\mid a)=m_i r_i(a)=0$, so they do not affect the compound posterior law at $q$.  The compound
channel
\[
C:=P\triangleright\{Q^i\}_{i=1}^n
\]
implements $\mu$ at prior $q$.  Let $T:\bigsqcup_{i=1}^n(\{i\}\times Z_i)\to\Delta(\{1,\dots,n\})$ be the kernel
$T(i\mid(i,z))=1$; then $CT=P$.  Clause \textup{(DP-rec)} of Condition~\textup{(P4)} applied to $C$ and $T$ gives
\[
q^0\succeq_{C\sqcup P}q^1.
\]
By Lemma~\ref{pt:lem:postsep},
\[
F_q(\mu)\ge F_q(\nu).
\]

To prove convexity, take $r,s_1,\dots,s_k\in\Delta(A)$ with
\[
r=\sum_{j=1}^k\lambda_j s_j,
\qquad
\lambda_j\ge 0,\quad \sum_j\lambda_j=1.
\]
Choose $0<\varepsilon<\min_a q(a)$ and set
\[
t:=\frac{q-\varepsilon r}{1-\varepsilon}.
\]
This belongs to $\Delta(A)$: for each $a$, $q(a)-\varepsilon r(a)\ge q(a)-\varepsilon>0$, and the coordinates sum to one.
The law
\[
\varepsilon\sum_{j=1}^k\lambda_j\delta_{s_j}+(1-\varepsilon)\delta_t
\]
is a refinement of
\[
\varepsilon\delta_r+(1-\varepsilon)\delta_t,
\]
and both belong to $\mathcal M_q$.  Refinement monotonicity and the integral representation give
\[
\varepsilon\sum_{j=1}^k\lambda_j\phi_q(s_j)+(1-\varepsilon)\phi_q(t)
\ge
\varepsilon\phi_q(r)+(1-\varepsilon)\phi_q(t).
\]
Cancelling the common term and dividing by $\varepsilon$ yields
\[
\phi_q(r)\le \sum_{j=1}^k\lambda_j\phi_q(s_j),
\]
so $\phi_q$ is convex.

Since $q$ is in the relative interior of the finite-dimensional simplex, the supporting-hyperplane theorem for finite
convex functions gives a relative subgradient $g\in\R^A$ of the continuous convex function $\phi_q$ at $q$.  Define
\[
\widetilde\phi_q(r):=\phi_q(r)-g\cdot(r-q).
\]
For every $\mu\in\mathcal M_q$,
\[
\int g\cdot(r-q)\,d\mu(r)=g\cdot(b(\mu)-q)=0,
\]
so $\widetilde\phi_q$ represents the same affine functional on $\mathcal M_q$.  The subgradient inequality gives
\[
\widetilde\phi_q(r)\ge \widetilde\phi_q(q)
\qquad\text{for all }r\in\Delta(A).
\]
Finally replace
\[
F_q(\mu)\quad\text{by}\quad F_q(\mu)-F_q(\delta_q)
\]
and replace $\widetilde\phi_q$ by
\[
\widetilde\phi_q-\widetilde\phi_q(q).
\]
Before this final constant shift, the representation with $\widetilde\phi_q$ gives
\[
F_q(\delta_q)=\int \widetilde\phi_q(r)\,d\delta_q(r)=\widetilde\phi_q(q).
\]
Thus the same constant is subtracted from both sides of the integral representation, and subtracting a constant from
$F_q$ preserves all ordinal comparisons on $\mathcal M_q$.  The final normalisation gives
\[
\phi_q(q)=F_q(\delta_q)=0,
\]
and makes $\phi_q$ non-negative on the whole simplex.  Hence
\[
F_q(\mu)=\int\phi_q(r)\,d\mu(r)\ge 0
\]
for every $\mu\in\mathcal M_q$.

Suppose now that $|A|\ge2$.  Full revelation implements $\chi_q$ and no
information implements $\delta_q$.  The strict local-orientation clause of
Condition~\textup{(P1)}, transported through the quotient comparison, gives
\[
F_q(\chi_q)>F_q(\delta_q)=0.
\]
Divide both $F_q$ and $\phi_q$ by the positive number $F_q(\chi_q)$.  This
preserves the integral representation, all comparisons, convexity, and
non-negativity, and gives $F_q(\chi_q)=1$.  These two anchor values make the
selected affine representative unique.  If $|A|=1$, every posterior law with
barycentre $q$ is $\delta_q$, so the zero representative is the only relevant
choice.
\end{proof}

For a finitely supported posterior law $\mu=\sum_i\lambda_i\delta_{r_i}$ and a bijection $\pi:A\to A'$, write
\[
\mu\pi:=\sum_i\lambda_i\delta_{r_i\pi}
\]
for the posterior law obtained by relabelling actions through $\pi$.

\begin{lemma}[Exact cross-fibre relabelling]\label{pt:lem:actionbase}
Action coarsening implies ordinal cross-fibre equivalence for bijections.  The
zero/full-revelation selection in Lemma~\ref{pt:lem:convnorm} makes that
equivalence exact: for every bijection $\pi:A\to A'$ and full-support
$q\in\Delta(A)$,
\[
F_{q\pi}(\mu\pi)=F_q(\mu)
\]
for all $\mu\in\mathcal M_q$.  In particular, no-information experiments
have the same selected value zero on every finite action set.
\end{lemma}

\begin{proof}
\emph{Ordinal equivalence.}  Let $\pi:A\to A'$ be a bijection.  View $\pi$ as the deterministic kernel
$S(a'\mid a)=\1_{a'=\pi(a)}$.  For any channel $P:A\to\Delta(X)$ and full-support $q\in\Delta(A)$, the Bayesian
pushforward satisfies $S^qP=P\pi^{-1}$ (explicitly: $(S^qP)(x\mid a')=P(x\mid\pi^{-1}(a'))$).  Clause
\textup{(DP-act)} of Condition~\textup{(P4)} gives
\[
q^0\succeq_{P\sqcup P\pi^{-1}}(q\pi)^1.
\]
For the reverse inequality, apply \textup{(DP-act)} to the bijection $\pi^{-1}:A'\to A$, the prior $q\pi\in\Delta(A')$,
and the channel $P\pi^{-1}:A'\to\Delta(X)$.  Then $(\pi^{-1})^{q\pi}(P\pi^{-1})=P$ (the inverse bijection undoes the
relabeling), and $(q\pi)\pi^{-1}=q$.  Hence
\[
(q\pi)^0\succeq_{P\pi^{-1}\sqcup P}q^1.
\]
By the irrelevant-blocks clause of Condition~\textup{(P3)}, block comparisons depend only on the pair of channels involved, not on
numeric labels.  The comparison $(q\pi)^0\succeq_{P\pi^{-1}\sqcup P}q^1$ asks whether $(q\pi,P\pi^{-1})$ is at least
as good as $(q,P)$.  Swapping the block labels (so that $P$ becomes block~0 and $P\pi^{-1}$ becomes block~1) gives
the equivalent statement $q^0\preceq_{P\sqcup P\pi^{-1}}(q\pi)^1$.  Combining with the first inequality gives
\[
q^0\sim_{P\sqcup P\pi^{-1}}(q\pi)^1.
\]
If $P$ implements $\mu\in\mathcal M_q$, then $P\pi^{-1}$ implements $\mu\pi\in\mathcal M_{q\pi}$: for every
positive-probability record $x$,
\[
m_{q\pi,P\pi^{-1}}(x)=m_{q,P}(x),
\qquad
r_x^{\,q\pi,P\pi^{-1}}=r_x^{\,q,P}\pi.
\]
Now take arbitrary $\mu,\nu\in\mathcal M_q$, and choose channels $P,Q$ implementing them at prior $q$.  Let $E$ be the
four-block environment with blocks $P$, $P\pi^{-1}$, $Q$, and $Q\pi^{-1}$.  Write $x,x',y,y'$ for the block-supported
lotteries $q$, $q\pi$, $q$, and $q\pi$ on these four blocks, respectively.  By Lemma~\ref{pt:lem:blockcoh}, the two
two-block indifferences give $x\sim_E x'$ and $y\sim_E y'$.  Hence, by transitivity of the weak order on $E$,
\[
x\succeq_E y
\quad\Longleftrightarrow\quad
x'\succeq_E y'.
\]
Applying Lemma~\ref{pt:lem:blockcoh} again to delete irrelevant blocks gives
\[
q^0\succeq_{P\sqcup Q}q^1
\quad\Longleftrightarrow\quad
(q\pi)^0\succeq_{P\pi^{-1}\sqcup Q\pi^{-1}}(q\pi)^1.
\]
By Lemma~\ref{pt:lem:postsep} on the two fibres, this is equivalent to
\[
F_q(\mu)\ge F_q(\nu)
\quad\Longleftrightarrow\quad
F_{q\pi}(\mu\pi)\ge F_{q\pi}(\nu\pi).
\]
The map $\mu\mapsto\mu\pi$ is an affine bijection from $\mathcal M_q$ to $\mathcal M_{q\pi}$, with inverse induced by
$\pi^{-1}$.  Therefore $G(\mu):=F_{q\pi}(\mu\pi)$ is a continuous affine representative of the same weak order on
$\mathcal M_q$ as $F_q$.  Herstein--Milnor uniqueness gives $G=aF_q+b$ with $a>0$.  Since
$\delta_q\pi=\delta_{q\pi}$ and Lemma~\ref{pt:lem:convnorm} gives
\[
F_q(\delta_q)=F_{q\pi}(\delta_{q\pi})=0,
\]
we have $b=0$.  If $|A|=1$, both representatives vanish identically.  If
$|A|\ge2$, relabelling sends full revelation to full revelation,
$\chi_q\pi=\chi_{q\pi}$, while the canonical selection gives
$F_q(\chi_q)=F_{q\pi}(\chi_{q\pi})=1$.  Substitution in $G=aF_q$ therefore
gives $a=1$.  Hence
\[
F_{q\pi}(\mu\pi)=F_q(\mu)
\]
for all $\mu\in\mathcal M_q$.

\emph{No-information equivalence.}  By Lemma~\ref{pt:lem:convnorm}, $F_q(\delta_q)=0$ for every full-support $q$ on every
finite action set, independently of cardinality.
\end{proof}

\begin{lemma}[Undetectable distinctions neutrality]\label{pt:lem:neutrality}
Let $S:A\to A'$ be an onto deterministic partition map, and let $P:A\to\Delta(X)$ be $S$-measurable:
\[
S(a)=S(b)\quad\Rightarrow\quad P(\cdot\mid a)=P(\cdot\mid b).
\]
Then for every full-support $q\in\Delta(A)$,
\[
q^0\sim_{P\sqcup S^qP}(qS)^1.
\]
In particular, if $G_S:A\to\Delta(A')$ reveals the block $S(a)$, then
\[
(q,G_S)\sim(qS,\Id_{A'}).
\]
\end{lemma}

\begin{proof}
Since $S$ is onto and $q$ has full support, $qS$ has full support on $A'$.  Thus the Bayesian pushforward $S^qP$ has no
zero-row ambiguity in this lemma.  View the deterministic map $S$ as its associated stochastic kernel.
One direction is exactly clause \textup{(DP-act)} of Condition~\textup{(P4)}:
\[
q^0\succeq_{P\sqcup S^qP}(qS)^1.
\]
For the reverse direction, define the Bayesian refinement kernel $R:A'\to\Delta(A)$ by
\[
R(a\mid a')=
\begin{cases}
q(a)/(qS)(a'),&S(a)=a',\\
0,&\text{otherwise}.
\end{cases}
\]
For each $a'\in A'$, $R(\cdot\mid a')$ is stochastic because
\[
\sum_{a\in A}R(a\mid a')
=
\frac{\sum_{a:S(a)=a'}q(a)}{(qS)(a')}
=1.
\]
Moreover, for every $a\in A$,
\[
((qS)R)(a)
=
\sum_{a'\in A'}(qS)(a')R(a\mid a')
=
(qS)(S(a))\frac{q(a)}{(qS)(S(a))}
=q(a)>0.
\]
Thus $(qS)R=q$.

Since $P$ is $S$-measurable, the row of $S^qP$ at each block $a'$ equals the common $P$-row on that block: if $S(a)=a'$,
then
\[
S^qP(x\mid a')
=
\frac{\sum_{b:S(b)=a'}q(b)P(x\mid b)}{(qS)(a')}
=
P(x\mid a).
\]
Now pushing $S^qP$ back through $R$ recovers $P$.  For each $a\in A$,
\[
\bigl(R^{qS}(S^qP)\bigr)(x\mid a)
=
\frac{\sum_{a'\in A'}(qS)(a')R(a\mid a')S^qP(x\mid a')}{((qS)R)(a)}
=
S^qP(x\mid S(a))
=
P(x\mid a).
\]
Therefore
\[
R^{qS}(S^qP)=P.
\]

Apply \textup{(DP-act)} to $(qS,S^qP)$ with the refinement kernel $R$:
\[
(qS)^0\succeq_{S^qP\sqcup P}q^1.
\]
By Condition~\textup{(P3)} applied to the two-block environment $P\sqcup S^qP$ with ordered pair $(1,0)$, this is equivalent
to
\[
(qS)^1\succeq_{P\sqcup S^qP}q^0.
\]
Together the two inequalities give the desired indifference.
\end{proof}

For a nonempty $B\subseteq A$ and a finite law $\mu$ whose barycentre is supported on $B$, non-negativity implies that
every atom $p$ of $\mu$ is supported on $B$.  Define its support projection by
\[
\pi_B(\mu)
:=
\sum_{p\in\operatorname{supp}(\mu)}\mu(\{p\})\,\delta_{p|_B}.
\]
It is a finite law on $\Delta(B)$ with barycentre $b(\mu)|_B$.

\begin{lemma}[Support restriction]\label{pt:lem:supprestrict}
Let $r\in\Delta(A)$ have support $B:=\operatorname{supp}(r)\subsetneq A$.  Then:
\begin{enumerate}
\item[(i)] \emph{Posterior-law reduction.}  For any channel $P:A\to\Delta(X)$, the posterior law $\mu_{r,P}$ depends only
on the rows $P(\cdot\mid b)$ for $b\in B$.  Rows $P(\cdot\mid a)$ with $a\notin B$ have zero contribution because $r(a)=0$.
\item[(ii)] \emph{Preference equivalence.}  Two channels $P,Q:A\to\Delta(X)$ that agree on $B$ (i.e.\ $P(\cdot\mid b)=Q(\cdot\mid b)$
for all $b\in B$) are indifferent at prior $r$.
\item[(iii)] \emph{Face identification.}  For every channel $P:A\to\Delta(X)$,
\[
(r,P)\sim(r|_B,P|_B)
\]
as a two-block comparison; consequently the comparison between channels $P,Q$ at prior $r$ agrees with the comparison
between $P|_B,Q|_B$ at the full-support prior $r|_B\in\Delta(B)$.  This is proved by projecting $A$ to $B$ and embedding
$B$ back into $A$, using the all-completions form of clause \textup{(DP-act)} of Condition~\textup{(P4)} for the embedding.
\item[(iv)] \emph{Representative.}  Define $F_r:\mathcal M_r\to\R$ by $F_r(\mu):=F_{r|_B}^B(\pi_B(\mu))$, where $\pi_B$
projects posterior laws from $\Delta(A)$ to $\Delta(B)$ and $F_{r|_B}^B$ is the Lemma~\ref{pt:lem:postsep} representative
on $\Delta(B)$.  This $F_r$ represents continuation comparisons at prior $r$ in action set $A$.
\end{enumerate}
\end{lemma}

\begin{proof}
In this proof, $(s,R)\succeq(t,T)$ abbreviates the block comparison
$s^0\succeq_{R\sqcup T}t^1$, and $\sim$ denotes both weak inequalities.  These comparisons involve different priors and
different action sets; they chain by transitivity in Lemma~\ref{pt:lem:blockcoh}(a), which is stated across such pairs.

Part~(i): The marginal probability $m_{r,P}(x)=\sum_a r(a)P(x\mid a)$ and the Bayes numerator $r(a)P(x\mid a)$ are both
zero for $a\notin B$.  Hence $\mu_{r,P}$ is determined by $\{P(\cdot\mid b):b\in B\}$.

Part~(ii): Let $P$ and $\widetilde P$ be channels on $A$ that agree on $B$.  We first show
$(r,P)\sim(r,\widetilde P)$.  Fix $b_0\in B$, and let $S:A\to B$ be the deterministic map with $S(b)=b$ for $b\in B$
and $S(a)=b_0$ for $a\notin B$.  Then $rS=r|_B$ and the Bayesian pushforward $S^rP$ is $P|_B$ on $B$.  Clause \textup{(DP-act)}
gives
\[
(r,P)\succeq(r|_B,P|_B).
\]
Let $I:B\hookrightarrow A$ be the inclusion.  Then $(r|_B)I=r$, and the Bayesian pushforward $I^{r|_B}(P|_B)$ is
pinned down on $B$ and unconstrained on $A\setminus B$.  By the all-completions convention in \textup{(DP-act)}, choose
the completion to be $\widetilde P$.  Clause \textup{(DP-act)} gives
\[
(r|_B,P|_B)\succeq(r,\widetilde P).
\]
Thus $(r,P)\succeq(r,\widetilde P)$.  Applying the same projection/inclusion argument with $\widetilde P$ as the original
channel and choosing $P$ as the inclusion completion gives $(r,\widetilde P)\succeq(r,P)$.  Hence
$(r,P)\sim(r,\widetilde P)$.

If $P,Q$ agree on $B$, choose a common extension $\widetilde P=\widetilde Q$ of this common restriction to all of $A$
(for instance, use one fixed off-support row).  The preceding paragraph gives
$(r,P)\sim(r,\widetilde P)$ and $(r,Q)\sim(r,\widetilde Q)$; since $\widetilde P=\widetilde Q$, common-environment
transitivity gives $(r,P)\sim(r,Q)$.

Part~(iii): We first prove $(r,P)\sim(r|_B,P|_B)$ for every channel $P:A\to\Delta(X)$.

\emph{Forward direction (\textup{(DP-act)} projection).}  Let $S:A\to B$ map all $a\notin B$ to some fixed $b_0\in B$.  The coarsened
prior is $rS=r|_B$ (since $r(a)=0$ for $a\notin B$), and the Bayesian pushforward $S^rP$ equals $P|_B$ on $B$.  By
\textup{(DP-act)},
\[
(r,P)\succeq(r|_B,P|_B).
\]

\emph{Reverse direction (\textup{(DP-act)} inclusion with completion).}  Let $I:B\hookrightarrow A$ be the inclusion.  Then
$(r|_B)I=r$.  The Bayesian pushforward $I^{r|_B}(P|_B)$ is pinned down on $B$ and arbitrary on $A\setminus B$.  Choose
the \textup{(DP-act)} completion to be any extension $\widetilde P:A\to\Delta(X)$ of $P|_B$.  Then
\[
(r|_B,P|_B)\succeq(r,\widetilde P).
\]
By Part~(ii), $(r,\widetilde P)\sim(r,P)$ because $\widetilde P$ and $P$ agree on $B=\operatorname{supp}(r)$.  Hence
$(r|_B,P|_B)\succeq(r,P)$.

Together the two directions give $(r,P)\sim(r|_B,P|_B)$.

\emph{Comparison equivalence.}  For any $P,Q$: $(r,P)\succeq(r,Q)$ iff (by the above indifferences $(r,P)\sim(r|_B,P|_B)$,
$(r,Q)\sim(r|_B,Q|_B)$, and Condition~\textup{(P1)} transitivity) $(r|_B,P|_B)\succeq(r|_B,Q|_B)$.

Part~(iv): If $\mu\in\mathcal M_r$, every posterior in the finite support of $\mu$ is supported on $B$; otherwise the
barycentre $r$ would assign positive mass to $A\setminus B$.  Hence $\pi_B(\mu)\in\mathcal M_{r|_B}$ is well defined.
By (iii), continuation comparisons at $r$ are isomorphic to those at $r|_B$.  The Lemma~\ref{pt:lem:postsep} representative
$F_{r|_B}^B$ on $\Delta(B)$ therefore induces a representative on $\mathcal M_r$ via $\pi_B$.
\end{proof}

\paragraph{Stage 3: product transport and the cross-prior bridge.}
\addcontentsline{toc}{subsubsection}{Stage 3: product transport and the cross-prior bridge}
An independent background is first shown to be ordinally inert.  Separate
affinity then yields the most general compatible product expression,
$x+y+\kappa xy$, after a coherent positive gauge.  Crucially,
Lemma~\ref{pt:lem:blockbridge} is proved before any coarse-revelation value is
used: product structure and face scales alone do not identify cross-prior
block comparisons.

\begin{lemma}[Derived background inertness]\label{pt:lem:background}
Assume Conditions~\textup{(P1)} and \textup{(P3)}--\textup{(P5)}.  Let $A_1,A_2$ be finite action sets and let
$q_1\in\Delta(A_1)$ and $q_2\in\Delta(A_2)$ have full support.  Then, for all channels $P_1,Q_1$ on $A_1$ and every
background channel $R_2$ on $A_2$,
\begin{equation}
(q_1\otimes q_2)^0
\succeq_{(P_1\otimes R_2)\sqcup(Q_1\otimes R_2)}
(q_1\otimes q_2)^1
\quad\Longleftrightarrow\quad
q_1^0\succeq_{P_1\sqcup Q_1}q_1^1,
\tag{BI}\label{pt:eq:backgroundinert}
\end{equation}
and symmetrically for comparisons in the second component with a common first-component background.  In particular the
left-hand comparison does not depend on the common background: for all backgrounds $R_2,S_2$ on $A_2$,
\[
(q_1\otimes q_2)^0
\succeq_{(P_1\otimes R_2)\sqcup(Q_1\otimes R_2)}
(q_1\otimes q_2)^1
\quad\Longleftrightarrow\quad
(q_1\otimes q_2)^0
\succeq_{(P_1\otimes S_2)\sqcup(Q_1\otimes S_2)}
(q_1\otimes q_2)^1.
\]
\end{lemma}

\begin{proof}
Throughout, for a prior $s$ on an action set $A$ and channels $V,W$ on $A$ write $V\succeq_sW$ for the block comparison
$s^0\succeq_{V\sqcup W}s^1$, and $(s,V)\sim(t,W)$ for the two-block indifference between the pairs.  By
Lemma~\ref{pt:lem:blockcoh}(a) these comparisons form a weak order for each fixed prior and chain transitively across
pairs with different priors and action sets.  By Lemma~\ref{pt:lem:plsuff}, at a full-support prior a channel may be replaced by any channel with the
same posterior law, in weak and in strict comparisons alike; we call this \emph{substitution}.

\emph{Preliminary reduction: a common foreground alphabet.}
Condition~\textup{(P5)} requires the two continuation channels used on a given branch to share a record alphabet.  If
$P_1:A_1\to\Delta(X)$ and $Q_1:A_1\to\Delta(X')$, replace them by their zero-padded versions on $X\sqcup X'$, as in
Step~2 of Lemma~\ref{pt:lem:publiccoin}.  Padding adds only null records, so it changes neither $\mu_{q_1,\cdot}$ nor
$\mu_{q_1\otimes q_2,\cdot\otimes R_2}$; both sides of~\eqref{pt:eq:backgroundinert} are therefore unaffected by
substitution at the full-support priors $q_1$ and $q_1\otimes q_2$.  Assume from now on that $P_1$ and $Q_1$ share the
record alphabet $X_1$.

For a channel $V:A_1\to\Delta(X_1)$ and a finite set $A_2$, write $\widetilde V$ for its lift to the product action set,
\[
\widetilde V(z\mid(a_1,a_2)):=V(z\mid a_1),
\]
and let $\widetilde R_2(x_2\mid(a_1,a_2)):=R_2(x_2\mid a_2)$ be the lift of the background.

\emph{Step 1: a product channel is a two-stage channel with background first stage.}
Take $\widetilde R_2$ as first stage on the action set $A_1\times A_2$ and the constant continuation profile
$\widetilde V$ in every branch, so that every branch uses the same continuation alphabet $X_1$ and the compound record
set $\bigsqcup_{x_2\in X_2}(\{x_2\}\times X_1)$ is $X_2\times X_1$.  By the definition of two-stage composition,
\[
\bigl(\widetilde R_2\triangleright\{\widetilde V\}_{x_2}\bigr)(x_2,z\mid(a_1,a_2))
=R_2(x_2\mid a_2)\,V(z\mid a_1)
=(V\otimes R_2)(z,x_2\mid (a_1,a_2)).
\]
The two channels therefore differ only by the bijection $(x_2,z)\mapsto(z,x_2)$ of record labels, which leaves every
record probability and every posterior unchanged.  Hence they induce the same posterior law at $q_1\otimes q_2$, and
substitution makes them interchangeable in every comparison at that prior.  Applying this to $V=P_1$ and $V=Q_1$, the
left-hand side of~\eqref{pt:eq:backgroundinert} is equivalent to
\begin{equation}
\widetilde R_2\triangleright\{\widetilde P_1\}
\;\succeq_{q_1\otimes q_2}\;
\widetilde R_2\triangleright\{\widetilde Q_1\},
\tag{$\ast$}\label{pt:eq:bgcompound}
\end{equation}
and likewise for the strict comparisons.

\emph{Step 2: every reached branch posterior is a product with first marginal $q_1$.}
Write $q:=q_1\otimes q_2$, $m(x_2):=m_{q,\widetilde R_2}(x_2)$ and, for $m(x_2)>0$, $r_{x_2}:=r_{x_2}^{\,q,\widetilde R_2}$.
Since the rows of $\widetilde R_2$ do not depend on $a_1$,
\[
m(x_2)=\sum_{a_2\in A_2}q_2(a_2)R_2(x_2\mid a_2)=m_{q_2,R_2}(x_2),
\]
and Bayes' rule gives
\[
r_{x_2}(a_1,a_2)
=\frac{q_1(a_1)q_2(a_2)R_2(x_2\mid a_2)}{m(x_2)}
=q_1(a_1)\,q_2^{x_2}(a_2),
\qquad
q_2^{x_2}:=r_{x_2}^{\,q_2,R_2}.
\]
Thus $r_{x_2}=q_1\otimes q_2^{x_2}$: the background signal may change beliefs about $a_2$, but in every reached branch
the two components stay independent and the first marginal is still $q_1$.  At least one branch has positive
probability.

\emph{Step 3: a lifted foreground channel is evaluated as on its own component.}
We claim that for every $s_2\in\Delta(A_2)$ and every channel $V:A_1\to\Delta(X_1)$,
\begin{equation}
(q_1\otimes s_2,\widetilde V)\sim(q_1,V).
\tag{$\dagger$}\label{pt:eq:liftneutral}
\end{equation}
First suppose $s_2$ has full support, so that $q_1\otimes s_2$ has full support on $A_1\times A_2$.  Let
$S:A_1\times A_2\to A_1$ be the coordinate projection, an onto deterministic map, and note that $\widetilde V$ is
$S$-measurable, since $S(a)=S(b)$ implies $\widetilde V(\cdot\mid a)=\widetilde V(\cdot\mid b)$.  Moreover
$(q_1\otimes s_2)S=q_1$ and, for every $a_1$,
\[
\bigl(S^{q_1\otimes s_2}\widetilde V\bigr)(z\mid a_1)
=\frac{\sum_{a_2\in A_2}q_1(a_1)s_2(a_2)V(z\mid a_1)}{q_1(a_1)}
=V(z\mid a_1),
\]
so the Bayesian pushforward of $\widetilde V$ under $S$ is $V$ itself.  Undetectable-distinctions neutrality
(Lemma~\ref{pt:lem:neutrality}) gives~\eqref{pt:eq:liftneutral}.

If $s_2$ is not full support, let $B_2:=\operatorname{supp}(s_2)$, so that $B:=A_1\times B_2$ is the support of
$q_1\otimes s_2$ and $B\subsetneq A_1\times A_2$.  The set $B$ \emph{is} the product action set $A_1\times B_2$, so no
relabelling is involved: the restricted prior is
$q_1\otimes(s_2|_{B_2})$ and the restricted channel $\widetilde V|_B$ is the lift of $V$ to $A_1\times B_2$.  Face
identification (Lemma~\ref{pt:lem:supprestrict}(iii)) gives
$(q_1\otimes s_2,\widetilde V)\sim(q_1\otimes s_2|_{B_2},\widetilde V|_B)$, and the full-support case just proved,
applied on the action set $A_1\times B_2$, gives $(q_1\otimes s_2|_{B_2},\widetilde V|_B)\sim(q_1,V)$.  Transitivity in a
common block environment yields~\eqref{pt:eq:liftneutral} in general.

Fix a branch with $m(x_2)>0$ and apply~\eqref{pt:eq:liftneutral} at $s_2=q_2^{x_2}$ to $V=P_1$ and to $V=Q_1$, using
$r_{x_2}=q_1\otimes q_2^{x_2}$ from Step~2.  Place the four pairs
$(r_{x_2},\widetilde P_1)$, $(q_1,P_1)$, $(r_{x_2},\widetilde Q_1)$, $(q_1,Q_1)$ as blocks of one finite block environment
(Condition~\textup{(P3)}, Lemma~\ref{pt:lem:blockcoh}).  The two indifferences and transitivity of the weak order in that
environment (Condition~\textup{(P1)}) give
\begin{equation}
\widetilde P_1\succeq_{r_{x_2}}\widetilde Q_1
\quad\Longleftrightarrow\quad
P_1\succeq_{q_1}Q_1,
\tag{$\ddagger$}\label{pt:eq:branchreduce}
\end{equation}
and, since the relation is complete, the same equivalence for the strict comparisons.  In particular the branch
comparison is the same in every reached branch and does not involve $R_2$.

\emph{Step 4: forward implication.}
Assume $P_1\succeq_{q_1}Q_1$.  By~\eqref{pt:eq:branchreduce}, the branch hypothesis
$r_{x_2}^0\succeq_{\widetilde P_1\sqcup\widetilde Q_1}r_{x_2}^1$ of Condition~\textup{(P5)} holds for every branch with
$m(x_2)>0$, with the common first stage $\widetilde R_2$ and the two constant continuation profiles.
Condition~\textup{(P5)} yields~\eqref{pt:eq:bgcompound}, hence the left-hand side of~\eqref{pt:eq:backgroundinert} by Step~1.

\emph{Step 5: converse, via the strictness clause.}
Assume the left-hand side of~\eqref{pt:eq:backgroundinert}, so that~\eqref{pt:eq:bgcompound} holds by Step~1, and suppose
$P_1\succeq_{q_1}Q_1$ fails.  By completeness and the reverse-block equivalence of Lemma~\ref{pt:lem:blockcoh},
$Q_1\succ_{q_1}P_1$, so by the strict form of~\eqref{pt:eq:branchreduce} every reached branch satisfies
$\widetilde Q_1\succ_{r_{x_2}}\widetilde P_1$.  At least one such branch exists and has positive probability, so the
strictness clause of Condition~\textup{(P5)}, applied with the two profiles exchanged, gives
\[
\widetilde R_2\triangleright\{\widetilde Q_1\}
\;\succ_{q_1\otimes q_2}\;
\widetilde R_2\triangleright\{\widetilde P_1\},
\]
contradicting~\eqref{pt:eq:bgcompound} because $\succeq_{q_1\otimes q_2}$ is a weak order
(Lemma~\ref{pt:lem:blockcoh}).  Hence $P_1\succeq_{q_1}Q_1$, proving~\eqref{pt:eq:backgroundinert}.

The second-component clause is the same argument with the roles of the two coordinates exchanged: the common
first-component background is used as the first stage and the second-component channels as the continuations.  Finally,
since both product comparisons in the last display of the lemma are equivalent to the single background-free comparison
$P_1\succeq_{q_1}Q_1$, they are equivalent to each other.
\end{proof}

The lemma shows that separability across independent components is not an independent behavioural commitment.  A product
channel is a two-stage channel whose first stage is the background experiment: it can teach the observer about the
background action, but it leaves the foreground component independent of it and its marginal unchanged, so in every
reached branch the continuation comparison is the foreground comparison itself.  Branchwise monotonicity then transmits
that single comparison to the compound, and its strictness clause transmits it back.  The derived statement is stronger
than mere invariance to changing the background: it identifies the product comparison with the comparison of the
foreground channels alone, which is the form used in Lemma~\ref{pt:lem:coherentnorm} below.

For finite posterior laws $\mu=\sum_i m_i\delta_{r_i}\in\mathsf X_A$ and
$\nu=\sum_j n_j\delta_{s_j}\in\mathsf X_B$, define
\[
\mu\otimes\nu
:=
\sum_{i,j}m_i n_j\,\delta_{r_i\otimes s_j}.
\]
Bayes' rule then gives the product factorisation
\begin{equation}
\mu_{q_1\otimes q_2,\,P_1\otimes P_2}
=
\mu_{q_1,P_1}\otimes\mu_{q_2,P_2}.
\tag{PL}\label{pt:eq:productlaw}
\end{equation}

\begin{lemma}[Coherent product quasi-additivity]\label{pt:lem:coherentnorm}
There is a choice of positive rescalings of the zero-normalised fibrewise representatives $F_q$, still denoted $F_q$,
and a scalar $\kappa\in\R$ such that, for all full-support priors $p\in\Delta(A)$ and $r\in\Delta(B)$ and all
$\mu\in\mathcal M_p$, $\nu\in\mathcal M_r$,
\begin{equation}
F_{p\otimes r}(\mu\otimes\nu)
=F_p(\mu)+F_r(\nu)+\kappa F_p(\mu)F_r(\nu).
\tag{QA}\label{pt:eq:QA}
\end{equation}
If one factor is a singleton, the same formula holds with its identically zero representative.  Moreover, for every
nondegenerate $r$ and every $\nu\in\mathcal M_r$,
\begin{equation}
1+\kappa F_r(\nu)>0.
\tag{POS}\label{pt:eq:QAslope}
\end{equation}
\end{lemma}

\begin{proof}
All representatives in this proof are zero-normalised: $F_q(\delta_q)=0$.  The proof first treats nondegenerate action
sets, meaning action sets with at least two elements.

\emph{Step 1: product-coordinate independence.}
Fix full-support priors $p$ and $r$ on nondegenerate finite action sets and set
\[
L_{p,r}(\mu,\nu):=F_{p\otimes r}(\mu\otimes\nu).
\]
This map is continuous and separately affine.  To identify a first-coordinate slice, let $P,Q$ implement
$\mu,\mu'\in\mathcal M_p$ and let $R$ implement $\nu\in\mathcal M_r$.  By the product factorisation
\eqref{pt:eq:productlaw}, $P\otimes R$ implements $\mu\otimes\nu$ and $Q\otimes R$ implements $\mu'\otimes\nu$ at the
full-support prior $p\otimes r$.  Hence Lemma~\ref{pt:lem:postsep}, applied on the fibre $\mathcal M_{p\otimes r}$ and on
the fibre $\mathcal M_p$, turns background inertness
\eqref{pt:eq:backgroundinert} of Lemma~\ref{pt:lem:background} into
\[
L_{p,r}(\mu,\nu)\ge L_{p,r}(\mu',\nu)
\quad\Longleftrightarrow\quad
F_p(\mu)\ge F_p(\mu').
\]
The symmetric clause of Lemma~\ref{pt:lem:background} identifies every second-coordinate slice with the $F_r$-order.  Local non-triviality, the
fibrewise representation, and the zero normalisation from Lemma~\ref{pt:lem:convnorm} make both slice orders nonconstant:
$F_p(\chi_p)>0$ and $F_r(\chi_r)>0$.

\emph{Step 2: the pair-specific bilinear form.}
Put $x(\mu):=F_p(\mu)$ and $y(\nu):=F_r(\nu)$.  For fixed $\nu$, affine-utility uniqueness gives
\[
L_{p,r}(\mu,\nu)=\alpha(\nu)x(\mu)+\gamma(\nu),
\qquad \alpha(\nu)>0.
\]
Taking $\mu=\delta_p$ shows $\gamma(\nu)=L_{p,r}(\delta_p,\nu)$.  This is a continuous affine representative of the
$F_r$-order and vanishes at $\delta_r$, so
\[
\gamma(\nu)=B_{p,r}y(\nu)
\]
for some $B_{p,r}>0$.  Likewise, the slice at $\nu=\delta_r$ gives
\[
L_{p,r}(\mu,\delta_r)=A_{p,r}x(\mu)
\]
for some $A_{p,r}>0$.

Let $\bar\mu:=\chi_p$ and $h_p:=F_p(\chi_p)>0$.  The function
$\nu\mapsto L_{p,r}(\bar\mu,\nu)$ is again a continuous affine representative of the $F_r$-order.  Hence there are
$D_{p,r}>0$ and a constant $E_{p,r}$ such that
\[
L_{p,r}(\bar\mu,\nu)=D_{p,r}y(\nu)+E_{p,r}.
\]
At $\nu=\delta_r$, $E_{p,r}=A_{p,r}h_p$.  Comparing this identity with
$L_{p,r}(\bar\mu,\nu)=\alpha(\nu)h_p+B_{p,r}y(\nu)$ gives
\[
\alpha(\nu)=A_{p,r}+C_{p,r}y(\nu),
\qquad
C_{p,r}:=\frac{D_{p,r}-B_{p,r}}{h_p}.
\]
Therefore
\begin{equation}
L_{p,r}(\mu,\nu)
=A_{p,r}x(\mu)+B_{p,r}y(\nu)+C_{p,r}x(\mu)y(\nu).
\tag{BIL}\label{pt:eq:pairbilinear}
\end{equation}
The coordinate-slice slopes are positive:
\[
A_{p,r}+C_{p,r}y(\nu)>0,
\qquad
B_{p,r}+C_{p,r}x(\mu)>0.
\]

\emph{Step 3: normalising the two linear coefficients.}
Let
\[
\alpha:((A\times B)\times C)\longrightarrow A\times(B\times C),
\qquad \alpha((a,b),c)=(a,(b,c)),
\]
be the canonical associating bijection.  It carries the left-grouped product
prior and posterior law to the right-grouped ones.  Exact cross-fibre
relabelling (Lemma~\ref{pt:lem:actionbase}) therefore makes the two evaluations
equal, with no unknown affine factor.  Expanding this common value under the
two groupings and comparing the coefficients of $x$, $y$, and $z$ gives
\begin{align}
A_{p\otimes r,s}A_{p,r}&=A_{p,r\otimes s}, \tag{C1}\label{pt:eq:Acoc1}\\
A_{p\otimes r,s}B_{p,r}&=B_{p,r\otimes s}A_{r,s}, \tag{C2}\label{pt:eq:Acoc2}\\
B_{p\otimes r,s}&=B_{p,r\otimes s}B_{r,s}. \tag{C3}\label{pt:eq:Acoc3}
\end{align}
Coefficient comparison is legitimate because the mixtures $(1-t)\delta_p+t\chi_p$ make $x$ range over a nondegenerate
interval, and similarly in the other coordinates.

The coordinate swap is another bijection, so exact cross-fibre relabelling gives
\[
F_{r\otimes p}(\nu\otimes\mu)=F_{p\otimes r}(\mu\otimes\nu).
\]
Comparing coefficients in~\eqref{pt:eq:pairbilinear},
\[
B_{r,p}=A_{p,r},
\qquad
A_{r,p}=B_{p,r},
\qquad
C_{r,p}=C_{p,r}.
\]
Thus, with $\rho(p,r):=B_{p,r}/A_{p,r}$,
\begin{equation}
\rho(r,p)=\rho(p,r)^{-1}.
\tag{SW}\label{pt:eq:rhoswap}
\end{equation}

Divide~\eqref{pt:eq:Acoc2} by~\eqref{pt:eq:Acoc1}:
\[
\rho(p,r)=\rho(p,r\otimes s)A_{r,s}.
\]
Fix a nondegenerate reference prior $q_0$ on a finite action set $A_0$ and put $g(r):=\rho(q_0,r)$.  Then
\begin{equation}
A_{r,s}=\frac{g(r)}{g(r\otimes s)}.
\tag{G}\label{pt:eq:gaugeA}
\end{equation}
The coefficients $A_{p,r},B_{p,r},C_{p,r}$ are unchanged when either prior is
transported by a bijection.  Indeed, transport the product law by the product
bijection, use Lemma~\ref{pt:lem:actionbase} on the two input fibres and the
product fibre, and compare the unique coefficients in
\eqref{pt:eq:pairbilinear}.  Hence $\rho$, and therefore $g$, is invariant
under relabelling.  In particular it takes the same value on the two
associations of a product prior.

Rescale every representative by $F_p^*:=g(p)F_p$.  Under this positive rescaling,
\[
A^*_{p,r}=A_{p,r}\frac{g(p\otimes r)}{g(p)},
\quad
B^*_{p,r}=B_{p,r}\frac{g(p\otimes r)}{g(r)},
\quad
C^*_{p,r}=C_{p,r}\frac{g(p\otimes r)}{g(p)g(r)}.
\]
Because $g$ is relabelling-invariant, these rescaled representatives retain
the exact relabelling equality of Lemma~\ref{pt:lem:actionbase}.  Thus the
associator and swap comparisons used below still carry no scalar defect.
Equation~\eqref{pt:eq:gaugeA} gives $A^*_{p,r}=1$ for every pair.  The
transformed swap identity gives $B^*_{p,r}=A^*_{r,p}=1$.  Thus, after
the positive gauge rescaling, $A_{p,r}=B_{p,r}=1$ for every nondegenerate pair
(we now drop the stars).  Positivity of the slice
slopes is preserved because the rescaling factors are positive.

For a boundary prior $r\in\Delta(A)$ with $B=\operatorname{supp}(r)$, redefine
the post-gauge boundary representative by exact support transport,
\[
F_r^A(\mu):=F_{r|_B}^B\bigl(\pi_B(\mu)\bigr),
\]
where the representative on the right is the newly gauged full-support one;
singleton support faces remain identically zero.  Thus the support identity of
Lemma~\ref{pt:lem:supprestrict} remains definitional after the gauge change.

\emph{Step 4: the interaction coefficient is universal.}
Write the remaining coefficient as $\kappa_{p,r}$.  Comparing the $xy$, $xz$, $yz$, and $xyz$ coefficients in the two
associative expansions gives
\begin{align}
\kappa_{p,r}&=\kappa_{p,r\otimes s}, \tag{K1}\label{pt:eq:kappa1}\\
\kappa_{p\otimes r,s}&=\kappa_{p,r\otimes s}, \tag{K2}\label{pt:eq:kappa2}\\
\kappa_{p\otimes r,s}&=\kappa_{r,s}, \tag{K3}\label{pt:eq:kappa3}\\
\kappa_{p\otimes r,s}\kappa_{p,r}&=\kappa_{p,r\otimes s}\kappa_{r,s}. \tag{K4}\label{pt:eq:kappa4}
\end{align}
The last equation follows from the first three, but it records the full coefficient comparison.  After $A=B=1$, exact
invariance under the coordinate swap and comparison of the bilinear coefficient give
\begin{equation}
\kappa_{p,r}=\kappa_{r,p}.
\tag{KS}\label{pt:eq:kappasym}
\end{equation}
Using~\eqref{pt:eq:kappa1}--\eqref{pt:eq:kappa3} with the reference prior $q_0$,
\[
\kappa_{p,r}
=\kappa_{p,r\otimes q_0}
=\kappa_{p\otimes r,q_0}
=\kappa_{r,q_0}.
\]
Hence $\kappa_{p,r}$ is independent of its first argument.  By symmetry it is also independent of its second argument.
Denote the common value by $\kappa$.  This proves~\eqref{pt:eq:QA} for nondegenerate factors.

The positive slice-slope condition now reads
\[
1+\kappa F_r(\nu)>0,
\]
which is~\eqref{pt:eq:QAslope}.

\emph{Step 5: singleton factors.}
Let $q_*=\delta_*$ be the prior on a singleton.  Its zero-normalised representative is identically zero.  If $p$ is
nondegenerate, the canonical bijection $A\to A\times\{*\}$ and exact
relabeling give
\[
F_{p\otimes q_*}(\mu\otimes\delta_{q_*})=F_p(\mu).
\]
The product gauge preserves this equality because $p\otimes q_*$ is a
relabelling of $p$.  Since $F_{q_*}=0$, equation~\eqref{pt:eq:QA} follows
immediately.  The reversed order is identical, and the two-singleton case is
$0=0$.
\end{proof}

\begin{corollary}[Exact relabelling after the product gauge]\label{pt:cor:permutationinvariance}
For every bijection $\pi:A\to A'$ and every full-support $q\in\Delta(A)$,
\[
F_{q\pi}(\mu\pi)=F_q(\mu)
\qquad(\mu\in\mathcal M_q).
\]
\end{corollary}

\begin{proof}
Before the product gauge this is Lemma~\ref{pt:lem:actionbase}.  The proof of
Lemma~\ref{pt:lem:coherentnorm} shows that the gauge factor $g(q)$ is
unchanged by every bijection, so multiplication by that factor preserves the
equality.  The singleton case is immediate because both representatives
vanish.
\end{proof}

\begin{lemma}[Cross-prior block representation]\label{pt:lem:blockbridge}
For arbitrary finite priors $q\in\Delta(A)$ and $r\in\Delta(B)$ and channels
$P:A\to\Delta(X)$ and $Q:B\to\Delta(Y)$,
\begin{equation}
q^0\succeq_{P\sqcup Q}r^1
\quad\Longleftrightarrow\quad
F_q(\mu_{q,P})\ge F_r(\mu_{r,Q}).
\tag{BC}\label{pt:eq:blockcomp}
\end{equation}
Boundary representatives are interpreted by Lemma~\ref{pt:lem:supprestrict}.
\end{lemma}

\begin{proof}
First suppose $q$ and $r$ have full support, allowing singleton action sets.  At prior $q\otimes r$, quasi-additivity and
zero normalisation give
\[
F_{q\otimes r}(\mu_{q,P}\otimes\delta_r)=F_q(\mu_{q,P}),
\qquad
F_{q\otimes r}(\delta_q\otimes\mu_{r,Q})=F_r(\mu_{r,Q}).
\]
These values are in the same fibre $\mathcal M_{q\otimes r}$, so Lemma~\ref{pt:lem:postsep} gives
\[
(q\otimes r)^0\succeq_{(P\otimes U_B)\sqcup(U_A\otimes Q)}(q\otimes r)^1
\quad\Longleftrightarrow\quad
F_q(\mu_{q,P})\ge F_r(\mu_{r,Q}).
\]

It remains to transfer the product-lifted comparison back to the original two blocks.  For the $A$-side pair,
Clause \textup{(DP-act)} of Condition~\textup{(P4)} applied to the projection $\pi_A:A\times B\to A$ gives
\[
(q\otimes r,P\otimes U_B)\succeq(q,P)
\]
in block comparisons.  Conversely, let $T:A\to\Delta(A\times B)$ be
\[
T((a,b)\mid a')=\1_{\{a=a'\}}r(b).
\]
Then $qT=q\otimes r$, and the Bayesian pushforward $T^qP$ has row $P(\cdot\mid a)$ at every $(a,b)$.  Thus
$T^qP$ is $P\otimes U_B$ after the canonical one-record identification of the $B$-experiment, and \textup{(DP-act)}
gives the reverse comparison.  Equivalently, if the harmless record identification is not made, Lemma~\ref{pt:lem:plsuff}
identifies $T^qP$ and $P\otimes U_B$ because they induce the same posterior law at $q\otimes r$.  Hence
$(q\otimes r,P\otimes U_B)$ and $(q,P)$ are indifferent in every relevant block comparison.  The same
projection/embedding argument identifies $(q\otimes r,U_A\otimes Q)$ with $(r,Q)$.

Place the four pairs $(q,P)$, $(q\otimes r,P\otimes U_B)$, $(r,Q)$, and $(q\otimes r,U_A\otimes Q)$ as blocks of one
finite block environment.  Lemma~\ref{pt:lem:blockcoh} and transitivity transfer the two pairwise indifferences to the
two-block comparison, proving~\eqref{pt:eq:blockcomp} for full-support priors.

For arbitrary priors, let $A_+:=\operatorname{supp}(q)$ and $B_+:=\operatorname{supp}(r)$.  Lemma~\ref{pt:lem:supprestrict}
identifies $(q,P)$ with $(q|_{A_+},P|_{A_+})$, identifies $(r,Q)$ with $(r|_{B_+},Q|_{B_+})$, and identifies the two
values $F_q(\mu_{q,P})$ and $F_r(\mu_{r,Q})$ with their support-face values.  Applying the full-support bridge on
$A_+$ and $B_+$ and pulling back by Lemma~\ref{pt:lem:supprestrict} proves the claim.  Singleton supports use the same
zero-representative convention as Lemma~\ref{pt:lem:coherentnorm}.
\end{proof}

\paragraph{Stage 4: branch coefficients, cocycle, and support faces.}
\addcontentsline{toc}{subsubsection}{Stage 4: branch coefficients, cocycle, and support faces}
Insertion into a reached branch is affine.  Its linear part is evaluated on
finite atomic signed posterior laws; the spanning argument shows that the
resulting positive coefficient depends only on the prior and reached
posterior.  Path independence gives the cocycle and the normalized chain
rule.  The subsequent cardinal correction aligns auxiliary chain scales on
nested faces; it does not add a behavioural assumption or change the ordinal
representative.

\begin{lemma}[Branch aggregation]\label{pt:lem:branchagg}
Let $P_1:A\to\Delta(X_1)$ be a first-stage channel at full-support $q$, with branch probabilities $m(x)$ and posteriors
$r_x$.  There are positive branch coefficients $\beta(q,r_x)$, depending only on $q$ and the reached posterior $r_x$,
such that for every continuation profile $\{Q^x\}_{x:m(x)>0}$,
\[
F_q(\mu_{q,P_1\triangleright\{Q^x\}})
=
F_q(\mu_{q,P_1})+
\sum_{x:m(x)>0}m(x)\,\beta(q,r_x)\,
F_{r_x}(\mu_{r_x,Q^x}).
\]
If $r_x$ has full support, $\beta(q,r_x)$ is the full-support tangent coefficient.  If
$1<|\operatorname{supp}(r_x)|<|A|$, it is the full-to-face coefficient.  If $|\operatorname{supp}(r_x)|=1$, the value
$F_{r_x}$ is identically zero and $\beta(q,r_x)$ is an arbitrary positive convention.
\end{lemma}

\begin{proof}
Fix the first-stage channel $P_1$ and its branch structure $(m(x),r_x)_{x:m(x)>0}$.  For any profile of continuation
channels $\{Q^x\}$, the compound posterior law at prior $q$ is
\[
\mu_{q,P_1\triangleright\{Q^x\}}
=
\sum_{x:m(x)>0}m(x)\,\mu_{r_x,Q^x}.
\]

\emph{Setup: linear part of affine functional.}  Since $F_q$ is affine on $\mathcal M_q$ (Lemma~\ref{pt:lem:postsep}),
it extends uniquely to an affine functional on the affine hull of $\mathcal M_q$.  Write $L_q$ for the linear part:
for any $\mu,\mu'\in\mathcal M_q$,
\[
F_q(\mu)-F_q(\mu')=L_q(\mu-\mu').
\]

\emph{Branch-affine structure.}  Fix a positive-probability branch $\bar x$ and fix continuation channels
$H^x:A\to\Delta(Z^x)$ in all other branches $x\ne\bar x$.  Let $r:=r_{\bar x}$ and $m:=m(\bar x)$.  If
$\operatorname{supp}(r)$ is a singleton, then $\mathcal M_r=\{\delta_r\}$ and $F_r(\delta_r)=0$ by the singleton
zero-normalisation convention; such branches contribute zero and their coefficient convention is recorded in
Case~1 below.  For the branch-affine argument in this paragraph, assume $|\operatorname{supp}(r)|\ge2$.  If $r$ is a
boundary posterior, $F_r$ and continuation comparisons at $r$ are understood through Lemma~\ref{pt:lem:supprestrict}.

Let
\[
\rho_H:=\sum_{x\ne\bar x}m(x)\,\mu_{r_x,H^x}.
\]
For any $\nu\in\mathcal M_r$, the aggregate posterior law produced by using branch law $\nu$ in branch $\bar x$ and
the fixed continuations elsewhere is
\[
T_H(\nu):=m\nu+\rho_H\in\mathcal M_q.
\]
Define $g:\mathcal M_r\to\R$ by
\[
g(\nu):=F_q(T_H(\nu)).
\]
Thus $g$ depends only on $\nu$, and $g=F_q\circ T_H$ is continuous and affine.

We now show that $g$ and $F_r$ represent the same weak order on $\mathcal M_r$.  For $\nu,\nu'\in\mathcal M_r$, choose
finite continuation channels implementing them at prior $r$; when $r$ is boundary these are implemented on
$B=\operatorname{supp}(r)$ and extended arbitrarily to $A$, as licensed by Lemma~\ref{pt:lem:supprestrict}.  Padding record
alphabets if needed, Condition~\textup{(P5)} compares the two continuation profiles that differ only in branch $\bar x$.
If $F_r(\nu)>F_r(\nu')$, Condition~\textup{(P5)} gives $g(\nu)>g(\nu')$.  Conversely, if $g(\nu)\ge g(\nu')$ but
$F_r(\nu)<F_r(\nu')$, applying the strict part of Condition~\textup{(P5)} to the swapped branch comparison gives
$g(\nu')>g(\nu)$, a contradiction.  Hence $g(\nu)\ge g(\nu')$ iff $F_r(\nu)\ge F_r(\nu')$.

Since $|\operatorname{supp}(r)|\ge2$, local non-triviality on the support face and Lemmas~\ref{pt:lem:supprestrict} and~\ref{pt:lem:convnorm} imply $F_r$ is nonconstant.  The continuous affine functions $g$ and $F_r$ are therefore nonconstant
affine representatives of the same weak order on $\mathcal M_r$.  Herstein--Milnor uniqueness gives
\[
g(\nu)=\alpha\,F_r(\nu)+\gamma
\]
for some $\alpha>0$ and $\gamma\in\R$.

\emph{The coefficient $\alpha$ is independent of the fixed continuations.}  Write $g^{(1)}$ and $g^{(2)}$ for the
functionals arising from two different choices $\{H^x\}$ and $\{\widetilde H^x\}$ in the other branches.  Both
represent the same order as $F_r$, so $g^{(1)}=\alpha_1 F_r+\gamma_1$ and
$g^{(2)}=\alpha_2 F_r+\gamma_2$ with $\alpha_1,\alpha_2>0$.  To show $\alpha_1=\alpha_2$, consider the
difference $g^{(1)}-g^{(2)}=(\alpha_1-\alpha_2)F_r+(\gamma_1-\gamma_2)$.  For any
$\nu\in\mathcal M_r$,
\[
g^{(1)}(\nu)-g^{(2)}(\nu)
=
F_q\!\left(m(\bar x)\nu+\textstyle\sum_{x\ne\bar x}m(x)\mu_{r_x,H^x}\right)
-
F_q\!\left(m(\bar x)\nu+\textstyle\sum_{x\ne\bar x}m(x)\mu_{r_x,\widetilde H^x}\right).
\]
Both arguments are posterior laws in $\mathcal M_q$, differing only in the fixed non-$\bar x$ components.  Since $L_q$
has already been defined on the affine-hull difference space,
\[
g^{(1)}(\nu)-g^{(2)}(\nu)=L_q(\rho_1-\rho_2),
\]
where $\rho_i$ is the fixed non-$\bar x$ background law for $g^{(i)}$.  This difference is independent of $\nu$.  Since
$F_r$ is nonconstant in the present nondegenerate case, $\alpha_1=\alpha_2$.

\emph{Path-independence of the branch coefficient via tangent-space decomposition.}
We now prove that $\alpha/m(\bar x)$ depends only on the pair $(q,r_{\bar x})$, not on the first-stage experiment.
This is the key step that converts ordinal branchwise monotonicity into a cardinal aggregation formula.

\emph{Tangent space at a posterior.}  For a posterior $r\in\Delta(A)$ with support $B:=\operatorname{supp}(r)$, define
the \emph{tangent space}:
\[
V_r:=\operatorname{span}\{\nu-\nu':\nu,\nu'\in\mathcal M_r\}.
\]
Elements of $V_r$ are signed measures with total mass zero and barycentre zero.

\emph{Tangent-space identification.}  For any $r$ with full support, we claim
\[
V_r=W:=\{\eta\in\operatorname{span}\Delta_f(\Delta(A)):\eta(\Delta(A))=0,\ \textstyle\int p\,d\eta(p)=0\},
\]
the space of finite-support signed measures on $\Delta(A)$ with zero total mass and zero barycentre.
Indeed, if $\nu,\nu'\in\mathcal M_r$, then $\nu-\nu'$ has total mass $1-1=0$ and barycentre $r-r=0$, so $V_r\subset W$.
Conversely, given any $\eta\in W$: if $\eta=0$ then $\eta\in V_r$ trivially.  Otherwise, since $\eta$ has finite support,
	write its Jordan decomposition $\eta=\eta^+-\eta^-$ where $\eta^+,\eta^-\ge 0$ are finite-support positive measures
	with equal total mass $c>0$.  Normalise $\zeta^+:=\eta^+/c$ and $\zeta^-:=\eta^-/c$; both are finite-support probability
	measures on $\Delta(A)$ with the same barycentre
	$b:=\int p\,d\zeta^+(p)=\int p\,d\zeta^-(p)$.  Since $b\in\Delta(A)$, the set $\{a:b(a)>0\}$ is nonempty.  Because $r$
	has full support,
	\[
	\theta:=\min\Bigl\{1,\min_{a:b(a)>0}\frac{r(a)}{b(a)}\Bigr\}>0.
	\]
	Choose $0<\lambda<\theta$ and set
	\[
	\tau:=\frac{r-\lambda b}{1-\lambda}.
	\]
	Then $1-\lambda>0$, $\sum_a\tau(a)=1$, and $\tau(a)\ge0$ for every $a$: if $b(a)>0$ this follows from
	$\lambda<r(a)/b(a)$, while if $b(a)=0$ the numerator is $r(a)>0$.  Hence $\tau\in\Delta(A)$.  Define
	\[
	\nu:=\lambda\zeta^++(1-\lambda)\delta_\tau,
	\qquad
	\nu':=\lambda\zeta^-+(1-\lambda)\delta_\tau.
	\]
	Both are finite-support probability measures on $\Delta(A)$ with barycentre $\lambda b+(1-\lambda)\tau=r$, so
	$\nu,\nu'\in\mathcal M_r$.  Then $\nu-\nu'=\lambda(\zeta^+-\zeta^-)=(\lambda/c)\eta$, hence
	$\eta=(c/\lambda)(\nu-\nu')\in V_r$.  Thus $W\subset V_r$ and $V_r=W$.
Since $W$ depends only on the full-support property (not on the specific prior), $V_q=V_r=V_s=W$ for all full-support priors.

\emph{Case 1: Degenerate branch posterior ($|B|=1$).}  If $r=\delta_a$ for some action $a$, then
	$\mathcal M_{\delta_a}=\{\delta_{\delta_a}\}$ is a singleton: the only Bayes-plausible law with barycentre $\delta_a$
	is the point mass at $\delta_a$ itself.  Hence $V_r=\{0\}$, and every continuation experiment at prior $\delta_a$ gives
	the same posterior law $\delta_{\delta_a}$.  We set $F_{\delta_a}(\delta_{\delta_a})=0$.  For such branches, the term
	$m(x)\beta(q,r_x)F_{r_x}(\mu_{r_x,Q^x})$ is zero under any convention for $\beta(q,r_x)$.

For the aggregation formula below, assign any positive value to $\beta(q,\delta_a)$; the choice is not behaviourally
identified because the corresponding continuation value is always zero.


\emph{Case 2: Non-full-support branch posterior ($1<|B|<|A|$).}  If $r$ has support $B\subsetneq A$ with $|B|\ge 2$,
apply Lemma~\ref{pt:lem:supprestrict}.  Continuation comparisons at prior $r$ reduce to those at $r|_B\in\Delta(B)$,
with representative $F_r:=F_{r|_B}^B\circ\pi_B$.

For any finite action set $C$, write
\[
W_C:=\{\sigma\in\operatorname{span}\Delta_f(\Delta(C)):\sigma(\Delta(C))=0,\ \textstyle\int_{\Delta(C)}p\,d\sigma(p)=0\}.
\]
Let $\jmath_B:\Delta(B)\to\Delta(A)$ be the face inclusion, extending each posterior by zero outside $B$, and let
$i_B:=(\jmath_B)_\#$ act on finite signed posterior laws.  By Lemma~\ref{pt:lem:supprestrict}(iv),
$F_r^A=F_{r|_B}^B\circ\pi_B$, so on face directions
\[
L_r^A(i_B\sigma)=L_{r|_B}^B(\sigma)\qquad(\sigma\in W_B).
\]
Since $|B|\ge2$ and $r|_B$ is full support on $B$, local non-triviality and Lemma~\ref{pt:lem:convnorm} give
\[
L_{r|_B}^B(\chi_{r|_B}-\delta_{r|_B})>0,
\]
so $L_{r|_B}^B\ne0$.

\emph{Full-to-face scalar.}  Fix any full-support $q\in\Delta(A)$.  The boundary posterior $r$ is reachable from $q$:
choose
\[
0<\varepsilon<\min\Bigl\{1,\min_{b\in B}\frac{q(b)}{r(b)}\Bigr\}
\]
and set $P_1(1\mid a)=\varepsilon r(a)/q(a)$, with $r(a)=0$ for $a\notin B$.  Then branch $1$ has positive probability
$\varepsilon$ and posterior $r$.  For any nonzero $\sigma\in W_B$, the tangent-space identification on the face $B$ gives
$\sigma=t(\nu-\nu')$ for some $t>0$ and $\nu,\nu'\in\mathcal M_{r|_B}^B$.  Implement $\nu,\nu'$ by continuation channels
on $B$ and extend them arbitrarily to $A$.  Lemma~\ref{pt:lem:supprestrict} identifies their comparison at the boundary prior
$r$ with their comparison at $r|_B$.

Condition~\textup{(P5)} applied to the branch reaching $r$, with all other branches fixed, gives the positive sign direction
for feasible differences.  The negative direction follows by swapping the two continuation laws, and equality follows by
applying the weak part of Condition~\textup{(P5)} in both directions.  Extending by positive scalar multiples from feasible
differences to all of $W_B$, $L_q^A\circ i_B$ and $L_{r|_B}^B$ have the same sign partition on $W_B$.  Since
$L_{r|_B}^B\ne0$, both are nonzero and there is a unique $\beta_A(q,r)>0$ such that
\[
L_q^A(i_B\sigma)=\beta_A(q,r)L_{r|_B}^B(\sigma)\qquad(\sigma\in W_B).
\]
This is the ambient full-to-face coefficient used for boundary branches.

\emph{Case 3: Full-support branch posterior ($B=A$).}  This is the main case.  Suppose $r$ has full support on $A$.

\emph{Claim 1: aggregate change from branch variation.}  If a first-stage branch reaches posterior $r$ with probability
$m$, and the branch continuation law changes from $\nu$ to $\nu'$ (both in $\mathcal M_r$), then the aggregate posterior
law changes by $m(\nu-\nu')\in m\cdot V_r$.  This difference depends only on $(m,\nu,\nu',r)$---not on which experiment
produced the branch, nor on the other branches.

\emph{Proof of Claim 1.}  The aggregate changes from $m\nu+\rho$ to $m\nu'+\rho$, where $\rho$ collects the other
branches.  The difference is $m(\nu-\nu')$.

\emph{Claim 2: $L_q|_{V_r}$ is a positive scalar multiple of $L_r|_{V_r}$.}

\emph{Proof of Claim 2.}  We first show that every direction in $V_r$ is proportional to a feasible difference.
By the Tangent-space identification above, any nonzero $\sigma\in V_r$ can be written as
$\sigma=t(\nu-\nu')$ for some $t>0$ and $\nu,\nu'\in\mathcal M_r$.  Hence every nonzero $\sigma\in V_r$ is a positive
scalar multiple of some feasible difference.

By Condition~\textup{(P5)} (branchwise monotonicity), if $\nu\succ_r\nu'$ (i.e.\ $F_r(\nu)>F_r(\nu')$, equivalently
$L_r(\nu-\nu')>0$), then the aggregate comparison ranks the $\nu$-continuation strictly higher.  Since the aggregate
value difference is $L_q(m(\nu-\nu'))=mL_q(\nu-\nu')$ with $m>0$, we have $L_q(\nu-\nu')>0$ whenever $L_r(\nu-\nu')>0$.

\emph{Reverse and equality directions.}  If $L_r(\nu-\nu')<0$, completeness of $\succeq_r$ gives $\nu'\succ_r\nu$, and
Condition~\textup{(P5)} applied to the reversed comparison gives $L_q(\nu'-\nu)>0$, i.e.\ $L_q(\nu-\nu')<0$.  If $L_r(\nu-\nu')=0$, then
$\nu\sim_r\nu'$, so both $\nu\succeq_r\nu'$ and $\nu'\succeq_r\nu$; applying Condition~\textup{(P5)} weakly in both directions gives
$L_q(\nu-\nu')\ge 0$ and $L_q(\nu-\nu')\le 0$, hence $L_q(\nu-\nu')=0$.

Since every nonzero $\sigma\in V_r$ is proportional to some feasible difference, the sign agreement extends from
feasible differences to all of $V_r$: for any $\sigma\in V_r$, $L_q(\sigma)>0$ iff $L_r(\sigma)>0$, $L_q(\sigma)<0$
iff $L_r(\sigma)<0$, and $L_q(\sigma)=0$ iff $L_r(\sigma)=0$.

Two nonzero linear functionals on a vector space that have the same positive half-space must be positive scalar
multiples of each other.  (If $L_r\ne 0$, then $\ker L_q=\ker L_r$ and they lie in the same half-space, so $L_q=cL_r$
for some $c>0$.)  Since $r$ has full support and local non-triviality gives $F_r(\chi_r)>F_r(\delta_r)=0$, we have
$L_r\ne 0$.  Hence there exists $\beta(q,r)>0$ such that
\[
L_q(\sigma)=\beta(q,r)\,L_r(\sigma)
\qquad\text{for all }\sigma\in V_r.
\]

\emph{Path-independence.}  Consider any first-stage experiment reaching posterior $r$ in some branch with probability
$m$.  Changing the branch continuation from $\nu$ to $\nu'$ changes the aggregate value by
\[
L_q(m(\nu-\nu'))=m\,\beta(q,r)\,L_r(\nu-\nu')=m\,\beta(q,r)\,(F_r(\nu)-F_r(\nu')).
\]
The coefficient $m\,\beta(q,r)$ depends only on $(q,r,m)$.  In particular, $\beta(q,r)$ is independent of which
experiment reached $r$ and independent of the other branches.

Comparing with the earlier notation $g(\nu)=\alpha F_r(\nu)+\gamma$: the value difference is
$g(\nu)-g(\nu')=\alpha(F_r(\nu)-F_r(\nu'))$, so $\alpha=m\,\beta(q,r)$ and hence $\alpha/m=\beta(q,r)$ universally.
The coefficient $\beta(q,r_{\bar x})$ is defined as the scalar from $L_q|_{V_{r_{\bar x}}}=\beta(q,r_{\bar x})L_{r_{\bar x}}|_{V_{r_{\bar x}}}$;
it depends only on $q$ and $r_{\bar x}$, not on the experiment that produced the branch.

\emph{Aggregation formula.}  Let
\[
\mu^0:=\sum_{x:m(x)>0}m(x)\delta_{r_x}=\mu_{q,P_1}.
\]
For each positive-probability branch let $B_x:=\operatorname{supp}(r_x)$.  If $B_x=A$, set
$v_x:=\mu_{r_x,Q^x}-\delta_{r_x}\in W_A$.  If $1<|B_x|<|A|$, set
\[
\bar v_x:=\pi_{B_x}(\mu_{r_x,Q^x})-\delta_{r_x|_{B_x}}\in W_{B_x},
\qquad
v_x:=i_{B_x}\bar v_x\in W_A.
\]
If $|B_x|=1$, set $v_x:=0$.  Then
\[
\mu_{q,P_1\triangleright\{Q^x\}}=\mu^0+\sum_{x:m(x)>0}m(x)v_x.
\]
Since $V_q=W_A$, affinity of $F_q$ gives
\[
F_q(\mu_{q,P_1\triangleright\{Q^x\}})-F_q(\mu^0)
=\sum_{x:m(x)>0}m(x)L_q^A(v_x).
\]
For a full-support branch $B_x=A$, Case~3 gives
\[
L_q^A(v_x)=\beta(q,r_x)L_{r_x}^A(v_x)
=\beta(q,r_x)F_{r_x}(\mu_{r_x,Q^x}).
\]
For a nondegenerate boundary branch $1<|B_x|<|A|$, Case~2 gives
\[
L_q^A(v_x)=L_q^A(i_{B_x}\bar v_x)
=\beta_A(q,r_x)L_{r_x|_{B_x}}^{B_x}(\bar v_x)
=\beta_A(q,r_x)F_{r_x}(\mu_{r_x,Q^x}).
\]
The last equalities use zero normalisation: Lemma~\ref{pt:lem:convnorm} in full support, Lemma~\ref{pt:lem:supprestrict} plus
Lemma~\ref{pt:lem:convnorm} on the support face in the boundary case, and the singleton convention when $|B_x|=1$.  Therefore
\[
F_q(\mu_{q,P_1\triangleright\{Q^x\}})
=
F_q(\mu_{q,P_1})+\sum_{x:m(x)>0}m(x)\,\beta(q,r_x)\,F_{r_x}(\mu_{r_x,Q^x}),
\]
where $\beta(q,r_x)$ denotes the Case~3 coefficient when $B_x=A$, the Case~2 full-to-face coefficient when
$1<|B_x|<|A|$, and any positive singleton convention when $|B_x|=1$.
\end{proof}

\begin{lemma}[Branch-coefficient cocycle and normalised chain rule]\label{pt:lem:chain}
Let the full-support and full-to-face coefficients be those constructed in Lemma~\ref{pt:lem:branchagg}, applied in the
relevant ambient simplex or intrinsically on the relevant support face.  If $q$ has full support and $r,s$ are
nondegenerate priors with
\[
\operatorname{supp}(s)\subseteq\operatorname{supp}(r)\subseteq\operatorname{supp}(q),
\]
then the coefficients satisfy the cocycle identity
\[
\beta(q,s)=\beta(q,r)\beta(r,s).
\]
Hence $\beta(q,r)=a_q/a_r$ for positive scales $a_q$, with singleton scales fixed by convention.  The rescaled values
\[
\widehat F_q(\mu):=\frac{F_q(\mu)}{a_q}
\]
satisfy the exact chain rule: for every first-stage channel $P_1$ at full-support $q$ and every continuation profile
$\{Q^x\}_{x:m(x)>0}$,
\[
\widehat F_q(\mu_{q,P_1\triangleright\{Q^x\}})
=
\widehat F_q(\mu_{q,P_1})
\;+\;
\sum_{x:m(x)>0}m(x)\,\widehat F_{r_x}(\mu_{r_x,Q^x}).
\]
\end{lemma}

\begin{proof}

\emph{Full-support cocycle.}  First consider reachable full-support $q,r,s\in\operatorname{int}\Delta(A)$ with $|A|\ge2$.
For such priors, all three tangent spaces coincide.
By the tangent-space identification in Lemma~\ref{pt:lem:branchagg}, $V_q=V_r=V_s=W$ for all full-support priors.
The full-support coefficient identities constructed there share a common domain $W$:
\[
L_q|_W=\beta(q,r)L_r|_W,
\qquad
L_r|_W=\beta(r,s)L_s|_W.
\]
Substituting the second into the first:
\[
L_q|_W=\beta(q,r)\beta(r,s)L_s|_W.
\]
Comparing with $L_q|_W=\beta(q,s)L_s|_W$ and using $L_s\ne 0$:
\[
\beta(q,s)=\beta(q,r)\beta(r,s).
\]

\emph{Reachability.}  The above uses the full-support coefficient construction for the pairs $(q,r)$, $(r,s)$, and
$(q,s)$.  Each invocation requires the pair to be reachable: the second prior must be a possible posterior of some
experiment at the first prior.
For any full-support $q,r$, construct a two-record experiment $P_1$ with $P_1(1\mid a)=\varepsilon r(a)/q(a)$ for
$0<\varepsilon<\min\{1,\min_a q(a)/r(a)\}$.  Record~1 has probability $\varepsilon$ and posterior~$r$; record~0 has posterior
$(q-\varepsilon r)/(1-\varepsilon)$.  Hence any full-support $r$ is reachable from any full-support $q$, so the construction
applies to all full-support pairs.  The cocycle therefore holds for all full-support triples.

Fix a full-support basepoint $q_0$, set $a_{q_0}=1$, and for full-support $q$ define
\[
a_q:=\frac{1}{\beta(q_0,q)}.
\]
By the full-support cocycle, this is equivalently $a_q=\beta(q,q_0)$, and
\[
\beta(q,r)=\frac{\beta(q,q_0)}{\beta(r,q_0)}=\frac{a_q}{a_r}.
\]

\emph{Boundary coefficients and face scales.}  If $r$ is nondegenerate with support $B\subsetneq A$, Lemma~\ref{pt:lem:branchagg}
constructs the ambient full-to-face coefficient $\beta_A(q,r)>0$ by
\[
L_q^A i_B=\beta_A(q,r)L_{r|_B}^B.
\]
When the ambient action set is clear, write this coefficient as $\beta(q,r)$.

When nondegenerate boundary priors $r,s$ have common support $B$, $\beta_A(r,s)$ means the unique $c>0$ satisfying
\[
L_r^A(i_B\sigma)=c\,L_s^A(i_B\sigma)\qquad(\sigma\in W_B).
\]
This is a scalar on the common face, not on all of $W_A$.  Since
\[
L_r^A\circ i_B=L_{r|_B}^B,\qquad L_s^A\circ i_B=L_{s|_B}^B,
\]
and the intrinsic face scalar satisfies
$L_{r|_B}^B(\sigma)=\beta_B(r|_B,s|_B)L_{s|_B}^B(\sigma)$, we have
\[
\beta_A(r,s)=\beta_B(r|_B,s|_B).
\]

Fix the same full-support basepoint $q_0$ on $\Delta(A)$.  If $r_0$ is any nondegenerate face basepoint with support $B$,
then for every nondegenerate $r$ with support $B$,
\[
\beta_A(q_0,r)=\beta_A(q_0,r_0)\,\beta_B(r_0|_B,r|_B).
\]
Indeed, both sides are the coefficient multiplying $L_{r|_B}^B$ to obtain $L_{q_0}^A\circ i_B$, and
$L_{r|_B}^B\ne0$.  Define the ambient boundary scale by
\[
a_r:=\frac{1}{\beta_A(q_0,r)}.
\]
If the face scale is based at $r_0$ with $a_{r_0}:=1/\beta_A(q_0,r_0)$, then
\[
a_{r|_B}^B:=\frac{a_{r_0}}{\beta_B(r_0|_B,r|_B)}
=\frac{1}{\beta_A(q_0,r_0)\beta_B(r_0|_B,r|_B)}
=\frac{1}{\beta_A(q_0,r)}
=a_r.
\]
Thus the ambient and intrinsic scales agree: treating $r$ as a boundary posterior in $\Delta(A)$ or as a full-support prior
on $\Delta(B)$ yields the same normalised representative $\widehat F_r$.

\emph{Extension to boundary posteriors.}  If $r$ is nondegenerate with support $B\subsetneq A$, Lemma~\ref{pt:lem:branchagg}
defines
$\beta(q,r)=\beta_A(q,r)$ by
\[
L_q^A i_B=\beta_A(q,r)L_{r|_B}^B.
\]
Using the full-support identity $L_q^A=\beta(q,q_0)L_{q_0}^A$ and restricting to $i_B(W_B)$,
\[
L_q^A i_B=\beta(q,q_0)L_{q_0}^A i_B
=\beta(q,q_0)\beta_A(q_0,r)L_{r|_B}^B.
\]
Comparing with the definition of $\beta_A(q,r)$ and using $L_{r|_B}^B\ne0$ gives
\[
\beta(q,r)=\beta(q,q_0)\beta(q_0,r)=\frac{a_q}{a_r},
\]
where $a_r=1/\beta(q_0,r)$ by the boundary-scale definition above.

\emph{Singleton scale conventions.}  If $r=\delta_a$ is a singleton prior, set $a_{\delta_a}:=a_{q_0}=1$ and
$\widehat F_{\delta_a}\equiv0$.  Define $\beta(q,\delta_a):=a_q/a_{\delta_a}$ as a convention.  No behavioural branch
coefficient is identified at a singleton target, but singleton terms in the aggregation and chain formulas are always zero.

\emph{Mixed-support cocycle.}  Let $r,s$ be nondegenerate with nested supports
$C:=\operatorname{supp}(s)\subseteq B:=\operatorname{supp}(r)$, and write $s|_B$ for $s$ viewed as a boundary prior on
the action set $B$.  Interpret
\[
\beta(r,s):=\beta_B(r|_B,s|_B),
\]
the intrinsic full-to-face coefficient on action set $B$.  Let $i_C^A:W_C\to W_A$ and $i_C^B:W_C\to W_B$ be the face
inclusions.  For every $\sigma\in W_C$,
\[
L_q^A i_C^A\sigma
=\beta(q,r)L_{r|_B}^B i_C^B\sigma
=\beta(q,r)\beta(r,s)L_{s|_C}^C\sigma.
\]
The direct full-to-face definition from Lemma~\ref{pt:lem:branchagg} for the pair $(q,s)$ gives
\[
L_q^A i_C^A\sigma=\beta(q,s)L_{s|_C}^C\sigma.
\]
Since $L_{s|_C}^C\ne0$, $\beta(q,s)=\beta(q,r)\beta(r,s)$.  If
$\operatorname{supp}(s)\not\subseteq\operatorname{supp}(r)$, then $\beta(r,s)$ is not behaviourally defined; singleton
targets are also excluded from the behavioural cocycle because their scales are conventions only.

\emph{Chain rule.}  Applying Lemma~\ref{pt:lem:branchagg} and substituting $\beta(q,r)=a_q/a_r$ gives
\[
F_q(\mu_{q,P_1\triangleright\{Q^x\}})
=
F_q(\mu_{q,P_1})+\sum_{x:m(x)>0}m(x)\,\frac{a_q}{a_{r_x}}\,F_{r_x}(\mu_{r_x,Q^x}).
\]
Dividing by $a_q$ and writing $\widehat F_q(\mu):=F_q(\mu)/a_q$ yields
\[
\widehat F_q(\mu_{q,P_1\triangleright\{Q^x\}})
=
\widehat F_q(\mu_{q,P_1})+\sum_{x:m(x)>0}m(x)\,\widehat F_{r_x}(\mu_{r_x,Q^x}).
\qedhere
\]
\end{proof}

\begin{lemma}[Coherent relabelling and face scales]\label{pt:lem:facescales}
The scales produced by Lemma~\ref{pt:lem:chain} can be chosen simultaneously across all finite action sets so that:
\begin{enumerate}
\item for every bijection $\pi:A\to A'$,
\[
a^{A'}_{q\pi}=a^A_q;
\]
\item if $\iota:B\hookrightarrow A$ is an injection with $|B|\ge2$, $r$ has full support on $B$, and $q$ has full support
on $A$, then the full-to-face branch coefficient for reaching the face $\iota(B)$ is
\begin{equation}
\beta_\iota(q,r)=\frac{a^A_q}{a^B_r}.
\tag{FS}\label{pt:eq:facescale}
\end{equation}
\end{enumerate}
Singleton target scales remain a convention because every continuation value on a singleton fibre is zero.
\end{lemma}

\begin{proof}
For every nondegenerate finite action set $A$, first choose the full-support scales supplied by the cocycle part of
Lemma~\ref{pt:lem:chain}, so that
\[
\beta_A(q,q')=\frac{a^A_q}{a^A_{q'}}
\qquad(q,q'\in\operatorname{int}\Delta(A)).
\]
They are unique up to multiplication by one positive constant depending on $A$.  Normalise initially by
$a^A_{u_A}=1$, where $u_A$ is uniform.  Exact relabelling invariance of the representatives implies relabelling
invariance of the uniquely determined branch coefficients.  Since bijections carry uniform priors to uniform priors,
the initial scales are exactly invariant under bijections.

Let $\iota:B\hookrightarrow A$ be an injection between nondegenerate finite sets.  By relabelling $B$ onto its image
$\iota(B)$ and using Corollary~\ref{pt:cor:permutationinvariance}, it is enough to define the coefficient for the inclusion
of the face $\iota(B)$; the notation below keeps the original symbol $\iota$.  Let $i_\iota$ denote the induced map on
signed posterior-law tangent spaces.  For full-support $q\in\Delta(A)$ and $r\in\Delta(B)$, let
$\beta_\iota(q,r)>0$ be the full-to-face scalar characterised by
\[
L_q^A\circ i_\iota=\beta_\iota(q,r)L_r^B
\quad\hbox{on }W_B.
\]
The number
\begin{equation}
K(\iota):=\beta_\iota(q,r)\frac{a^B_r}{a^A_q}
\tag{K}\label{pt:eq:embeddingdefect}
\end{equation}
is independent of both $q$ and $r$.  Independence of $q$ follows by comparing two full-support ambient priors and using
$L_q^A=(a_q^A/a_{q'}^A)L_{q'}^A$ on $W_A$.  Independence of $r$ follows from same-face compatibility: if $r,s$ have
full support on $B$, then $L_r^B=(a_r^B/a_s^B)L_s^B$, and hence
$\beta_\iota(q,s)=\beta_\iota(q,r)a_r^B/a_s^B$.

The defects multiply under composition.  If
$C\xhookrightarrow{\jmath}B\xhookrightarrow{\iota}A$ and all three sets in the comparison are nondegenerate, then
restricting first to $B$ and then to $C$ gives, for full-support $q,r,s$ on $A,B,C$ respectively,
\[
\beta_{\iota\circ\jmath}(q,s)=\beta_\iota(q,r)\beta_\jmath(r,s).
\]
Substituting the definition of $K$ and cancelling the intermediate scale $a_r^B$ gives
\begin{equation}
K(\iota\circ\jmath)=K(\iota)K(\jmath).
\tag{KC}\label{pt:eq:embeddingcocycle}
\end{equation}
Exact relabelling invariance shows that $K(\iota)$ depends only on $n:=|A|$ and $m:=|B|$; write it as $K_{n,m}$.
Bijections give $K_{n,n}=1$, and
\[
K_{n,\ell}=K_{n,m}K_{m,\ell}
\qquad(2\le\ell\le m\le n).
\]
Put $t_n:=K_{n,2}$, with $t_2=1$.  Then $K_{n,m}=t_n/t_m$.  Finally replace every scale on an $n$-point action set by
\[
\widetilde a_q^A:=t_n a_q^A.
\]
Under this change, the defect in~\eqref{pt:eq:embeddingdefect} becomes
\[
\widetilde K_{n,m}=K_{n,m}\frac{t_m}{t_n}=1.
\]
The rescaling is common to all priors on a fixed action set, so it leaves all within-set cocycle ratios unchanged; because
$t_n$ depends only on cardinality, it also preserves exact relabelling invariance.  Dropping tildes gives
\eqref{pt:eq:facescale}.  Singleton target scales may be set arbitrarily because all singleton continuation terms are zero.
\end{proof}

\paragraph{Stage 5: compatibility collapse and entropy identification.}
\addcontentsline{toc}{subsubsection}{Stage 5: compatibility collapse and entropy identification}
The product and sequential evaluations of full revelation must agree.  Their
compatibility produces a positive grouping-weight equation; evaluating one
distribution by two groupings forces the product interaction to vanish and
the chain scale to be universal.  Entropy reduction and Faddeev's recursion
then identify the remaining functional.

\begin{lemma}[Interaction collapse and universal chain scale]\label{pt:lem:scalecoherence}
The universal interaction coefficient in~\eqref{pt:eq:QA} is zero.  Moreover, there exists $a>0$ such that the coherently
chosen chain scales satisfy $a_q=a$ for every nondegenerate full-support prior on every finite action set.  The same
value may be assigned to singleton priors.  Consequently,
\begin{equation}
F_{p\otimes r}(\mu\otimes\nu)=F_p(\mu)+F_r(\nu)
\tag{PA}\label{pt:eq:PArecovered}
\end{equation}
and the branch formula has a universal chain scale.
\end{lemma}

\begin{proof}
Define, for every finite prior $q$ read on its support,
\[
H(q):=F_q(\chi_q),
\qquad
Z(q):=1+\kappa H(q).
\]
For a nondegenerate prior, local non-triviality gives $H(q)>0$, and~\eqref{pt:eq:QAslope} gives
\begin{equation}
Z(q)>0.
\tag{Z+}\label{pt:eq:Zpositive}
\end{equation}
For a singleton, $H(q)=0$ and $Z(q)=1$.

\emph{Step 1: product revelation links the chain scale to $Z$.}
Let $q\in\Delta(A)$ and $r\in\Delta(B)$ be nondegenerate and full support.  Quasi-additivity applied to full revelation
gives
\begin{equation}
H(q\otimes r)=H(q)+H(r)+\kappa H(q)H(r).
\tag{QH}\label{pt:eq:QH}
\end{equation}
Compute the same value sequentially.  First reveal the $A$-coordinate and then reveal the $B$-coordinate in every branch.
The first-stage value is $F_{q\otimes r}(\chi_q\otimes\delta_r)=H(q)$ by~\eqref{pt:eq:QA}.  After observing $a$, the
posterior is $e_a\otimes r$ on the face $\{a\}\times B$; support restriction and relabelling identify its continuation
full-revelation value with $H(r)$.  For the injection
\[
\iota_a:B\hookrightarrow A\times B,\qquad b\mapsto(a,b),
\]
Lemma~\ref{pt:lem:facescales} gives the branch coefficient
\[
\beta_{\iota_a}(q\otimes r,r)=\frac{a^{A\times B}_{q\otimes r}}{a^B_r}.
\]
The branch weights are $q(a)$ and the same coefficient and continuation value appear in every branch, so
$\sum_a q(a)=1$.  Lemma~\ref{pt:lem:branchagg} therefore yields
\begin{equation}
H(q\otimes r)=H(q)+\frac{a_{q\otimes r}}{a_r}H(r).
\tag{SA}\label{pt:eq:seqA}
\end{equation}
Comparing~\eqref{pt:eq:QH} and~\eqref{pt:eq:seqA}, and using $H(r)>0$,
\begin{equation}
\frac{a_{q\otimes r}}{a_r}=Z(q).
\tag{R1}\label{pt:eq:ratio1}
\end{equation}
Revealing $B$ first gives symmetrically
\begin{equation}
\frac{a_{q\otimes r}}{a_q}=Z(r).
\tag{R2}\label{pt:eq:ratio2}
\end{equation}
Hence $a_rZ(q)=a_qZ(r)$, so there is a universal constant $C>0$ such that
\begin{equation}
a_q=CZ(q)
\tag{AS}\label{pt:eq:ascaleZ}
\end{equation}
for every nondegenerate full-support prior on every finite action set.  Assign $a_q:=C$ at singleton priors, where
$Z(q)=1$ and all continuation values are zero.

If $\kappa=0$, then $Z\equiv1$, equation~\eqref{pt:eq:ascaleZ} already gives universal scales, and~\eqref{pt:eq:QA} is exact
product additivity.  For the remaining argument suppose $\kappa\ne0$.

\emph{Step 2: the induced weight equation.}
Let $q$ be full support on $A$, let $\mathcal P=\{A_1,\ldots,A_m\}$ be a partition into nonempty cells, and put
\[
p_i:=q(A_i),
\qquad
q^i:=q(\cdot\mid A_i),
\qquad
p:=(p_1,\ldots,p_m).
\]
Let $G_{\mathcal P}$ reveal the block.  Undetectable-distinctions neutrality gives
$(q,G_{\mathcal P})\sim(p,\Id_{\{1,\ldots,m\}})$.  The early cross-prior representation
Lemma~\ref{pt:lem:blockbridge}, importantly for the \emph{unscaled} $F$'s, therefore gives
\begin{equation}
F_q(\mu_{q,G_{\mathcal P}})=H(p).
\tag{BG}\label{pt:eq:blockvalue}
\end{equation}
Append full revelation in every branch.  The compound posterior law is $\chi_q$, so the chain formula and
\eqref{pt:eq:ascaleZ} give
\begin{equation}
H(q)=H(p)+Z(q)\sum_{i=1}^m p_i\frac{H(q^i)}{Z(q^i)}.
\tag{GR}\label{pt:eq:groupingH}
\end{equation}
Define
\[
w(q):=\frac1{Z(q)}>0.
\]
Using $\kappa H(x)=Z(x)-1$ in~\eqref{pt:eq:groupingH} gives
\[
\frac{Z(p)}{Z(q)}=\sum_{i=1}^m p_i\frac1{Z(q^i)}.
\]
Equivalently,
\begin{equation}
w(q)=w(p)\sum_{i=1}^m p_iw(q^i).
\tag{W}\label{pt:eq:weight}
\end{equation}
No continuity assertion is needed here.

\emph{Step 3: a two-grouping argument forces $w\equiv1$.}
All distributions in this paragraph are read on their positive-probability supports, so the weight equation applies after
support restriction.  First,~\eqref{pt:eq:weight} applied to a product distribution gives
\begin{equation}
w(u\otimes v)=w(u)w(v),
\tag{M}\label{pt:eq:wmult}
\end{equation}
because the first-coordinate block probabilities are $u$ and every conditional distribution is $v$.

Take arbitrary finite probability vectors $u$ and $v$, possibly on different alphabets, and write
\[
x:=w(u),\qquad y:=w(v),\qquad p:=(1/2,1/2).
\]
On the labelled disjoint union $(U\times U)^0\sqcup(V\times V)^1$, form the distribution
\[
T:=\tfrac12(u\otimes u)^0\;\sqcup\;\tfrac12(v\otimes v)^1.
\]
Grouping first by the top block and using~\eqref{pt:eq:wmult},
\begin{equation}
w(T)=w(p)\frac{x^2+y^2}{2}.
\tag{E1}\label{pt:eq:twoeval1}
\end{equation}
Now group first by the pair consisting of the top-block label and the first within-block coordinate.  Its marginal is
\[
S:=\tfrac12u^0\;\sqcup\;\tfrac12v^1,
\]
and~\eqref{pt:eq:weight} gives
\[
w(S)=w(p)\frac{x+y}{2}.
\]
Conditional on a point in the first part of $S$, the remaining coordinate has distribution $u$; conditional on a point in
the second part, it has distribution $v$.  A second application of~\eqref{pt:eq:weight} therefore gives
\begin{equation}
w(T)=w(S)\frac{x+y}{2}
=w(p)\left(\frac{x+y}{2}\right)^2.
\tag{E2}\label{pt:eq:twoeval2}
\end{equation}
Since $w(p)>0$, comparison of~\eqref{pt:eq:twoeval1} and~\eqref{pt:eq:twoeval2} yields
\[
\frac{x^2+y^2}{2}=\left(\frac{x+y}{2}\right)^2,
\]
so $(x-y)^2=0$.  As $u$ and $v$ were arbitrary, $w$ is constant on all finite probability vectors.  At a singleton prior,
$H=0$, hence $w=1$.  Therefore $w(q)\equiv1$.

Thus $Z(q)=1$ and $\kappa H(q)=0$ for every $q$.  Local non-triviality gives a nondegenerate $q$ with $H(q)>0$, hence
$\kappa=0$, contradicting the temporary alternative unless the interaction coefficient is zero.  Equation
\eqref{pt:eq:ascaleZ} now gives $a_q=C$ universally.  Finally~\eqref{pt:eq:QA} reduces to exact product additivity
\eqref{pt:eq:PArecovered}.
\end{proof}

\begin{lemma}[Entropy reduction and Faddeev identification]\label{pt:lem:faddeevsketch}
By Lemma~\ref{pt:lem:scalecoherence}, $a_q\equiv a$ is universal.  Define
\[
\widehat F_q:=F_q/a,
\qquad
H(q):=\widehat F_q(\chi_q),
\]
the rescaled value of full revelation.  Then, for every channel $P$,
\[
\widehat F_q(\mu_{q,P})
=
H(q)-\sum_{x:m_{q,P}(x)>0}m_{q,P}(x)\,H(r_x^{\,q,P}).
\tag{ER}
\]
Moreover $H$ satisfies Faddeev's recursion, so
\[
H(q)=\alpha \Sh(q)
\]
for some $\alpha>0$.  Consequently,
\[
\widehat F_q(\mu_{q,P})=\alpha I_{q,P}(A;X).
\]
\end{lemma}

\begin{proof}
\emph{Entropy-reduction formula.}  Boundary values of $H$ are read through support restriction: if $q$ has support
$B$, then $H(q):=H(q|_B)$, and $H(\delta_a)=0$.  Let $P:A\to\Delta(X)$ be any channel at full-support prior $q$.  Append
full revelation $\Id_A:A\to\Delta(A)$ after $P$: the compound $P\triangleright\{\Id_A\}_x$ (using $\Id_A$ in every branch)
has $\mu_{r_x,\Id_A}=\chi_{r_x}$ and the overall posterior law is full revelation $\chi_q$.  Applying the chain rule
(Lemma~\ref{pt:lem:chain}) gives
\[
\widehat F_q(\chi_q)
=
\widehat F_q(\mu_{q,P})+\sum_{x:m(x)>0}m(x)\,\widehat F_{r_x}(\chi_{r_x}).
\]
Hence
\[
\widehat F_q(\mu_{q,P})=H(q)-\sum_{x:m(x)>0}m(x)\,H(r_x).
\]
If $q$ is on the boundary, restrict $q$ and $P$ to $B=\operatorname{supp}(q)$, apply the full-support formula on
$\Delta(B)$, and pull the value and comparison back by Lemma~\ref{pt:lem:supprestrict}; zero-probability rows and records
do not affect either side of the displayed identity.

\emph{Scaled cross-prior bridge.}
Lemma~\ref{pt:lem:blockbridge} was proved for the unscaled values $F_q$.  Since Lemma~\ref{pt:lem:scalecoherence} has made
$a_q\equiv a$ universal, it is equivalently valid for
\[
V(q,P):=\widehat F_q(\mu_{q,P})=F_q(\mu_{q,P})/a:
\]
for all finite priors and channels,
\[
q^0\succeq_{P\sqcup Q}r^1
\quad\Longleftrightarrow\quad
V(q,P)\ge V(r,Q).
\]

\emph{Faddeev recursion.}  First let $q$ be full support and let $\mathcal P=\{A_1,\dots,A_m\}$ be a partition into
nonempty cells.  Then $p_i:=q(A_i)>0$ for every $i$.  Let $G_{\mathcal P}:A\to\Delta(\{1,\dots,m\})$ reveal only the
block index: $G_{\mathcal P}(i\mid a)=\1_{a\in A_i}$.  Write
\[
p_i:=q(A_i),\qquad q^i:=q(\cdot\mid A_i),
\]
and let $p=(p_1,\dots,p_m)\in\Delta(\{1,\dots,m\})$.  Consider the deterministic kernel
$S:A\to\Delta(\{1,\dots,m\})$ given by
$S(i\mid a)=\1_{a\in A_i}$, so $qS=p$.

By Lemma~\ref{pt:lem:neutrality}, since $G_{\mathcal P}$ is $S$-measurable,
\[
(q,G_{\mathcal P})\sim(p,\Id_{\{1,\dots,m\}})
\]
in the two-block comparison.  Applying the scaled form of Lemma~\ref{pt:lem:blockbridge} to the forward comparison gives
$V(q,G_{\mathcal P})\ge V(p,\Id_{\{1,\dots,m\}})$.  Applying the reverse-block form from Lemma~\ref{pt:lem:blockcoh} and then
the same bridge to the reverse comparison gives $V(p,\Id_{\{1,\dots,m\}})\ge V(q,G_{\mathcal P})$.  Therefore
\[
\widehat F_q(\mu_{q,G_{\mathcal P}})
=\widehat F_p(\chi_p)=H(p).
\]
The entropy-reduction formula applied to $G_{\mathcal P}$ gives
\[
H(q)=\widehat F_q(\chi_q)
=\widehat F_q(\mu_{q,G_{\mathcal P}})+\sum_{i=1}^m p_i\widehat F_{q^i}(\chi_{q^i})
=H(p)+\sum_{i=1}^m p_iH(q^i).
\]
This is Faddeev's recursion.  For boundary $q$, restrict to $\operatorname{supp}(q)$, apply the same full-support recursion
to the nonempty positive-mass intersections, and pull back by Lemma~\ref{pt:lem:supprestrict}; equivalently, the recursion is
read over positive-mass blocks only.

We now verify that $H$ is continuous on each $\Delta_n:=\{q\in\Delta(A):|A|=n\}$.

\emph{Binary continuity via erasure calibrators.}
Define the binary entropy function $h:[0,1]\to\R$ by $h(t):=H(t,1-t)$.  We show $h$ is continuous using
prior-independent calibrators.  By Lemma~\ref{pt:lem:convnorm}, support restriction, and $a>0$, $H(q)\ge0$ for every $q$,
so $h(t)\ge0$ on $[0,1]$.  If $c<0$, then
\[
\{t:h(t)\ge c\}=[0,1],
\qquad
\{t:h(t)\le c\}=\varnothing,
\]
both closed.  Now fix $c\ge0$.  By local non-triviality, there exists a binary prior $\theta=(s,1-s)$ with
$u:=H(\theta)>0$.  Choose $n\in\mathbb{N}$ and $\lambda\in[0,1]$ such that $c=\lambda nu$
(take $n$ large enough that $c\le nu$, then set $\lambda:=c/(nu)$).

Let $w_n:=\theta^{\otimes n}$ be the $n$-fold product prior on $B_n:=\{0,1\}^n$.  By product additivity,
$H(w_n)=nH(\theta)=nu$.  Define the erasure experiment $E_{n,\lambda}:B_n\to\Delta(\{*,1,\dots,2^n\})$ which
reveals the state fully with probability $\lambda$ and reveals nothing with probability $1-\lambda$.  By the
entropy-reduction formula,
\[
\widehat F_{w_n}(\mu_{w_n,E_{n,\lambda}})=\lambda H(w_n)=\lambda nu=c.
\]
Now consider the block action space $A_c:=\{0,1\}\sqcup B_n$ and the block-diagonal channel
$K_c:=\Id_{\{0,1\}}\sqcup E_{n,\lambda}$.  For $t\in[0,1]$, let $q_t^0$ be the prior $(t,1-t)$ supported on
block $\{0,1\}$, and let $w_n^1$ be $w_n$ supported on block $B_n$.  By the scaled form of
Lemma~\ref{pt:lem:blockbridge} (which extends to boundary priors via Lemma~\ref{pt:lem:supprestrict}),
\[
q_t^0\succeq_{K_c}w_n^1\quad\Longleftrightarrow\quad V(q_t,\Id_{\{0,1\}})=h(t)\ge c=V(w_n,E_{n,\lambda}).
\]
Since $K_c$ is fixed and the map $t\mapsto q_t^0$ is continuous in $\ell_1$, Condition~\textup{(P2)} implies that the
set $\{t\in[0,1]:h(t)\ge c\}$ is closed.  Similarly, $\{t\in[0,1]:h(t)\le c\}$ is closed.  Thus all real superlevel and
sublevel sets of $h$ are closed, so $h$ is both upper and lower semicontinuous on $[0,1]$, hence continuous.

\emph{Continuity on higher simplices by induction.}
For $n\ge3$, write any $q\in\Delta_n$ as $q=(q_1,(1-q_1)\bar q)$ where
$\bar q:=(q_2/(1-q_1),\dots,q_n/(1-q_1))\in\Delta_{n-1}$ for $q_1<1$.  The Faddeev recursion with the
binary partition $\{\{1\},\{2,\dots,n\}\}$ gives
\[
H(q)=h(q_1)+(1-q_1)H(\bar q).
\]
On the region $0<q_1<1$, this is continuous: $h$ is continuous by binary continuity, $H|_{\Delta_{n-1}}$ is
continuous by the induction hypothesis, and $q\mapsto(q_1,\bar q)$ is continuous.  If $q_1=0$, support restriction,
equivalently the positive-mass-block convention in the recursion, gives $H(q)=H(\bar q)$ and $h(0)=0$,
so continuity at that face follows from the induction hypothesis.  At $q_1=1$, the induction hypothesis and compactness
of $\Delta_{n-1}$ imply $H$ is bounded on $\Delta_{n-1}$, so
$(1-q_1)H(\bar q)\to0$ as $q_1\to1$.  Thus $H(q)\to h(1)=0=H(\delta_1)$.  The other boundary faces are handled
by support restriction, relabelling on the support, and pulling back.  This completes the induction.

We have verified the hypotheses of the strong-additivity form of Faddeev's finite entropy characterisation, as stated
for example by Baez--Fritz--Leinster \cite[Theorem~6]{BaezFritzLeinster2011}, following Faddeev's original theorem
\cite{Faddeev1956}: a function on finite probability vectors that is continuous on each closed simplex, invariant under
bijections, expansive under adjoining zero-probability states, zero on point masses, and strongly additive is a
nonnegative multiple of Shannon entropy.  Here invariance under bijections follows from
Corollary~\ref{pt:cor:permutationinvariance}, with boundary cases obtained by restricting to the support and pulling back;
expansibility follows from support restriction; and the grouping recursion above is precisely strong additivity, read
over positive-mass blocks.  Faddeev's theorem therefore gives $H(q)=\alpha\Sh(q)$ for some $\alpha\ge0$.
Local non-triviality (Condition~\textup{(P1)}) forces $\alpha>0$.

\emph{Mutual information.}  Substituting $H=\alpha\Sh$ into the entropy-reduction formula yields
\[
\widehat F_q(\mu_{q,P})
=\alpha\Sh(q)-\sum_{x:m_{q,P}(x)>0}m_{q,P}(x)\alpha\Sh(r_x^{\,q,P})
=\alpha\bigl(\Sh(q)-\E[\Sh(r_X)]\bigr)
=\alpha I_{q,P}(A;X).
\qedhere
\]
\end{proof}

\begin{lemma}[Fixed-environment mutual-information representation]\label{pt:lem:globalsketch}
For every finite channel $P:A\to\Delta(X)$ and all $q,q'\in\Delta(A)$,
\[
q\succeq_P q'
\quad\Longleftrightarrow\quad
I_{q,P}(A;X)\ge I_{q',P}(A;X).
\]
\end{lemma}

\begin{proof}
Let
\[
V(s,R):=\widehat F_s(\mu_{s,R}),
\]
with boundary priors interpreted by Lemma~\ref{pt:lem:supprestrict}.  Lemma~\ref{pt:lem:blockbridge}, divided by the universal
scale from Lemma~\ref{pt:lem:scalecoherence}, gives the cross-prior block-comparison bridge: for finite action sets,
priors $s,t$, and channels $R,S$,
\[
s^0\succeq_{R\sqcup S}t^1
\quad\Longleftrightarrow\quad
V(s,R)\ge V(t,S),
\]
including boundary priors by support restriction.  Lemma~\ref{pt:lem:faddeevsketch} gives, for every finite
prior/channel pair,
\[
V(s,R)=\alpha I_{s,R}
\]
with $\alpha>0$.

Fix $P:A\to\Delta(X)$ and $q,q'\in\Delta(A)$.  By the duplication clause of Condition~\textup{(P3)},
\[
q\succeq_P q'
\quad\Longleftrightarrow\quad
q^0\succeq_{P\sqcup P}(q')^1.
\]
Applying this bridge with $(s,R)=(q,P)$ and $(t,S)=(q',P)$ gives
\[
q^0\succeq_{P\sqcup P}(q')^1
\quad\Longleftrightarrow\quad
V(q,P)\ge V(q',P).
\]
Substituting $V=\alpha I$ and using $\alpha>0$ yields
\[
q\succeq_P q'
\quad\Longleftrightarrow\quad
I_{q,P}(A;X)\ge I_{q',P}(A;X).
\]

If $q$ or $q'$ is on the boundary, the bridge is read after restricting each pair to its positive-probability support
and then pulling the comparison back by Lemma~\ref{pt:lem:supprestrict}.  Deleting
zero-probability actions does not change mutual information, and the Condition~\textup{(P3)} equivalence remains the ambient
fixed-channel comparison for $P$.
\end{proof}

\begin{proof}[Proof of Proposition~\ref{pt:prop:main}]
That \textup{(ii)} implies \textup{(i)} is the necessity verification at the start of this appendix.  That
\textup{(i)} implies \textup{(ii)} is Lemma~\ref{pt:lem:globalsketch}.  For the moreover clause, apply \textup{(ii)} to
the fixed block environment $\bigsqcup_{k\in K}P_k$: a lottery supported on block $i$ has constant block label, so
\[
I_{q_i^i,\bigsqcup_{k\in K}P_k}=I_{q_i,P_i}(A_i;X_i),
\]
and similarly for block $j$.
\end{proof}
%</leanpureproof>

\subsection{Uniqueness}

\begin{proof}[Proof of Proposition~\ref{pt:prop:uniqueness}]
First note that the range of $(q,P)\mapsto I_{q,P}$ is $[0,\infty)$.  For $q\in\Delta(A)$ and $t\in[0,1]$, let
$E_t:A\to\Delta(A\sqcup\{*\})$ reveal the realised action with probability $t$ and otherwise report $*$:
\[
E_t(a'\mid a)=t\1_{a'=a},
\qquad
E_t(*\mid a)=1-t.
\]
Then $I_{q,E_t}=t\Sh(q)$.  In particular, if $u_n$ is uniform on $n\ge2$ actions, the values at $u_n$ fill
$[0,\log n]$.  These intervals are unbounded, so every non-negative value is attained; mutual information is
non-negative, so there are no other values.

For part \textup{(i)}, let $V$ be a common-scale representation and define
$\varphi:[0,\infty)\to\R$ by $\varphi(v):=V(q,P)$ for any pair satisfying $I_{q,P}=v$.  The definition is independent
of the chosen pair.  Indeed, if $I_{q,P}=I_{r,Q}$, the moreover clause of Proposition~\ref{pt:prop:main} gives
$q^0\sim_{P\sqcup Q}r^1$, and the common-scale property gives $V(q,P)=V(r,Q)$.  If
$I_{q,P}>I_{r,Q}$, the block comparison is strict---the reverse weak comparison would require the opposite
information inequality---so $V(q,P)>V(r,Q)$.  Thus $\varphi$ is strictly increasing and $V=\varphi\circ I$.
Conversely, Proposition~\ref{pt:prop:main} shows that every $\varphi\circ I$ with $\varphi$ strictly increasing is a
common-scale representation.

For part \textup{(ii)}, every $aI$ with $a>0$ is a common-scale representation by part \textup{(i)} and is
branch-additive by the mutual-information chain rule (N3).  Conversely, let the common-scale representation
$V=\varphi\circ I$ be branch-additive.  Composing the one-record uninformative channel with itself gives
$\varphi(0)=2\varphi(0)$, hence $\varphi(0)=0$.

Fix $n\ge2$, put $q=u_n$, and write $h_n:=\Sh(u_n)=\log n$.  For $0<\lambda<1$, let the first stage be a public coin
with two records and identical rows $(\lambda,1-\lambda)$.  Its information is zero, each posterior is $q$, and its
branch weights are $\lambda$ and $1-\lambda$.  For any $x\in[0,h_n]$, use the erasure channel $E_{x/h_n}$ in the first
branch and the uninformative channel in the second.  The chain rule (N3) makes the compound information $\lambda x$,
while branch additivity gives
\[
\varphi(\lambda x)
=\varphi(0)+\lambda\varphi(x)+(1-\lambda)\varphi(0)
=\lambda\varphi(x).
\]
The same identity holds at $\lambda=0,1$ because $\varphi(0)=0$.  Taking $x=h_n$ and $\lambda=z/h_n$ therefore shows
that
\[
\varphi(z)=\frac{\varphi(h_n)}{h_n}z
\qquad(0\le z\le h_n).
\]
The intervals $[0,h_n]$ are nested and unbounded; since they carry the same function $\varphi$, their slopes agree.
Hence $\varphi(z)=az$ on $[0,\infty)$ for a single constant $a$, and $a>0$ because $\varphi$ is strictly increasing.
\end{proof}

\subsection{Independence of the pure-record conditions}

\begin{proof}[Proof of Proposition~\ref{pt:prop:independence}]
For each condition, we construct a family of weak orders that satisfies the other four and violates it.

\emph{Condition \textup{(P1)}.}  Let every pair of lotteries be indifferent in every environment.  The graphs in
\textup{(P2)} are then the whole parameter space, and \textup{(P3)}--\textup{(P5)} hold trivially; the strict clause of
\textup{(P5)} is vacuous.  Local non-triviality fails.

\emph{Condition \textup{(P2)}.}  A discontinuous refinement of mutual information gives the required family.  Choose an
additive map $g:\R\to\R$ that is not of the form $g(x)=cx$, and, for a finite distribution $p$, set
\[
H_g(p):=-\sum_{a:p(a)>0}p(a)g(\log p(a)).
\]
Because $g(x+y)=g(x)+g(y)$, this entropy obeys the grouping identity and is additive on products, although it need not
be continuous, concave, or non-negative.  Define the associated information reduction by
\[
I^g_{q,P}:=H_g(q)-\sum_{x:m_{q,P}(x)>0}m_{q,P}(x)H_g(r_x^{q,P}),
\]
and compare procedures lexicographically:
\[
q\succeq_Pq'
\quad\Longleftrightarrow\quad
\bigl(I_{q,P},I^g_{q,P}\bigr)
\ge_{\rm lex}
\bigl(I_{q',P},I^g_{q',P}\bigr).
\]
This is a weak order, and the first coordinate gives local non-triviality.  A block-supported lottery has the same two
coordinates on its own block, so \textup{(P3)} holds.  Both coordinates obey the chain rule; hence branchwise
lexicographic improvements aggregate, with strictness, and \textup{(P5)} holds.

For \textup{(P4)}, let $(A,X,A',X')$ carry the joint law in \eqref{pt:eq:dpjoint}.  Shannon data processing settles the
comparison whenever $I(A;X)>I(A';X')$.  In the equality case,
\[
I(A;X)-I(A';X')=I(A;X\mid A')+I(A';X\mid X')=0.
\]
Thus $A\perp X\mid A'$ and $A'\perp X\mid X'$.  Product additivity makes the corresponding two conditional
$I^g$ terms zero as well, so $I^g(A;X)=I^g(A';X')$.  The second coordinate therefore ties whenever the first one ties,
which proves \textup{(P4)}.

It remains to choose $g$ so that continuity fails.  Let
$p=(1/2,1/4,1/8,1/8)$, and choose an interior $r\in\Delta(\{1,2,3,4\})$ on the same Shannon-entropy level set such that
$\log2,\log r_1,\ldots,\log r_4$ are linearly independent over $\mathbb Q$; generic $r$ on that level set has this property.
Extend these numbers to a Hamel basis and prescribe
\[
g(\log2)=0,
\qquad
g(\log r_1)=-1/r_1,
\qquad
g(\log r_i)=0\quad(i=2,3,4).
\]
Then $\Sh(p)=\Sh(r)$ but $H_g(p)=0<1=H_g(r)$.  Under full revelation, therefore, $p\not\succeq_{\Id_4}r$.
If $p_n\to p$ moves from $p$ toward the uniform distribution, strict concavity gives
$\Sh(p_n)>\Sh(p)=\Sh(r)$, so $p_n\succ_{\Id_4}r$ for every $n$.  The graph in \textup{(P2)} is not closed.

\emph{Condition \textup{(P3)}.}  The labelled copies used in block environments make a simple label-dependent family
possible.  Treat labels as syntactic tags, with the outermost tag taking precedence, and define
\[
w(a^0)=1,
\qquad
w(a^i)=0\quad(i\ne0),
\]
with $w(a)=0$ for untagged actions.  For a fixed $\varepsilon>0$, represent each environment by
\[
V_\varepsilon(q,P):=I_{q,P}+\varepsilon\sum_aq(a)w(a).
\]
This continuous criterion defines a weak order and satisfies local non-triviality.  Every comparison in \textup{(P4)}
places the original procedure in block $0$ and the processed procedure in block $1$, so its value difference is
\[
I(A;X)-I(A';X')+\varepsilon>0.
\]
For \textup{(P5)}, write $\Delta_x:=I_{r_x,Q^x}-I_{r_x,R^x}$.  A branch comparison is weak precisely when
$\Delta_x+\varepsilon\ge0$, while the aggregate value difference is
\[
\sum_x m(x)\Delta_x+\varepsilon
=\sum_x m(x)(\Delta_x+\varepsilon).
\]
The weak and strict conclusions of \textup{(P5)} follow immediately.  Duplication coherence fails: take an untagged
binary action set, $P=\Id_2$, a degenerate $q$, and a non-degenerate $q'$ close enough to $q$ that
$0<\Sh(q')<\varepsilon$.  Then $q\not\succeq_Pq'$, whereas
$q^0\succ_{P\sqcup P}(q')^1$ because the former receives the block premium.

\emph{Condition \textup{(P4)}.}  Represent every environment by
$V(q,P):=\bigl(1+\Sh(q)\bigr)I_{q,P}$.  Continuity and block coherence are immediate.  In each branch comparison the
entropy factor is common to the two continuations, so the mutual-information chain rule gives \textup{(P5)}, including
strictness.  To violate \textup{(P4)}, take $A=\{1,2\}$, $q=(0.99,0.01)$, $P=\Id_A$, and split action $1$ uniformly
between two synonymous reported labels while sending action $2$ to a third.  The reported act still determines the
original act, so mutual information is unchanged, but $\Sh(qS)>\Sh(q)$.  The processed procedure is strictly preferred,
contrary to \textup{(DP-act)}.

\emph{Condition \textup{(P5)}.}  Let $V(q,P):=\min\{I_{q,P},c\}$ for a fixed $c>0$.  This continuous, weakly increasing
transform of mutual information satisfies \textup{(P1)}--\textup{(P4)}.  For $c=1$ bit, take $q$ uniform on four
actions and let the first stage reveal which of two pairs contains the realised act.  Full revelation within each pair
strictly improves every branch from zero to one bit, but the two compound procedures carry respectively one and two bits
before capping and are therefore indifferent.  The strict conclusion of \textup{(P5)} fails.
\end{proof}



\end{document}
